%%%%%%%%%%%%%%%%%%%%%%%%%%%%%%%%%%%%%%%%%%%%%%%%%%%%%%%%%%%%%%%%%
\chapter{Unit 1: Exponential Properties}
\label{chap:unit1}
%%%%%%%%%%%%%%%%%%%%%%%%%%%%%%%%%%%%%%%%%%%%%%%%%%%%%%%%%%%%%%%%%

%%===============================================================
\section{Exponential Properties}
\label{sec:exponential-properties}
%%===============================================================

\paragraph{Exponential Properties}

Exponential properties can be used to manipulate equations involving exponential expressions and/or functions. Here are some important exponential properties:

\begin{itemize}
    \item \textbf{Negative exponent property}:
\end{itemize}
  $$
  a^{-n} = \frac{1}{a^n} \quad \text{and} \quad \frac{1}{a^{-n}} = a^n
  $$

\begin{itemize}
    \item \textbf{Product of like bases}:
\end{itemize}
  $$
  a^m a^n = a^{m+n}
  $$

\begin{itemize}
    \item \textbf{Quotient of like bases}:
\end{itemize}
  $$
  \frac{a^m}{a^n} = a^m a^{-n} = a^{m-n}
  $$

\begin{itemize}
    \item \textbf{Multiple powers}:
\end{itemize}
  $$
  (a^m)^n = a^{mn} = (a^n)^m
  $$

\begin{itemize}
    \item \textbf{Product to a power}:
\end{itemize}
  $$
  (ab)^n = a^n b^n
  $$

\begin{itemize}
    \item \textbf{Quotient to a power}:
\end{itemize}
  $$
  \left(\frac{a}{b}\right)^n = \frac{a^n}{b^n}
  $$

\paragraph{Special cases involving 0:}

\begin{itemize}
    \item For any nonzero number $a$, $a^0 = 1$.
    \item For any positive number $m$, $0^m = 0$.
    \item $0^m$ does not exist if $m$ is negative (since $0^{-n} = \frac{1}{0^n} = \frac{1}{0}$).
    \item $0^0$ is indeterminate or undefined depending on the context.
\end{itemize}

\paragraph{Resources}

\begin{itemize}
    \item \href{https://tutoring.asu.edu/sites/default/files/exponentialandlogrithmicproperties.pdf}{Exponential and Logarithmic Properties, Arizona State University}
    \item \href{https://www.khanacademy.org/math/cc-eighth-grade-math/cc-8th-numbers-operations/cc-8th-exponent-properties/a/exponent-properties-review}{Exponent Properties Review, Khan Academy}
\end{itemize}

\subsection{Learning Objectives}

\begin{enumerate}
    \item \textbf{Understanding Exponential Properties}: Students will identify and apply the properties of exponents in various mathematical contexts.
    \item \textbf{Manipulating Exponential Expressions}: Students will simplify and manipulate exponential expressions using the correct properties.
    \item \textbf{Special Cases with Exponents}: Students will recognize and correctly handle special cases involving exponents, including zero and negative exponents.
\end{enumerate}

\subsection{Assessment Instrument}

\begin{enumerate}
    \item Part 1: Identifying Exponential Properties
    \item State the negative exponent property and provide an example.
    \item Explain the product of like bases property and simplify $2^3 \cdot 2^4$.
    \item Describe the quotient of like bases property and simplify $\frac{5^7}{5^2}$. Part 2: Simplifying Exponential Expressions
    \item Simplify the following expressions:
    \item $(3^2)^4$
    \item $(ab)^3$ where $a = 2$ and $b = 5$
    \item $\left(\frac{4}{2}\right)^3$
    \item Simplify and combine the following expressions using the properties of exponents:
    \item $2^5 \cdot 2^{-3}$
    \item $\frac{7^3}{7^1}$
    \item $(x^2)^3 \cdot x^{-4}$ Part 3: Applying Special Cases
    \item Evaluate the following expressions:
    \item $10^0$
    \item $0^4$
    \item Explain why $0^{-1}$ does not exist.
    \item Determine if $0^0$ is defined and explain your reasoning.
    \item Solve for $x$ in the equation $3^x = 27$.
    \item Determine the value of $y$ in the expression $y = 2^{-3} \cdot 2^5$. Part 4: Word Problems Involving Exponents
    \item If a bacteria population doubles every hour, write an exponential expression to represent the population after $t$ hours if the initial population is 100. Calculate the population after 5 hours.
    \item A radioactive substance decays at a rate of $\frac{1}{2}$ per year. Write an exponential expression to represent the remaining quantity after $t$ years if the initial quantity is 200 grams. Calculate the remaining quantity after 3 years.
\end{enumerate}

%%===============================================================
\section{Algebraic Expressions}
\label{sec:algebraic-expressions}
%%===============================================================

\paragraph{Algebraic Expressions}



\paragraph{Introduction}



Algebraic expressions are made up of terms. A term is a constant or the product of a constant and one or more variables. Some examples of terms are $7$, $y$, $5x^2$, $9a$, and $13xy$. The constant that multiplies the variable(s) in a term is called the coefficient. We can think of the coefficient as the number in front of the variable. For example, the coefficient of the term $3x$ is 3. When we write $x$, the coefficient is 1, since $x = 1(x)$. To evaluate an algebraic expression means to find the value of the expression when the variable is replaced by a given number. To evaluate an expression, we substitute the given number for the variable in the expression, and then simplify the expression using the order of operations.



Some terms share common traits. Terms are called like terms if they have the same variables and exponents. All constant terms are also like terms. Consider the terms $5x$, $7$, $n^2$, $4$, $3x$, $9n^2$, $-2xy$, and $8xy$. From this list, 7 and 4 are like terms, $5x$ and $3x$ are like terms, $-2xy$ and $8xy$ are like terms, and $n^2$ and $9n^2$ are like terms.



We can simplify an expression by combining the like terms. To do this, we add the coefficients and keep the same variable. For example, say we want to simplify the expression $x^2 + 4x^2 - 3x + 6 + x$. The coefficients of the like terms $x^2$ and $4x^2$ are 1 and 4, and add up to 5. The coefficients of the like terms $-3x$ and $x$ are 1 and -3, and add to -2. 6 is not like to any of the other terms, so the constant term remains as 6. So, we can combine like terms and get that $x^2 + 4x^2 - 3x + 6 + x = (x^2 + 4x^2) + (-3x + x) + 6 = 5x^2 - 2x + 6$ by adding up the coefficients of like terms.



\paragraph{Expressions and Terms}



All algebraic equations can be considered as a series of expressions. On each side of the equals ''='' sign there will be an expression. For example:



$$

5x + 2x = 7

$$



contains the expressions $(5x + 2x)$ and $(7)$.



Within each expression, there will be a number of terms. A term is any variable, number, or algebraic letter. In the equation above we can say that there are three terms. These three terms are: $5x$, $2x$, and $7$.



Now what is key to all equations is simplifying what they say, that is reducing the length of the expression by collecting like terms together. We have seen that the equation above technically has three terms. Let us break down a term further.



The term $5x$ has a coefficient of $5$ and a variable of $x$.



When you simplify an expression, you simply collect the like terms together. In other words, you collect the terms with the same variables together.



So taking our very first equation: $5x + 2x = 7$ we can see that the terms $5x$ and $+2x$ contain the same variable $x$. They can therefore be added together to obtain $7x$.



So, from above we can see that like terms can only be collected together in order to simplify an expression and hence an equation. Terms which do not contain identical variables cannot be simplified. The examples below show collecting like and unlike terms together.



\paragraph{Example 1}



Simplify the expression $5x^2 + 2x + 3x + 5 + 7$.



We can see that:



\begin{itemize}
    \item There is only 1 term in $x^2$ ($5x^2$) and so it cannot be simplified.
    \item There are two terms in $x$ ($+2x$ and $+3x$) so these can be added.
    \item There are two terms with no variable ($+5$ and $+7$) and these can be added.
\end{itemize}

Therefore, after simplifying the expression we finish with $5x^2 + 5x + 12$.



\paragraph{Example 2}



Simplify the expression $3x^2 + 2x + 2y - 4y - 3x + x^2 - xy$.



Using the rules stated above, the expression can be simplified to:



$$

4x^2 - x - 2y - xy

$$



\paragraph{Single Variables}



In example 2 above we ended up with the term $-x$. It is important to note that when there is just a variable, with no apparent coefficient, the coefficient is in fact 1.



$-x$ is the same as $-1x$.



$xy$ is the same as $1xy$.





\paragraph{Resources}



\begin{itemize}
    \item \href{https://opentextbc.ca/prealgebraopenstax/chapter/evaluate-simplify-and-translate-expressions/#:\textasciitilde{}:text=To\%20evaluate\%20an\%20algebraic\%20expression,using\%20the\%20order\%20of\%20operations.}{Evaluate, Simplify, and Translate Expressions, OpenStax}
    \item \href{https://courses.lumenlearning.com/waymakercollegealgebra/chapter/evaluate-and-simplify-algebraic-expressions/}{Evaluate and Simplify Algebraic Expressions, Lumen Learning}
    \item \href{https://www.palmbeachstate.edu/prepmathlw/Documents/algebraicexpressions.pdf}{Evaluating Algebraic Expressions, Palm Beach State College}
\end{itemize}

\paragraph{Licensing}



Content obtained and/or adapted from:



\begin{itemize}
    \item \href{https://en.wikibooks.org/wiki/A-level_Mathematics/Edexcel/Core_1/Algebra/Expressions,_Terms_and_Simplifying_Equations}{Algebraic Expressions, Terms, and Simplification, Wikibooks: A-level Mathematics} under a CC BY-SA license
    \item \href{https://en.wikibooks.org/wiki/A-level_Mathematics/Edexcel/Core_1/Algebra/Expressions,_Terms_and_Simplifying_Equations}{Simplifying Algebraic Expressions, Lumen Learning Prealgebra} under a CC BY license
\end{itemize}

\subsection{Learning Objectives}

\begin{enumerate}
    \item \textbf{Classifying Numbers}: Students will accurately classify numbers as natural, whole, integers, rational, or irrational.
    \item \textbf{Evaluating Algebraic Expressions}: Students will correctly evaluate algebraic expressions for given values of variables.
    \item \textbf{Simplifying Algebraic Expressions}: Students will proficiently simplify algebraic expressions by combining like terms.
\end{enumerate}

\subsection{Assessment Instrument}

\begin{enumerate}
    \item Part 1: Classifying Numbers
    \item Classify each of the following numbers as natural, whole, integer, rational, or irrational:
    \item $7$
    \item $0$
    \item $-5$
    \item $2.75$
    \item $\sqrt{2}$ Part 2: Evaluating Algebraic Expressions
    \item Evaluate the following algebraic expressions for the given values of variables:
    \item $x + 5$ for $x = 3$
    \item $\frac{4}{3} \pi r^3$ for $r = 2$
    \item $\sqrt{2m^3n^2}$ for $m = 2$ and $n = 3$ Part 3: Simplifying Algebraic Expressions
    \item Simplify each expression using the properties of real numbers:
    \item $-\left( \frac{23}{5} \right) \cdot 11 \cdot \left( -\frac{5}{23} \right)$
    \item $5 \cdot (6.2 + 0.4)$
    \item $18 - (7 - 15)$
    \item $\frac{17}{18} + \frac{4}{9} + \left( -\frac{17}{18} \right)$
    \item $6 \cdot (-3) + 6 \cdot 3$ Part 4: Applying Properties of Real Numbers
    \item Solve the following equations:
    \item $2x + 1 = 7$
    \item $3y - 4 = 11$
    \item $\frac{t}{2t - 1} = 1$ for $t = 10$
\end{enumerate}

%%===============================================================
\section{Scientific Notation}
\label{sec:scientific-notation}
%%===============================================================

\paragraph{Scientific Notation}

Scientific notation is a way of expressing numbers that are too large or too small (usually would result in a long string of digits) to be conveniently written in decimal form. It may be referred to as scientific form or standard index form, or standard form in the UK. This base ten notation is commonly used by scientists, mathematicians, and engineers, in part because it can simplify certain arithmetic operations. On scientific calculators, it is usually known as "SCI" display mode.

In scientific notation, nonzero numbers are written in the form

$$
m \times 10^n
$$

or $m$ times ten raised to the power of $n$, where $n$ is an integer, and the coefficient $m$ is a nonzero real number (usually between 1 and 10 in absolute value, and nearly always written as a terminating decimal). The integer $n$ is called the exponent and the real number $m$ is called the significand or mantissa. The term "mantissa" can be ambiguous where logarithms are involved, because it is also the traditional name of the fractional part of the common logarithm. If the number is negative, then a minus sign precedes $m$, as in ordinary decimal notation. In normalized notation, the exponent is chosen so that the absolute value (modulus) of the significand $m$ is at least 1 but less than 10.

Decimal floating point is a computer arithmetic system closely related to scientific notation.

\paragraph{Significant Figures}

A significant figure is a digit in a number that adds to its precision. This includes all nonzero numbers, zeroes between significant digits, and zeroes indicated to be significant. Leading and trailing zeroes are not significant digits because they exist only to show the scale of the number. Unfortunately, this leads to ambiguity. For example, the number 1230400 is usually read to have five significant figures: 1, 2, 3, 0, and 4, with the final two zeroes serving only as placeholders and adding no precision. The same number, however, would be used if the last two digits were also measured precisely and found to equal 0—resulting in seven significant figures.

When a number is converted into normalized scientific notation, it is scaled down to a number between 1 and 10. All of the significant digits remain, but the placeholding zeroes are no longer required. Thus, 1230400 would become 1.2304e6 if it had five significant digits. If the number were known to six or seven significant figures, it would be shown as 1.23040e6 or 1.230400e6. Thus, an additional advantage of scientific notation is that the number of significant figures is unambiguous.

\paragraph{Estimated Final Digits}

It is customary in scientific measurement to record all the definitely known digits from the measurement and to estimate at least one additional digit if there is any information at all available on its value. The resulting number contains more information than it would without the extra digit, which may be considered a significant digit because it conveys some information leading to greater precision in measurements and in aggregations of measurements (adding them or multiplying them together).

Additional information about precision can be conveyed through additional notation. It is often useful to know how exact the final digit is. For instance, the accepted value of the mass of the proton can properly be expressed as (51), which is shorthand for 0.00000000051.

\paragraph{E Notation}

Most calculators and many computer programs present very large and very small results in scientific notation, typically invoked by a key labelled EXP (for exponent), EEX (for enter exponent), EE, EX, E, or $\times 10^x$ depending on vendor and model. Because superscripted exponents like 107 cannot always be conveniently displayed, the letter E (or e) is often used to represent "times ten raised to the power of" (which would be written as "× 10^n") and is followed by the value of the exponent; in other words, for any two real numbers $m$ and $n$, the usage of "mEn" would indicate a value of $m \times 10^n$. In this usage, the character e is not related to the mathematical constant $e$ or the exponential function $e^x$ (a confusion that is unlikely if scientific notation is represented by a capital E). Although the E stands for exponent, the notation is usually referred to as (scientific) E notation rather than (scientific) exponential notation. The use of E notation facilitates data entry and readability in textual communication since it minimizes keystrokes, avoids reduced font sizes, and provides a simpler and more concise display, but it is not encouraged in some publications.

\paragraph{Resources}

\begin{itemize}
    \item \href{https://openstax.org/books/college-algebra/pages/1-2-exponents-and-scientific-notation}{Exponents and Scientific Notation, OpenStax}
    \item \href{https://www.chem.tamu.edu/class/fyp/mathrev/mr-scnot.html}{Scientific Notation, Texas A\&M University}
    \item \href{https://www.mathsisfun.com/numbers/scientific-notation.html}{Scientific Notation, Math Is Fun}
\end{itemize}

\paragraph{Licensing}

Content obtained and/or adapted from:

\begin{itemize}
    \item \href{https://en.wikipedia.org/wiki/Scientific_notation}{Scientific notation, Wikipedia} under a CC BY-SA license
\end{itemize}

\subsection{Learning Objectives}

\begin{enumerate}
    \item Understanding Scientific Notation: Students will be able to express large and small numbers in scientific notation, identifying the significand and exponent, and interpret values given in E notation.
    \item Applying Significant Figures: Students will be able to determine the number of significant figures in a measurement, convert numbers to scientific notation while preserving precision, and distinguish between significant and placeholder zeroes.
    \item Evaluating Precision: Students will be able to estimate and record significant digits in scientific measurements, apply the concept of precision in calculations, and understand the role of final estimated digits in enhancing measurement accuracy.
\end{enumerate}

\subsection{Assessment Instrument}

\begin{enumerate}
    \item Part 1: Scientific Notation \textbf{Convert to Scientific Notation}: Convert the following numbers to scientific notation:
    \item $0.000345$
    \item $12300000$
    \item $0.00456$ \textbf{E Notation Interpretation}: Express the following in E notation and convert to standard scientific notation:
    \item $5.67E-4$
    \item $3.45E6$ Part 2: Significant Figures \textbf{Identify Significant Figures}: Determine the number of significant figures in each of the following numbers:
    \item $450.0$
    \item $0.00789$
    \item $1.0200$ \textbf{Convert with Precision}: Convert the following numbers to scientific notation, maintaining the specified number of significant figures:
    \item $789000$ (to 3 significant figures)
    \item $0.0002345$ (to 4 significant figures) Part 3: Application \textbf{Measurement Precision}: Given the measurement of $0.0001234$ meters, record it with 3 significant figures in scientific notation. \textbf{Estimate and Record}: If a proton mass is given as $0.00000000051$ kg, express this value using E notation and explain the significance of the additional digit in scientific measurement.
\end{enumerate}

%%===============================================================
\section{Radical \& Rational Exponent}
\label{sec:radical-rational-exponent}
%%===============================================================

\paragraph{Radical \& Rational Exponent}

The principal $n$th root of $a$ is written as $\sqrt[n]{a}$, where $n$ is a positive integer greater than or equal to 2. In the radical expression, $n$ is called the index of the radical. Radical expressions can also be written without using the radical symbol, using rational (fractional) exponents. The index must be a positive integer, and if the index $n$ is even, then $a$ cannot be negative.

$$\sqrt[n]{a} = a^{\frac{1}{n}}$$

We can also have rational exponents with numerators other than 1. In these cases, the exponent must be a fraction in lowest terms. We raise the base to a power and take an $n$th root. The numerator tells us the power and the denominator tells us the root.

$$a^{\frac{m}{n}} = \left(\sqrt[n]{a}\right)^m = \sqrt[n]{a^m}$$

All of the properties of exponents that we learned for integer exponents also hold for rational exponents.

\paragraph{Resources}

\begin{itemize}
    \item \href{https://www.khanacademy.org/test-prep/sat/sat-math-practice/new-sat-passport-advanced-mathematics/v/sat-math-p3-easier}{Radicals and rational exponents - Basic example, Khan Academy}
    \item \href{https://openstax.org/books/college-algebra/pages/1-3-radicals-and-rational-exponents}{Radicals and Rational Exponents, OpenStax}
\end{itemize}

\paragraph{Licensing}

Content obtained and/or adapted from:

\begin{itemize}
    \item \href{https://openstax.org/books/college-algebra/pages/1-3-radicals-and-rational-exponents}{Radicals and Rational Exponents, OpenStax College Algebra} under a CC BY-SA license
\end{itemize}

\subsection{Learning Objectives}

\begin{enumerate}
    \item \textbf{Understand Radical Expressions:}
    \item Students will be able to express the principal $n$th root of a number $a$ using both radical notation ($\sqrt[n]{a}$) and rational exponents ($a^{\frac{1}{n}}$).
    \item \textbf{Apply Rational Exponents:}
    \item Students will simplify expressions involving rational exponents, including cases where the exponent is a fraction ($a^{\frac{m}{n}}$), and understand the connection between fractional powers and roots.
    \item \textbf{Use Exponent Properties:}
    \item Students will apply properties of exponents to rational exponents, ensuring consistent use of rules across integer and fractional exponents. ---
\end{enumerate}

\subsection{Assessment Instrument}

\begin{enumerate}
    \item \textbf{Part 1: Radical Expressions}
    \item \textbf{Convert to Exponent Form:}
    \item Express the following radical expressions as rational exponents:
    \item $\sqrt[3]{27}$
    \item $\sqrt[4]{16}$
    \item $\sqrt{25}$
    \item \textbf{Convert to Radical Form:}
    \item Express the following rational exponents as radical expressions:
    \item $27^{\frac{1}{3}}$
    \item $16^{\frac{1}{4}}$
    \item $25^{\frac{1}{2}}$ \textbf{Part 2: Simplifying Expressions}
    \item \textbf{Simplify Using Rational Exponents:}
    \item Simplify the following expressions:
    \item $\left(81^{\frac{1}{4}}\right)^2$
    \item $\left(64^{\frac{2}{3}}\right)^{\frac{3}{2}}$
    \item $\frac{27^{\frac{5}{3}}}{3^{\frac{4}{3}}}$
    \item \textbf{Simplify and Compare:}
    \item Compare and simplify the following expressions to verify equality:
    \item $\sqrt[5]{32} \text{ and } 32^{\frac{1}{5}}$
    \item $\left(\sqrt{16}\right)^{\frac{3}{2}} \text{ and } 16^{\frac{3}{4}}$ \textbf{Part 3: Applying Exponent Properties}
    \item \textbf{Apply Properties to Rational Exponents:}
    \item Solve the following equations and simplify:
    \item $\left(a^{\frac{2}{3}}\right)^3$
    \item $\frac{b^{\frac{5}{2}}}{b^{\frac{1}{2}}}$
    \item $\left(\sqrt[3]{x^6}\right)^{2}$
\end{enumerate}

%%===============================================================
\section{Polynomials}
\label{sec:polynomials}
%%===============================================================

\paragraph{Polynomials}



A polynomial of one variable, $x$, is an algebraic expression that is a sum of one or more monomials. The degree of the polynomial is the highest degree of the monomials in the sum. A polynomial $P(x)$ can generically be expressed in the form:



$$P(x) = a\_n x^n + \cdots + a\_i x^i + \cdots + a\_2 x^2 + a\_1 x + a\_0$$





or



$$

P(x) = \sum\_{i=0}^n a\_i x^i$$







The constants $a\_i$ are called the coefficients of the polynomial. Each of the individual monomials in the above sum, whose coefficient $a\_i \neq 0$, is called a term of the polynomial. When $i = 0$, $x^i = 1$ and the corresponding term simply equals the constant $a\_i$. Also, when $i = 1$, the corresponding term equals $a\_i x$.



A polynomial having two terms is called a binomial. A polynomial having three terms is called a trinomial.



\paragraph{Polynomial Equations}



We refer to all functions with one independent variable as $P(x)$. Each instance of $P(x)$ can be represented by an equation (either a monomial or a polynomial) which may have one or more places where the dependent variable is equal to zero. These places are called roots and they represent the number(s) whose value(s) for $x$ make the function $P(x) = 0$ true. These roots are called the zeroes of the polynomial (singular is zero). 



A polynomial of degree 1 will always look like a line when you graph it and always has 1 real zero. A polynomial of degree 2, a quadratic function, can have 0, 1, or 2 real zeroes. A polynomial of degree 3 (a cubic function) can have 1 or 3 real zeroes. A polynomial of degree 4 can have 0, 2, or 4 real zeroes. Complex (unreal) zeroes, when present, always come in pairs. In general, a polynomial of degree $n$, where $n$ is odd, can have from 1 to $n$ real zeroes. A polynomial of degree $n$, where $n$ is even, can have from 0 to $n$ real zeroes.



When we graph polynomials, each zero is a place where the polynomial crosses the x-axis. A polynomial of degree one can be generically written as:



$$

P(x) = Mx + C

$$



where $M$ and $C$ can be any real number. Quadratic functions are curves. The curve can bend before it ever touches the x-axis, in which case it has no zeroes; it can bend just as it touches the x-axis, in which case it can have just one zero; or it can open up above or below the x-axis, in which case it will have two zeroes. Polynomials with an odd degree (1, 3, 5, ...) must cross the x-axis at least once. Polynomials with an even degree (2, 4, 6, ...) might always be positive or negative and never have a zero.



Normally we represent a function in the form $P(x) = y$, but when we are looking for the roots of the function, we want $y$ to be equal to zero so we solve for the equation of $P(x)$ where $P(x) = 0$.







\paragraph{Solving Polynomial Equations}



Some polynomial equations can be solved by factoring, and all equations of degrees 1-4 can be solved completely by formulas. Above degree 4, there are no formulas for solving completely, and you must rely on numerical analysis or factoring. This means that for polynomials of degree greater than 4, it is often impossible to find exact solutions.



\paragraph{Rational Roots of Polynomial Equations}



Often we are interested in the rational roots of polynomials. A root is much like a factor of a number. For instance, all even numbers have a factor of two. This means you can write even numbers as two times another number. That is, the numbers 2, 4, 6, 8, ... can be written as $2 \cdot 1$, $2 \cdot 2$, $2 \cdot 3$, $2 \cdot 4$, ... . This fact is helpful when you have a fraction of two even numbers. Given a fraction of two even numbers called $N/M$, you could reduce the fraction by re-writing it as $2 \cdot n / 2 \cdot m$. By keeping fractions in lowest terms, it's easier to know when you can add or subtract them without looking for a common denominator.



\paragraph{An Example of a Use of a Polynomial Equation}



There is a story that in grade school the mathematician Gauss was asked to add the numbers 1 to 100 sequentially. He is said to have intuited the sum could be expressed with the formula $n(n+1)/2$ and quickly gave the answer 5050. The basis of this formula is that the numbers 1 through 49 added to the numbers 99 through 51 each yield 100. It is interesting to look at how this formula works for the values 9 and 10. For 10, we add the numbers $1+9$, $2+8$, $3+7$, $4+6$ to get 40 and we add the two remaining terms 5 and 10 to get 55. For 9, we add the terms $1 + 8$, $2 + 7$, $3 + 6$, $4+5$ to get $4 \cdot 9 = 36 + 9 = 45$. In the first case, $n + 1$ is the odd number and represents adding the 10 and the middle number, the 5. In the second case, $n$ is the odd number and the $n+1$ represents the sum for the preceding terms in the formula. You may or may not find stories like this intriguing based on how your personality reacts to what is known as the foundational crisis of mathematics. Learning mathematics is a lot like learning a foreign language. Some people seem more adept at learning languages than others, but with hard work, learning a new language is something we can all do.



\paragraph{Multiplying Polynomials Together}



When we multiply polynomials together, we rely heavily on the distributive property.



For instance, when we multiply 67 by 5, we can divide the equation into $(60 + 7) \cdot 5 = (300 + 35) = 335$. Additionally, we can apply the commutative property to multiply multidigit numbers. $67 \cdot 25 = (60 + 7)(20 + 5) = ((60 + 7) \cdot 20) + ((60 + 7) \cdot 5) = (60 \cdot 20) + (7 \cdot 20) + (60 \cdot 5) + (7 \cdot 5) = 1200 + 140 + 300 + 35 = 1675$. These properties are the foundation for the different forms of the mechanical calculating tool, the abacus.



When multiplying polynomials together, we do similar operations. We use the commutative property to divide the multiplier into its component parts and multiply the multiplicand by each of these parts. For instance, to multiply $x^2 + x$ by $x + 1$, we first write the multiplicand and multiplier in terms of powers of $x$. This gives us $x^2 + x + 0x^0$ and $x + 1x^0$. The terms raised to the zero power represent constant integer terms in our equations. Next, we apply the commutative property to rewrite the equations as:



$$

[(x^2 + x + 0x^0) \cdot 1x^1] + [(x^2 + x + 0x^0) \cdot 1x^0]

$$



We simplify these equations to be:



$$

[x^3 + x^2 + 0x^1] + [x^2 + x + 0x^0]

$$



(notice how our integer term drops out). Finally, we combine like terms to get the answer $x^3 + 2x^2 + x + 0x^0$. 



By breaking a polynomial into its roots, if we have a polynomial $P(x)$:



$$

P(x) = a\_n x^n + \cdots + a\_i x^i + \cdots + a\_2 x^2 + a\_1 x + a\_0

$$



The only possible rational roots (roots of the form $\frac{p}{q}$) are in the form:



$$

\frac{p}{q} = \frac{\text{factor of } a\_0}{\text{factor of } a\_n}

$$





\paragraph{Binomials}



A binomial is a sum or difference of two monomials. These can also be called polynomials, but to specify, these are binomials.



\paragraph{Examples}



\begin{itemize}
    \item $2x + 2$
    \item $2y - 7$
\end{itemize}

\paragraph{How to Factor}



To factor binomials, find the greatest common factor between the terms and factor.



\textbf{Example:}



$4x + 2$



The greatest common factor between these terms is 2 because both of the terms can be divided by it and the coefficient and constant is still an integer. The example factored would become:



$2(2x + 1)$



\paragraph{Symmetrical Polynomials}



In mathematics, a polynomial is considered to be symmetrical if you take the roots of the original polynomial and then interchange any root with another root, the polynomial will remain the same. For example, the polynomial 



$8x^3 + 16x^2 - 22x - 30$ 



is symmetrical because its factorized form is 



$(2x - 3)(2x + 2)(2x + 5)$ 



and if you interchange the roots, the resulting polynomial will be the same. However, the polynomial 



$4x^3 + 12x^2 - 7x - 30$ 



is not symmetrical because its factorized form is 



$(2x - 3)(x + 2)(2x + 5)$ 



and if you interchange the roots, the resulting polynomial will be 



$4x^3 + 18x^2 - 16x - 30$ 



if you switch the 2 and the 5 around.



\paragraph{Roots of Quadratic Polynomials}



If we need to find the roots of a given quadratic function, we have two formulae that can help us find the roots of a quadratic equation.



Let $\alpha$ and $\beta$ be the roots of $ax^2 + bx + c = 0$. Then,



$$

\alpha + \beta = -\frac{b}{a}, \quad \alpha \beta = \frac{c}{a}

$$



\textbf{Example:}



Find the values of $a$ and $b$ of the equation $ax^2 + bx - 48$ if $\alpha + \beta = 6$ and $\alpha \beta = -16$.



First, we need to find the value of $a$ and $b$. We use the relationships of the roots to find $a$ and $b$.



$$

\alpha \beta = -16 = \frac{-48}{a}

$$



From this, we can determine that $a = 3$.



$$

\alpha + \beta = 6 = -\frac{b}{a}

$$



Now that we have determined that $a = 3$, we can write the second relationship as:



$$

\alpha + \beta = 6 = -\frac{b}{3}

$$



So we can determine that $b = -18$.



Now we can write the complete equation:



$$

3x^2 - 18x - 48

$$



\paragraph{Roots of Cubic Equations}



If we need to find the roots of a given cubic function, we have three formulae that can help us find the roots of a cubic equation.



Let $\alpha$, $\beta$, and $\gamma$ be the roots of $ax^3 + bx^2 + cx + d = 0$. Then,



$$

\sum \alpha = -\frac{b}{a}, \quad \sum \alpha \beta = \frac{c}{a}, \quad \alpha \beta \gamma = -\frac{d}{a}

$$



Where:



$$

\sum \alpha = \alpha + \beta + \gamma

$$



And:



$$

\sum \alpha \beta = \alpha \beta + \alpha \gamma + \beta \gamma

$$



\textbf{Example:}



In this example, we consider the special case of the cubic 



$$

x^3 + 21x^2 + cx + 280 = 0

$$ 



where $c$ is to be determined, and we are given the additional information that its 3 roots are in arithmetic progression. Thus, we can write the roots in the form $p$, $p + q$, $p - q$. Also, factorize the equation.



First, to find $p$, we use the $\sum \alpha$:



$$

\sum \alpha = p + (p + q) + (p - q) = -21

$$



$$

\sum \alpha = 3p = -21

$$



$$

p = -7

$$



Then we need to find the value of $q$:



$$

-7(-7 + q)(-7 - q) = -280

$$



$$

7q^2 - 343 = -280

$$



$$

7q^2 = 63

$$



$$

q^2 = 9

$$



$$

q = 3

$$



Now we can write out our roots:



$$

(-7 - 3), -7, (-7 + 3)

$$



$$

-10, -7, -4

$$



We can now find $c$:



$$

\sum \alpha \beta = -7 \times -4 + -7 \times -10 + -10 \times -4 = c

$$



$$

138 = c

$$



The complete equation is:



$$

x^3 + 21x^2 + 138x + 280 = 0

$$



Finally, we write out the factorized equation:



$$

(x + 7)(x + 4)(x + 10) = 0

$$



\paragraph{Simple Substitution of Roots}



If you increase each root in a polynomial equation by the number $n$, you can calculate the resulting equation by replacing each $x$ term in the original polynomial equation with $(x - n)$. This leads to binomial expansion, so make sure that you are well-versed in it.



\textbf{Example:}



Suppose that the cubic equation 



$$

x^3 + x^2 - 22x - 40 = 0

$$ 



has roots $\alpha$, $\beta$, and $\gamma$. Find a cubic equation with the roots $\alpha - 2$, $\beta - 2$, and $\gamma - 2$.



If $x = \alpha - 2$, then $\alpha = x + 2$. Since $\alpha$ is a root of the original equation, you can replace each $x$ term with $x + 2$:



$$

(x + 2)^3 + (x + 2)^2 - 22(x + 2) - 40 = 0

$$



Using Binomial Expansion, we can easily find the terms:



$$

(x^3 + 6x^2 + 12x + 8) + (x^2 + 4x + 4) - 22(x + 2) - 40 = 0

$$



Finally, we combine all the terms and we have:



$$

x^3 + 7x^2 - 6x - 72

$$





\paragraph{Resources}



\begin{itemize}
    \item \href{https://en.wikibooks.org/wiki/Algebra/Polynomials}{Polynomials, Wikibooks: Algebra}
    \item \href{https://en.wikibooks.org/wiki/A-level_Mathematics/OCR/FP1/Roots_of_Polynomial_Equations}{Roots of Polynomial Equations, Wikibooks: A-level Mathematics}
    \item \href{https://mathresearch.utsa.edu/wikiFiles/MAT1053/Intro_to_Polynomial_and_Power_Functions/MAT1053_M2.1Intro_to_Power_Functions_and_Polynomial_Functions.pdf}{Intro to Power Functions and Polynomial Functions, Book Chapter}
    \item \href{https://mathresearch.utsa.edu/wikiFiles/MAT1053/Intro_to_Polynomial_and_Power_Functions/MAT1053_M2.1Intro_to_Power_Functions_and_Polynomial_FunctionsGN.pdf}{Guided Notes}
    \item \href{https://www.youtube.com/watch?v=nJLzFBsSX3o}{Multiplying Polynomials. Produced by TA Catherine Sporer, UTSA}
\end{itemize}

\paragraph{Licensing}



Content obtained and/or adapted from:



\begin{itemize}
    \item \href{https://en.wikibooks.org/wiki/Algebra/Polynomials}{Polynomials, Wikibooks: Algebra} under a CC BY-SA license
    \item \href{https://en.wikibooks.org/wiki/A-level_Mathematics/OCR/FP1/Roots_of_Polynomial_Equations}{Roots of Polynomial Equations, Wikibooks: A-level Mathematics} under a CC BY-SA license
\end{itemize}

\subsection{Learning Objectives}

\begin{enumerate}
    \item \textbf{Understand Polynomial Definitions:} Students will define polynomials, binomials, and trinomials, and explain their components, including coefficients and terms.
    \item \textbf{Solve Polynomial Equations:} Students will solve polynomial equations of degrees 1-4 using factoring, the quadratic formula, and other methods, and understand the conditions for finding rational roots.
    \item \textbf{Analyze Polynomial Graphs:} Students will analyze and interpret polynomial graphs, identifying real zeroes, and determine the degree and number of real zeroes based on polynomial properties.
\end{enumerate}

\subsection{Assessment Instrument}

\begin{enumerate}
    \item \textbf{Multiple Choice:}
    \item Define a polynomial.
    \item A) An expression with one term
    \item B) An expression with two terms
    \item C) An algebraic expression with one or more terms
    \item D) A sum of fractions
    \item \textbf{True/False:}
    \item A polynomial of degree 3 can have up to 3 real zeroes.
    \item True
    \item False
    \item \textbf{Short Answer:}
    \item Given the polynomial \( P(x) = 3x^2 - 12x + 9 \), factor the polynomial.
    \item \textbf{Problem Solving:}
    \item Solve the quadratic equation \( x^2 - 5x + 6 = 0 \) and find the roots.
    \item \textbf{Graph Analysis:}
    \item Graph the polynomial \( P(x) = x^3 - 3x^2 - 4x + 12 \) and identify the number of real zeroes.
\end{enumerate}

%%===============================================================
\section{Factoring Polynomials}
\label{sec:factoring-polynomials}
%%===============================================================

\paragraph{Factoring Polynomials}

\paragraph{Useful Formulas}

Computing factors of polynomials requires knowledge of different formulas and some experience to find out which formula to apply. Below, we provide some important formulas:

\begin{itemize}
    \item $x^2 - y^2 = (x + y)(x - y)$
    \item $x^2 + 2xy + y^2 = (x + y)^2$
    \item $x^2 - 2xy + y^2 = (x - y)^2$
    \item $x^3 - y^3 = (x - y)(x^2 + xy + y^2)$
    \item $x^3 + y^3 = (x + y)(x^2 - xy + y^2)$
    \item $a^3 + b^3 + c^3 - 3abc = (a + b + c)(a^2 + b^2 + c^2 - ab - ac - bc)$
\end{itemize}

\paragraph{Methods}

Given the expression $x^2 + 3x + 2$, one may ask "What are the values of $x$ that make this expression 0?" If we factor, we obtain:

$$
x^2 + 3x + 2 = (x + 2)(x + 1)
$$

If $x = -1$ or $x = -2$, then one of the factors on the right becomes zero. Therefore, the whole expression must be zero. By factoring, we have discovered the values of $x$ that render the expression zero. These values are termed "roots."

In general, given a quadratic polynomial $px^2 + qx + r$ that factors as:

$$
px^2 + qx + r = (ax + c)(bx + d)
$$

Then, the roots of the polynomial are:

$$
x = -\frac{c}{a} \text{ and } x = -\frac{d}{b}
$$

A special case to be aware of is the difference of two squares, $a^2 - b^2$. In this case, we can always factor as:

$$
a^2 - b^2 = (a + b)(a - b)
$$

For example, consider $4x^2 - 9$. On initial inspection, we see that both $4x^2$ and $9$ are squares of $2x$ and $3$, respectively. Applying the previous rule, we have:

$$
4x^2 - 9 = (2x + 3)(2x - 3)
$$

\paragraph{The AC Method}

The AC method simplifies the process of factoring. Suppose a quadratic polynomial has the form:

$$
ax^2 + bx + c
$$

If there are numbers $r$ and $q$ that satisfy:

$$
r \cdot q = ac
$$
$$
r + q = b
$$

Then, we can rewrite the polynomial as:

$$
ax^2 + bx + c = ax^2 + rx + qx + c
$$

And factor out a common term from $(ax^2 + rx)$ and $(qx + c)$ to factor the polynomial.

For example, take the polynomial $2x^2 + 7x - 15$:

$$
ac = 2 \cdot (-15) = -30
$$
$$
10 \cdot (-3) = -30 \text{ and } 10 + (-3) = 7
$$

So we can rewrite the polynomial as:

$$
2x^2 + 10x - 3x - 15
$$

From here, we can factor $2x$ from the first two terms and $-3$ from the last two, which gives us:

$$
2x(x + 5) - 3(x + 5) = (2x - 3)(x + 5)
$$

Thus, the factored form of $2x^2 + 7x - 15$ is:

$$
(2x - 3)(x + 5)
$$

\paragraph{The Quadratic Formula}

Given any quadratic equation $ax^2 + bx + c = 0$, where $a \neq 0$, all solutions are given by the quadratic formula:

$$
x = \frac{-b \pm \sqrt{b^2 - 4ac}}{2a}
$$

Note that the value of $b^2 - 4ac$ will affect the number of real solutions:

\begin{itemize}
    \item If $b^2 - 4ac > 0$, there are two real solutions.
    \item If $b^2 - 4ac = 0$, there is one real solution.
    \item If $b^2 - 4ac < 0$, there are no real solutions.
\end{itemize}

\paragraph{Remainder and Factor Theorem}

The polynomial division algorithm is as follows: Suppose $d(x)$ and $p(x)$ are nonzero polynomials where the degree of $p(x)$ is greater than or equal to the degree of $d(x)$. Then there exist two unique polynomials, $q(x)$ and $r(x)$, such that:

$$
p(x) = d(x)q(x) + r(x)
$$

Where either $r(x) = 0$ or the degree of $r(x)$ is strictly less than the degree of $d(x)$.

\paragraph{Remainder Theorem}

Suppose $p(x)$ is a polynomial of degree at least 1 and $c$ is a real number. When $p(x)$ is divided by $(x - c)$, the remainder is $p(c)$.

\textbf{Proof:} By the division algorithm:

$$
p(x) = (x - c)q(x) + r
$$

Where $r$ must be a constant since $d(x) = x - c$ has a degree of 1. Setting $x = c$, we get:

$$
p(c) = (c - c)q(x) + r = r
$$

Thus, the remainder $r = p(c)$.

\paragraph{Factor Theorem}

Suppose $p(x)$ is a nonzero polynomial. The real number $c$ is a zero of $p(x)$ if and only if $(x - c)$ is a factor of $p(x)$.

By the division algorithm, $x - c$ is a factor of $p(x)$ if and only if $r = 0$. Since $p(c) = r$ when $p(x)$ is divided by $x - c$, $x - c$ is a factor of $p(x)$ if and only if $p(c) = 0$; that is, if $c$ is a zero of $p(x)$.

\paragraph{Factor Theorem Example Problem}

Determine if $x + 2$ is a factor of $2x^2 + 3x - 2$.

Since $c$ is positive instead of negative, use this basic identity:

$$
x + 2 = x - (-2)
$$

Now use the factor theorem:

$$
2(-2)^2 + 3(-2) - 2 = 8 - 6 - 2 = 0
$$

Since the result is 0, $(x + 2)$ is a factor of $2x^2 + 3x - 2$.

Thus, it can be re-stated as:

$$
2x^2 + 3x - 2 = (x + 2)(ax + b)
$$

Expanding the right-hand side:

$$
2x^2 + 3x - 2 = ax^2 + x(2a + b) + 2b
$$

Equating like terms, we get:

$$
2 = a
$$
$$
2a + b = 3
$$
$$
2b = -2
$$

Solving these, $a = 2$ and $b = -1$. Thus:

$$
2x^2 + 3x - 2 = (x + 2)(2x - 1)
$$

\paragraph{Example Problems}

\textbf{EXAMPLE 1:} Find all the roots of $4x^2 + 7x - 2$.

Finding the roots is equivalent to solving $4x^2 + 7x - 2 = 0$. Applying the quadratic formula with $a = 4$, $b = 7$, $c = -2$, we have:

$$
x = \frac{-7 \pm \sqrt{7^2 - 4 \cdot 4 \cdot (-2)}}{2 \cdot 4}
$$
$$
x = \frac{-7 \pm \sqrt{49 + 32}}{8}
$$
$$
x = \frac{-7 \pm \sqrt{81}}{8}
$$
$$
x = \frac{-7 \pm 9}{8}
$$
$$
x = \frac{2}{8} \text{ or } x = \frac{-16}{8}
$$
$$
x = \frac{1}{4} \text{ or } x = -2
$$

\textbf{EXAMPLE 2:} Factor the polynomial $4x^2 + 7x - 2$.

We know the roots are $x = \frac{1}{4}$ and $x = -2$. Therefore, the polynomial factors as:

$$
4x^2 + 7x - 2 = 4(x - \frac{1}{4})(x + 2)
$$

Or:

$$
4x^2 + 7x - 2 = (4x - 1)(x + 2)
$$

\textbf{EXAMPLE 3:} Factor $4x^2 + 12x + 9$ using the perfect square trinomial method.

We observe that $4x^2 + 12x + 9$ can be rewritten as:

$$
(2x + 3)^2
$$

To verify, expand $(2x + 3)^2$:

$$
(2x + 3)^2 = 4x^2 + 12x + 9
$$

\paragraph{Resources}

\begin{itemize}
    \item \href{https://tutorial.math.lamar.edu/classes/alg/factoring.aspx}{Factoring Polynomials with Example Problems, Paul''s Online Notes}
    \item \href{https://courses.lumenlearning.com/suny-beginalgebra/chapter/5-2-1-factor-trinomials/}{Factoring Polynomials, Lumen Learning}
    \item \href{https://www.youtube.com/watch?v=shHS-ERFVrc}{Difference of Squares. Produced by TA Catherine Sporer, UTSA}
    \item \href{http://www.algebralab.org/lessons/lesson.aspx?file=algebra_factoring.xml}{6 Methods of Factoring with Practice Problems, AlgebraLAB}
    \item \href{https://txwes.edu/media/twu/content-assets/images/academics/academic-success-center/A-C-Method.pdf}{AC Factoring Method, Texas Wesleyan University}
\end{itemize}

\paragraph{Licensing}

Content obtained and/or adapted from:

\begin{itemize}
    \item \href{https://en.wikibooks.org/wiki/Calculus/Algebra}{Algebra, Wikibooks: Calculus} under a CC BY-SA license
    \item \href{https://en.wikibooks.org/wiki/Algebra/Factoring_Polynomials}{Factoring Polynomials, Wikibooks: Algebra} under a CC BY-SA license
\end{itemize}

\subsection{Learning Objectives}

\begin{enumerate}
    \item \textbf{1. Identify and Apply Formulas:}
    \item Students will be able to identify and apply key factoring formulas, such as the difference of squares, perfect square trinomials, and the sum and difference of cubes, to factor polynomials. \textbf{2. Utilize the AC Method:}
    \item Students will demonstrate proficiency in using the AC method to factor quadratic polynomials by identifying appropriate values for `r` and `q` that satisfy the conditions for rewriting and factoring the polynomial. \textbf{3. Solve Quadratic Equations:}
    \item Students will use the quadratic formula to solve quadratic equations and understand how the discriminant affects the number of real solutions, ensuring accurate factorization and root identification.
\end{enumerate}

\subsection{Assessment Instrument}

\begin{enumerate}
    \item Instructions Complete the following problems to demonstrate your understanding of factoring polynomials. Show all work and provide clear explanations for each step. Part A: Factor Polynomials
    \item \textbf{Factor the following expressions completely:} a. $x^2 - 16$ b. $x^2 + 6x + 9$ c. $x^3 - 27$ d. $x^3 + 8$
    \item \textbf{Use the AC method to factor the following quadratic polynomial:} a. $6x^2 + 11x + 3$
    \item \textbf{Determine if the following expressions are factors of the given polynomials using the Remainder Theorem:} a. Is $x + 3$ a factor of $2x^3 + 5x^2 - x - 6$? b. Is $x - 2$ a factor of $x^3 - 3x^2 - 4x + 12$? Part B: Solve and Interpret
    \item \textbf{Solve the quadratic equations using the quadratic formula and interpret the number of real solutions based on the discriminant:} a. $2x^2 - 4x - 6 = 0$ b. $x^2 + 4x + 4 = 0$ c. $3x^2 - 5x + 2 = 0$
    \item \textbf{Given the polynomial $4x^2 + 7x - 2$, factor the polynomial and verify your result by expanding the factored form.} Part C: Application and Explanation
    \item \textbf{Explain the difference between the difference of squares and perfect square trinomials, providing examples of each and how they are factored.}
    \item \textbf{Describe the process of using the quadratic formula, including how the discriminant affects the solutions of the quadratic equation.} Evaluation Criteria
    \item \textbf{Accuracy:} Correctly factoring polynomials and solving quadratic equations.
    \item \textbf{Clarity:} Clear and logical presentation of steps and explanations.
    \item \textbf{Understanding:} Demonstrating comprehension of factoring methods and the quadratic formula.
\end{enumerate}

%%===============================================================
\section{Rational Expressions}
\label{sec:rational-expressions}
%%===============================================================

\paragraph{Rational Expression}

Rational expressions are fractions that have a polynomial in the numerator, denominator, or both. Although rational expressions can seem complicated because they contain variables, they can be simplified using techniques similar to those used for simplifying expressions such as 

$$
\frac{4x^3}{12x^2}
$$

combined with techniques for factoring polynomials.

\paragraph{Simplifying Rational Expressions}

To simplify a rational expression, follow these steps:

1. \textbf{Determine the domain.} The excluded values are those values for the variable that result in the expression having a denominator of $0$.
2. \textbf{Factor the numerator and denominator.}
3. \textbf{Find common factors for the numerator and denominator and simplify.}

Consider two polynomials

$$
p(x) = a\_n x^n + a\_{n-1} x^{n-1} + \cdots + a\_1 x + a\_0
$$

and

$$
q(x) = b\_m x^m + b\_{m-1} x^{m-1} + \cdots + b\_1 x + b\_0
$$

When we take the quotient of the two we obtain

$$
\frac{p(x)}{q(x)} = \frac{a\_n x^n + a\_{n-1} x^{n-1} + \cdots + a\_1 x + a\_0}{b\_m x^m + b\_{m-1} x^{m-1} + \cdots + b\_1 x + b\_0}
$$

The ratio of two polynomials is called a rational expression. Many times we would like to simplify such a beast. For example, say we are given

$$
\frac{x^2 - 1}{x + 1}
$$

We may simplify this in the following way:

$$
\frac{x^2 - 1}{x + 1} = \frac{(x + 1)(x - 1)}{x + 1} = x - 1, \quad x \neq -1
$$

This is nice because we have obtained something we understand quite well, $x - 1$, from something we didn''t.

\paragraph{Resources}

\begin{itemize}
    \item \href{https://courses.lumenlearning.com/beginalgebra/chapter/6-1-1-introduction-to-rational-expressions/}{Identify and Simplify Rational Expressions, Lumen Learning}
    \item \href{https://www.youtube.com/watch?v=uVpsz-xpnPo}{Simplifying Rational Expressions, The Organic Chemistry Tutor}
\end{itemize}

\paragraph{Licensing}

Content obtained and/or adapted from:

\begin{itemize}
    \item \href{https://en.wikibooks.org/wiki/Calculus/Algebra}{Algebra, Wikibooks: Calculus} under a CC BY-SA license
\end{itemize}

\subsection{Learning Objectives}

\begin{enumerate}
    \item \textbf{Factor Rational Expressions:} Students will be able to factor polynomials in the numerator and denominator of a rational expression to simplify it.
    \item \textbf{Determine Domain:} Students will determine the domain of a rational expression by identifying values that make the denominator zero.
    \item \textbf{Apply Simplification Techniques:} Students will simplify rational expressions by canceling common factors and explaining the process used.
\end{enumerate}

\subsection{Assessment Instrument}

\begin{enumerate}
    \item Instructions Complete the following problems to demonstrate your understanding of simplifying rational expressions. Show all work and provide clear explanations for each step. Part A: Simplify Rational Expressions
    \item \textbf{Simplify the following rational expressions:} a. $ \frac{3x^2 - 12}{3x} $ b. $ \frac{x^2 - 4}{x^2 - 2x} $ c. $ \frac{2x^2 + 8x}{4x^2 + 16x} $ d. $ \frac{x^3 - 8}{x^2 - 4x + 4} $
    \item \textbf{Determine the domain of the following rational expressions:} a. $ \frac{5x}{x^2 - 9} $ b. $ \frac{x^2 - 1}{x^2 + x - 6} $ Part B: Apply and Interpret
    \item \textbf{Simplify the following rational expressions and interpret any domain restrictions:} a. $ \frac{x^2 + 6x + 9}{x + 3} $ b. $ \frac{x^2 - 4}{x - 2} $
    \item \textbf{Explain how you simplified each expression and any restrictions on the variable:} a. $ \frac{x^2 - 16}{x^2 - 4x} $ b. $ \frac{2x^2 - 10x}{4x^2 - 20x} $ Evaluation Criteria
    \item \textbf{Accuracy:} Correctly simplifying rational expressions and identifying domain restrictions.
    \item \textbf{Clarity:} Clear and logical presentation of steps and explanations.
    \item \textbf{Understanding:} Demonstrating comprehension of factoring, simplifying, and domain restrictions.
\end{enumerate}

%%===============================================================
\section{Graph of equations}
\label{sec:graph-of-equations}
%%===============================================================

\paragraph{Graphs}

$y = 2x$
% [Image: Graph of y = 2x]
 

It is sometimes difficult to understand the behavior of a function given only its definition; a visual representation or graph can be very helpful. A graph is a set of points in the Cartesian plane, where each point $(x, y)$ indicates that $f(x) = y$. In other words, a graph uses the position of a point in one direction (the vertical-axis or $y$-axis) to indicate the value of $f$ for a position of the point in the other direction (the horizontal-axis or $x$-axis).

Functions may be graphed by finding the value of $f$ for various $x$ and plotting the points $(x, f(x))$ in a Cartesian plane. For the functions that you will deal with, the parts of the function between the points can generally be approximated by drawing a line or curve between the points. Extending the function beyond the set of points is also possible, but becomes increasingly inaccurate.

\paragraph{Linear functions}

Graphing linear functions is easy to understand and do. Because we know that two points can form a line, only two points are needed for us to graph a linear function if those two points are on the function. Oppositely, we can write down the equation of a linear function if we only know two points that are on the function.

The following section mainly talks about different forms of linear function notations so that you can easily identify or graph the function.

\paragraph{Introduction}

Plotting points like this is laborious. Fortunately, many functions'''''''' graphs fall into general patterns. For a simple case, consider functions of the form

$$
f(x) = 3x + 2
$$

The graph of $f$ is a single line, passing through the point $(0, 2)$ with slope 3. Thus, after plotting the point, a straightedge may be used to draw the graph. This type of function is called linear and there are a few different ways to present a function of this type.

\paragraph{Slope}

The slope is the backbone of linear functions because it shows how much the output of a function changes when the input changes. For example, if the slope of a function is 2, then it means when the input of a function increases by 1 unit, the output of the function increases by 2 units. Now, let''''''''s look at a more mathematical example.

Consider this function: 

$$
f(x) = -5x + 8
$$

What does the number $-5$ mean?

It means that when $x$ increases by 1, $f(x)$ decreases by 5.

Using mathematical terms:

$$
f(x+1) - f(x) = (-5(x+1) + 8) - (-5x + 8) = -5
$$

It is easy to calculate the slope because the slope is like the speed of a vehicle. If we divide the change in distance and the corresponding change in time, we get the speed. Similarly, if we divide the change in $f(x)$ over the corresponding change in $x$, we get the slope. If given two points, $(x\_1, y\_1)$ and $(x\_2, y\_2)$, we may then compute the slope of the line that passes through these two points. Remember, the slope is determined as "rise over run." That is, the slope is the change in $y$-values divided by the change in $x$-values. In symbols:

$$
\text{slope} = \frac{\text{change in } y}{\text{change in } x} = \frac{\Delta y}{\Delta x} = \frac{y\_2 - y\_1}{x\_2 - x\_1}
$$

Slope in a linear function: If two points $(x\_1, y\_1)$ and $(x\_2, y\_2)$ are on a linear function, then the slope of the linear function is

$$
m = \frac{y\_2 - y\_1}{x\_2 - x\_1}
$$

Interestingly, there is a subtle relationship between the slope and the angle between the graph of the function and the positive $x$-axis, $\theta$. The relationship is:

$$
m = \tan \theta
$$

It is an obvious relationship, but it can be ignored relatively easily.

\paragraph{Slope-intercept form}

When we see a function presented as

$$
y = f(x) = mx + b
$$

we call this presentation the slope-intercept form. This is because, not surprisingly, this way of writing a linear function involves the slope, $m$, and the $y$-intercept, $b$.

\paragraph{Example 1}

Graph the function 

$$
f(x) = 3x + 2
$$

The slope of the function is 3, and it intercepts the $y$-axis at point $(0, 2)$. In order to graph the function, we need another point. Since the slope of the function is 3, then

$$
f(x + 1) = f(x) + 3
$$

$$
f(1) = f(0 + 1) = f(0) + 3 = 5
$$

Knowing that the function goes through points $(0, 2)$ and $(1, 5)$, the function can be easily graphed.

\paragraph{Example 2}

Now, consider another unknown linear function that goes through points $(-3, -4)$ and $(0, 5)$. What is the equation for this function?

The slope can be calculated with the formula mentioned above.

$$
m = \frac{y\_2 - y\_1}{x\_2 - x\_1} = \frac{5 - (-4)}{0 - (-3)} = 3
$$

And since the $y$-axis interception is $(0, 5)$, we can know that

$$
b = 5
$$

Thus, the equation of this linear function should be

$$
g(x) = 3x + 5
$$

\paragraph{Point-slope form}

If someone walks up to you and gives you one point and a slope, you can draw one line and only one line that goes through that point and has that slope. Said differently, a point and a slope uniquely determine a line. So, if given a point $(x\_0, y\_0)$ and a slope $m$, we present the graph as

$$
y - y\_0 = m(x - x\_0)
$$

We call this presentation the point-slope form. The point-slope and slope-intercept form are essentially the same. In the point-slope form, we can use any point the graph passes through. Whereas, in the slope-intercept form, we use the $y$-intercept, that is the point $(0, b)$. The point-slope form is very important. Although it is not used as frequently as its counterpart, the slope-intercept form, the concept of knowing a point and drawing the line in the direction of the slope will be encountered when we go into vector equations for lines and planes in future chapters.

\paragraph{Example 1}

If a linear function goes through points $(3, -4)$ and $(1, 5)$, what is the equation for this function?

The slope is:

$$
m = \frac{y\_2 - y\_1}{x\_2 - x\_1} = \frac{5 - (-4)}{1 - 3} = -\frac{9}{2}
$$

Since we know two points, the following answers are all correct:

$$
y + 4 = -\frac{9}{2}(x - 3)
$$

or

$$
y - 5 = -\frac{9}{2}(x - 1)
$$

The two-point form is another form to write the equation for a linear function. It is similar to the point-slope form. Given points $(x\_1, y\_1)$ and $(x\_2, y\_2)$, we have the equation

$$
y - y_1 = \frac{y\_2 - y\_1}{x\_2 - x\_1}(x - x\_1)
$$

This presentation is in the two-point form. It is essentially the same as the point-slope form except we substitute the expression $\frac{y\_2 - y\_1}{x\_2 - x\_1}$ for $m$. However, this expression is not widely used in mathematics because in most practical cases, the slope or $y$-intercept is known.

\paragraph{Intercept form}

Finally, another form is the intercept form. In this form, the equation of the line is presented in terms of $a$ and $b$, where $a$ and $b$ are the $x$-intercept and $y$-intercept, respectively.

$$
\frac{x}{a} + \frac{y}{b} = 1
$$

where the point $(a, 0)$ and $(0, b)$ are both on the line.

\paragraph{Example 1}

To find the intercept form of the equation for a line passing through $(1, 3)$ and $(4, 7)$.

The $x$-intercept is found by setting $y = 0$ in the equation. Substituting $y = 0$ into the general equation of a line

$$
y = m x + b
$$

we get:

$$
0 = m x + b
$$

Thus, the $x$-intercept is given by

$$
a = - \frac{b}{m}
$$

Similarly, to find the $y$-intercept, we set $x = 0$ and solve for $b$.

\paragraph{Summary}

The main takeaways from this section are:

1. \textbf{Linear Functions}: The graph of a linear function is always a straight line. Its slope determines how steep the line is, and its $y$-intercept determines where it crosses the $y$-axis.

2. \textbf{Forms of Linear Functions}: Several forms represent linear functions, including slope-intercept, point-slope, two-point, and intercept forms. Each provides a different way to express or work with linear functions.

3. \textbf{Graphing}: With the right form, you can easily graph a linear function using just a few key points or parameters.

\paragraph{Quadratic functions}

Quadratic functions are functions of the form

$$
f(x) = ax^2 + bx + c
$$

where $a$, $b$, and $c$ are constants. They produce parabolas when graphed. The shape and position of the parabola depend on the values of $a$, $b$, and $c$. Understanding the forms and characteristics of quadratic functions is crucial for solving various mathematical problems.

\paragraph{Standard form}

The standard form of a quadratic function is

$$
f(x) = ax^2 + bx + c
$$

where $a$, $b$, and $c$ are constants. The parabola opens upwards if $a > 0$ and downwards if $a < 0$. The vertex form can also be used to provide a more intuitive understanding of the function''''''''s graph.

\paragraph{Vertex form}

The vertex form of a quadratic function is

$$
f(x) = a(x - h)^2 + k
$$

where $(h, k)$ is the vertex of the parabola. This form is particularly useful for graphing because it directly shows the vertex and the direction of the parabola.

\paragraph{Factored form}

The factored form of a quadratic function is

$$
f(x) = a(x - r\_1)(x - r\_2)
$$

where $r\_1$ and $r\_2$ are the roots of the quadratic equation. This form is useful for finding the roots of the quadratic function and understanding the x-intercepts of the graph.

\paragraph{Exponential and Logarithmic functions}

Exponential and logarithmic functions are inverses of each other and are used to model a wide range of real-world phenomena. Exponential functions have the general form

$$
f(x) = a \cdot b^x
$$

where $a$ is a constant and $b$ is the base of the exponential function. Logarithmic functions have the general form

$$
f(x) = a \cdot \log\_b(x)
$$

where $a$ is a constant and $b$ is the base of the logarithm. Understanding these functions is important for solving problems involving growth and decay, among other applications.

\paragraph{Resources}

\paragraph{Basic Information}

\begin{itemize}
    \item \href{https://www.mathsisfun.com/data/cartesian-coordinates.html}{Introduction to Cartesian Coordinates and Graphs, mathisfun.com}
    \item \href{https://www.khanacademy.org/math/algebra/x2f8bb11595b61c86:foundation-algebra/x2f8bb11595b61c86:algebra-overview-history/v/descartes-and-cartesian-coordinates}{Introduction to the Coordinate Plane, Khan Academy}
    \item \href{https://courses.lumenlearning.com/ivytech-collegealgebra/chapter/plotting-ordered-pairs-in-the-cartesian-coordinate-system/}{Plotting Points on a Graph, Lumen Learning}
\end{itemize}

\paragraph{Graphing Linear Functions}

\begin{itemize}
    \item \href{https://courses.lumenlearning.com/waymakercollegealgebra/chapter/graph-linear-functions/}{Graphing Linear Functions, Lumen Learning}
    \item \href{https://www.khanacademy.org/math/algebra/x2f8bb11595b61c86:linear-equations-graphs}{Linear Equations and Graphs, Khan Academy}
    \item \href{https://www.palmbeachstate.edu/prepmathlw/Documents/graphing\%20linear\%20equations.pdf}{Graphing Linear Equations, Palm Beach State College}
    \item \href{https://www.youtube.com/watch?v=lkRyxnqOb8M}{Graphing a Line Using a Point and Slope, patrickJMT}
    \item \href{https://www.youtube.com/watch?v=mxBoni8N70Y}{Graphing a Line with X and Y Intercepts, patrickJMT}
\end{itemize}

\paragraph{Licensing}

Content obtained and/or adapted from:

\begin{itemize}
    \item \href{https://en.wikibooks.org/wiki/Calculus/Graphing_functions}{Graphing functions, Wikibooks: Calculus} under a CC BY-SA license
\end{itemize}

\subsection{Learning Objectives}

\begin{enumerate}
    \item \textbf{Identify and Graph Linear Functions}: Students will be able to identify linear functions in slope-intercept, point-slope, and intercept forms and accurately graph them on a Cartesian plane.
    \item \textbf{Analyze and Interpret Quadratic Functions}: Students will be able to convert quadratic functions between standard, vertex, and factored forms and interpret the features of their graphs, including vertices and intercepts.
    \item \textbf{Apply Exponential and Logarithmic Functions}: Students will be able to model real-world problems using exponential and logarithmic functions, and solve problems involving growth, decay, and inverse relationships.
\end{enumerate}

\subsection{Assessment Instrument}

\begin{enumerate}
    \item Assessment Instrument \textbf{1. Graphing Linear Functions}
    \item \textbf{Task}: Given the linear function $f(x) = -2x + 4$, graph the function using slope-intercept form. Label the $y$-intercept and another point on the line. \textbf{2. Converting and Analyzing Quadratic Functions}
    \item \textbf{Task}: Convert the quadratic function $f(x) = 2x^2 - 8x + 6$ from standard form to vertex form and identify the vertex and axis of symmetry. \textbf{3. Solving Exponential and Logarithmic Problems}
    \item \textbf{Task}: Solve the exponential equation $2^x = 16$ and the logarithmic equation $\log_2(x) = 5$. Provide the solutions and explain the steps.
\end{enumerate}

%%===============================================================
\section{Linear Equations}
\label{sec:linear-equations}
%%===============================================================

\paragraph{Linear Equations}



\paragraph{What are Linear Equations?}



In the functions section we talk about how a function is like a box that takes an independent input value and uses a rule defined mathematically to create a unique output value. The value for the output is dependent on the value that is put in the box. We call the values that are going into the box the independent values or the domain. We call the values coming out of the box the dependent values or the range.



Unless we specify differently, on Cartesian graphs the domain is the real numbers. In the Cartesian Plane section we saw how running different values through a function to identify the points on the Cartesian plane by picking the first (x) value of the point the domain and the second (y) value from the range. To restate this: by convention the two variables for a function on the Cartesian Plane are x for the domain, the independent variable, and y for the range, the dependent variable. The variable y is the same as writing f(x). Mathematicians recognize this equivalence but generally prefer to write y because it is shorter. Because a function definition has an input and an output it must also contain an equal sign. The section in this book solving equations showed the various operations we can perform on both sides of an equal sign and still maintain the notion of equivalence. In this section we plug different values into the independent variable and solve to find the associated dependent variable. For instance if we start with the equation:



$$y = x$$



$$ \text{We can add a } -x \text{ to both sides to get the equation}$$



$$y - x = x - x$$



$$ \text{which we then simplify to}$$



$$y - x = 0$$



$$ \text{Or we can add a } -y \text{ to both sides to get the equation}$$



$$y - y = x - y$$



$$ \text{which we then simplify to}$$



$$0 = x - y$$



$$ \text{Since } y - x = 0  \equiv 0 = x - y \text{ then:}$$



$$y - x = 0 = x - y$$



$$ \text{Using the transitive property}$$



$$y - x  = x - y$$



$$ \text{adding } x + y \text{ to both sides gives us}$$



$$(y - x )+( y + x)  = (x - y) + (y + x)$$



$$ \text{using the associative property we change this to}$$



$$(y - y )+( x + x)  = (x - x) + (y + y)$$



$$ \text{which simplifies to}$$



$$2x  = 2y$$



$$ \text{And divide both sides by } 2$$



$$\frac{2x}{2}  = \frac{2y}{2}$$



$$ \text{To simply reverse the order and show}$$



$$x = y$$



We have not really proved anything mathematically above, but these operations allow us to manipulate equations to get the dependent variable by itself on one side of the equals sign. Then we can plug numbers into the independent variable to discover the function values for those numbers. Then we can draw these values as points on the Cartesian plane and get a feel for what the function would look like if we could see all the points defined by the function at once.



\paragraph{Horizontal Linear Equations}



Equations of the form $y = C\_2$ are linear functions of the general form $y = mx + b$ where slope $m = 0$ and the constant $C\_2$ is the y-intercept $b$ (in the general form). The graph of this zero-slope function is a straight horizontal line, intercepting the y-axis at $C\_2$, including zero and extends infinitely in the positive and negative directions for all $R$ values of x.



% [Image: graph1]



The domain for such functions is $R$ covering all real numbers (unless otherwise specified), but the range is just the set $\{ c \}$. The equation $y = 0$ is the x-axis.



\paragraph{Vertical Linear Equations}



Equation $x = C\_1$, x is one single value $C\_1$ and y, being unrestricted, is every $R$ number. The graph of $x = C\_1$ is a straight vertical line where $x = C\_1$, covering all positive, negative and zero values of y.



% [Image: x = c, where c = 3]



Its domain is set $\{ C\_1 \}$ and range is set $R$ (unless otherwise specified). $x = C\_1$ is technically not a function (there is more than one value of y for each value of x), but it''s a relation. Vertical lines have no slope ($m =$ divide by zero, undefined, plus and minus infinity). These are the only types of linear equations of the general form shown previously which are not linear functions. The equation $x = 0$ is exactly the y-axis. Lines with steepness approaching vertical have very large-magnitude slopes but are still functions.



\paragraph{General Linear Equations}



We are going to start by looking at simple functions called linear equations. When none of the instances of x and y in the algebraic expression defining the function rule have exponents then all the instances of x and y can be combined into just two occurrences. The graph of the expression can be represented as a straight line. The equation that expresses the function is considered a linear equation with two variables. The following equation is a simple example of such a linear equation:



$$y - x = 2$$



Since y is the dependent variable it is standing in for the function. We can re-write the expression as $f(x) - x = 2$. If we add an x to both sides the equality property holds and we get the expression $f(x) - x + x = x + 2$. Simplifying we get $f(x) = x + 2$. In the following table we''ll pick 3 values for x, and then calculate the dependent (y) values from f(x).



$$

\begin{array}{|c|c|c|}

\hline

\text{x value} & \text{y value (abscissa)} & \text{Coordinates (x,y)} \\\\ 

\hline

-1 & 1 & (-1, 1) \\\\ 

\hline

0 & 2 & (0, 2) \\\\ 

\hline

1 & 3 & (1, 3) \\\\ 

\hline

\end{array}

$$





where x and y are variables to be plotted in a two-dimensional Cartesian coordinate graph as shown here:



% [Image: 3ptsyeqxpl2.png]



% [Image: Y equals x plus 2.PNG]



This function is equivalent to the previous example of a linear equation, $y - x = 2$. The arrows at each end of the line indicate that the line extends infinitely in both directions. All linear functions of a single input variable have or can be algebraically arranged to have the general form:



$$ y = f(x) = mx + b $$



where x and y are variables, f(x) is the function of x, m is a constant called the slope of the line, and b is a constant which is the ordinate of the y-intercept (i.e. the value of y where the function line crosses the y-axis). The slope indicates the steepness of the line. In the previous example where $y = x + 2$, the slope $m = 1$ and the y-intercept ordinate $b = 2$. The $y = mx + b$ form of a linear function is called the slope-intercept form.



Unless a domain for x is otherwise stated, the domain for linear functions will be assumed to be all real numbers and so the lines in graphs of all linear functions extend infinitely in both directions. Also in linear functions with all real number domains, the range of a linear function will cover the entire set of real numbers for y, unless the slope $m = 0$ and the function equals a constant. In such cases, the range is simply the constant.



Conversely, if a function has the general form $y = mx + b$ or if it can be arranged to have that form, the function is linear. A linear equation with two variables has or can be algebraically rearranged to have the general form:



$$ A x + B y = C $$



where x and y represent the linear variables, and the letters A, B, and C can represent any real constants, either positive or negative. Conversely, an equation with two variables x and y having that general form, or being able to be arranged in that form, would be linear as long as A and B are not both equal to 0. In the preceding equation, capital letters are to avoid confusion with other constants in this chapter and for consistency with Reference 1.



If one divides the preceding equation by B (when B is not 0) and solves for y, the following form can be obtained:



$$ y = \frac{-A}{B} x + \frac{C}{B} $$



If one equates $-A/B$ to the slope m and $C/B$ to the y-intercept ordinate b, it can be seen that the general form for a linear equation and the slope-intercept form for a linear function are practically interconvertible except for the fact that, in a linear function, the B constant in the linear equation form cannot equal 0.



\paragraph{X and Y axis intercepts}



An axis intercept point is a point where the graph of a function, relation, or equation intersects the X or Y axes. This section is about finding out how a particular set of functions: linear functions cross the axes.



We know that the domains of most lines are infinite because they are defined at every value of $X$. The exception is lines that are defined as $X = c$ where $c$ is a number we choose when we write the function. By definition, this line is only defined for one value of $X$. Since the domain maps onto more than one value for the range, this is actually a relationship and not a function. We''ve seen the graph of this relationship is a vertical line that passes through the point $(c,Y)$.



Lines with the equation $X = c$ intersect the X axis once and the Y axis 0 or infinite times.



We can also restrict the range of a function by simply writing $Y = c$, where $c$ is again any number we choose. The graph of this line is a horizontal line that passes through the point $(X, c)$. When $Y = 0$, this line is the same as the X axis. When $c \neq 0$, then the line can never intercept the X axis.



Lines with the equation $Y = c$ intersect the Y axis once and the X axis 0 or infinite times.



When looking at Cartesian graphs and linear equations, we run into a mathematical axiom: "Two points determine a line." We will see how this axiom affects the slope-intercept definition of a line $y = f(x) = mx + b$ in the next section. When two lines intersect, they intersect at a point. If a line is not horizontal or perpendicular, it will have to intersect the X and Y axes once, but only once.



In this page, we are going to accept the statement "At most one line can be drawn through any point not on a given line parallel to the given line in a plane."



We''ve seen that for the equation $Y = mX + b$, the Y intercept will always be at $b$ because that is where $X = 0$.



Using Algebra, we can subtract $b$ from both sides: $Y - b = mX$



and multiply by $\frac{1}{m}$



$\frac{Y}{m} - \frac{b}{m} = X$



we can see that the X intercept is going to be $-\frac{b}{m}$.



An axis intercept may simply refer to the number value on the axis where the intersection occurs. For brevity, we may say the line has an X intercept of 1 and a Y intercept of 2. After graphing just a few lines, you will be able to tell this line points down and runs through quadrants II, I, and IV. With a little more practice, you will be able to know that the equation for the line is $Y = -2x + 2$. We will see that by specifying the two points, we are actually implying the slope of the line. There is an exception to this rule. If we say a line crosses the axes at 0, we know that the line will pass through 2 quadrants instead of 3, but we won''t know which quadrants or how steep the line is. When we look at slope in the next section, we will see why the equations above specify a point and a slope.



When you are trying to graph a linear equation, finding the axes intercepts is often the easiest way to go about doing it. To find the x-intercept, set $y = 0$ and solve for $x$. To find the y-intercept, set $x = 0$ and solve for $y$. For most examples, the intercepts are different points, and a line can be drawn through the two intercepts. If both intercepts are $(0, 0)$, then another point must be determined to graph the line. If the equation is in the form $x = c$ or $y = c$, the horizontal or vertical lines are very simple to plot.



\paragraph{Slope}



\paragraph{Algebra/Slope}



Slope is the change in the vertical distance of a line on a coordinate plane over the change in horizontal difference. In other words, it is the “rise” over the “run” or the steepness of a line. Slope is usually represented by the symbol $m$ like in the equation $y = mx + b$, where $m$, the coefficient of $x$, represents the slope of the line.



Slope is computed by measuring the change in vertical distance divided by the change in horizontal difference, i.e.:



$$

m = \frac{\Delta y}{\Delta x} = \frac{\text{rise}}{\text{run}}

$$



The Greek uppercase letter $\Delta$ represents change, in this case, change in the $y$-coordinates divided by the change in $x$-coordinates.



\paragraph{Positive Slope/ Negative Slope}



If a line goes up from left to right, then the slope has to be positive. For example, a slope of $\frac{3}{4}$ would have a “rise” of 3, or go up 3; and a “run” of 4, or go right 4. Both numbers in the slope are either negative or positive in order to have a positive slope.



If a line goes down from left to right, then the slope has to be negative. For example, a slope of $-\frac{3}{4}$ would have a “rise” of -3, or go down 3; and a “run” of 4, or go right 4. Only one number in the slope can be negative for a line to have a negative slope.



\paragraph{Other Types of Slope}



There are two special circumstances: no slope and slope of zero. A horizontal line has a slope of 0, and a vertical line has an undefined slope.



Horizontal lines have the form: $y = a$ ; where $a$ is a constant, i.e. $a \in \mathbb{R}$.



Vertical lines have the form: $x = a$ ; where $a$ is a constant, i.e. $a \in \mathbb{R}$.



\paragraph{Determining Slope}



To determine the slope, you need some information. This can include two (or more) coordinates, a parallel slope and a coordinate, a perpendicular slope and a coordinate, or the y-intercept and slope.



For completely horizontal lines, the difference in $y$ coordinates between any two points is 0, so the slope $m = 0$, indicating no steepness in the line at all. If the line extends between right-upper $(+,+)$ and left-lower $(-,-)$ directions, then the slope is positive. As the slope increases, the line becomes steeper until the line is almost vertical when the slope is very large. When the slope $m = 1$, the line is diagonal with an angle halfway between the $x$ and $y$ axes. If the line extends between left-upper $(-,+)$ and right-lower $(+,-)$ directions, then the slope is negative. As the slope changes from 0 to very negative numbers, the steepness in the opposite direction increases. Compare the slope $(m)$ values in the following graph of functions $y = 1$ (where $m = 0$), $y = \frac{1}{2}x + 1$, $y = x + 1$, $y = 2x$, $y = -\frac{1}{2}x + 1$, $y = -x + 1$, and $y = -2x + 1$. For all two-variable linear equations that can be converted to linear functions, the same calculation applies to slopes for those lines.



\textbf{Various Linear Slopes}



% [Image: slopes]



For the most part, finding the slope when given information is a simple matter. Simply take the slope equation $y = mx + b$ and replace the variable with whatever information you know, and solve.



\paragraph{Two Coordinates}



To find the slope with two coordinates, you must first find the slope. Use the standard equation $\frac{y\_2 - y\_1}{x\_2 - x\_1}$. Put that into the equation as $m$, and replace $x$ and $y$ with $x$ and $y$ from one of the coordinates. Solve for $b$. Put that into the equation, and you''re done.



Example: $(1, 4)$ $(4, 8)$



$$

m = \frac{y\_2 - y\_1}{x\_2 - x\_1} = \frac{4 - 1}{4 - 8} = \frac{3}{-4}

$$



Plug that right in.



$$

y = \frac{3}{-4}x + b

$$



$$

4 = \frac{3}{-4}(1) + b

$$



$$

4 - \frac{3}{-4} = b

$$



$$

\frac{-19}{4} = b

$$



Put that in the equation, and you''re done.



$$

y = \frac{3}{-4}x + \frac{-19}{4}

$$





\paragraph{Parallel Lines}



If you have to find the slope of a line (Let''s say $AB$) which is parallel to line (Let''s say $XY$), then using the coordinates of line $XY$, you can find the coordinates of the slope of line $AB$. As, the slope of a line parallel to another line is equal, i.e. Slope of $AB$ = Slope of $XY$.



\textbf{Example:}



$AB$ and $XY$ are PARALLEL Lines. Find the slope of $AB$.



Let line $XY$ have the coordinates:



$X(2, 4) = (x\_1, y\_1)$ and $Y(3, 6) = (x\_2, y\_2)$



\textbf{Slope of $XY$:}



$$

m = \frac{y\_2 - y\_1}{x\_2 - x\_1} = \frac{6 - 4}{3 - 2} = \frac{2}{1} \Rightarrow m = 2

$$



Slope of $AB = $ Slope of $XY$, so Slope of $AB = 2$



---



\paragraph{Resources:}



\begin{itemize}
    \item \href{https://courses.lumenlearning.com/beginalgebra/chapter/introduction-to-solving-linear-equations/}{Solving Linear Equations}, Lumen Learning
    \item \href{https://openstax.org/books/intermediate-algebra-2e/pages/2-1-use-a-general-strategy-to-solve-linear-equations}{General Strategy to Solving Linear Equations}, OpenStax
    \item \href{https://openstax.org/books/intermediate-algebra-2e/pages/2-3-solve-a-formula-for-a-specific-variable}{Solving Linear Equations for a Specific Variable}, OpenStax
    \item \href{https://www.youtube.com/watch?v=bsGyJ0GeRy8}{Solving Linear Equations: Example 1.} Produced by TA Catherine Sporer, UTSA
    \item \href{https://www.youtube.com/watch?v=nGZcfHQEdqo}{Solving Linear Equations: Example 2.} Produced by TA Catherine Sporer, UTSA
\end{itemize}

\paragraph{Licensing}



Content obtained and/or adapted from:



\begin{itemize}
    \item \href{https://en.wikibooks.org/wiki/Algebra/Linear_Equations_and_Functions}{Linear Equations and Functions, Wikibooks: Algebra} under a CC BY-SA license
    \item \href{https://en.wikibooks.org/wiki/Algebra/Intercepts}{Intercepts, Wikibooks: Algebra} under a CC BY-SA license
    \item \href{https://en.wikibooks.org/wiki/Algebra/Slope}{Slope, Wikibooks: Algebra} under a CC BY-SA license
\end{itemize}

\subsection{Learning Objectives}

\begin{enumerate}
    \item \textbf{1. Understanding and Application of Linear Equations:}
    \item Students will be able to define and explain linear equations, including their components such as slope and y-intercept.
    \item Students will apply their knowledge to solve linear equations and interpret the results in real-world contexts. \textbf{2. Graphing Linear Equations:}
    \item Students will be able to graph linear equations on the Cartesian plane, identifying key features such as the slope, y-intercept, and x-intercept.
    \item Students will interpret and analyze graphs to understand the relationship between the variables. \textbf{3. Manipulating and Transforming Linear Equations:}
    \item Students will be able to perform algebraic manipulations on linear equations, including addition, subtraction, multiplication, and division of both sides of the equation.
    \item Students will demonstrate the ability to transform equations to isolate variables and solve for specific values.
\end{enumerate}

\subsection{Assessment Instrument}

\begin{enumerate}
    \item \textbf{Part 1: Multiple Choice (20 points)}
    \item Which of the following represents the slope-intercept form of a linear equation?
    \item a) $y = mx + b$
    \item b) $y = ax^2 + bx + c$
    \item c) $x = my + b$
    \item d) $y = m(x - b)$
    \item If the equation of a line is $y = 2x + 3$, what is the slope?
    \item a) $2$
    \item b) $3$
    \item c) $-2$
    \item d) $-3$ \textbf{Part 2: Short Answer (20 points)}
    \item Explain the significance of the $y$-intercept in the equation $y = mx + b$.
    \item How can you determine if two lines are parallel by looking at their equations? \textbf{Part 3: Graphing (30 points)}
    \item Graph the following equations on the Cartesian plane:
    \item $y = 2x + 1$
    \item $y = -x + 4$
    \item $x = 3$ \textbf{Part 4: Problem Solving (30 points)}
    \item Solve for $y$ in the equation $3y - 2x = 6$.
    \item Given the points $(2, 3)$ and $(4, 7)$, find the equation of the line passing through these points in slope-intercept form.
    \item A line passes through the point $(1, -2)$ and has a slope of $3$. Write the equation of the line. \textbf{Scoring Rubric:}
    \item Multiple Choice: 1 point for each correct answer.
    \item Short Answer: 5 points each. Clarity, accuracy, and completeness are considered.
    \item Graphing: 10 points for each graph. Correctness, labeling, and neatness are considered.
    \item Problem Solving: Points awarded based on correctness and method used for solving the problem. \textbf{Total Points: 100}
\end{enumerate}

%%===============================================================
\section{Rational Equations}
\label{sec:rational-equations}
%%===============================================================

\paragraph{Rational Equations}

Rational equations are equations containing rational expressions (or expressions with fractions that contain real numbers and/or variables). Some examples of rational equations:

1. $ \frac{1}{8}x + 2 = \frac{1}{4}x $
2. $ \frac{1}{x+1} = \frac{2}{3} $
3. $ \frac{y^{2}}{3y + 7} = \frac{y}{6} $

\paragraph{Steps to Solving Rational Equations:}

1. Note any value of the variable that would make any denominator zero.
2. Find the least common denominator (LCD) of all denominators in the equation.
3. Clear the fractions by multiplying both sides of the equation by the LCD.
4. Solve the resulting equation.
5. Check: If any values found in step 1 are algebraic solutions, discard them. Check any remaining solutions in the original equation.

\paragraph{Example Problem:}

$$
1 - \frac{1}{x} = \frac{2}{x^2}
$$

1. If $x = 0$, the denominator of $\frac{1}{x}$ and $\frac{2}{x^2}$ will be zero.
2. The least common denominator of all terms in the equation is $x^2$.
3. Multiplying each side of the equation $1 - \frac{1}{x} = \frac{2}{x^2}$ by $x^2$ gives us:

$$
x^2 - x = 2
$$

$$
x^2 - x = 2 \rightarrow x^2 - x - 2 = 0 \rightarrow (x - 2)(x + 1) = 0 \rightarrow x = -1, x = 2
$$

4. None of these solutions were noted in step 1, so we can check our two solutions:

\begin{itemize}
    \item For $x = -1$:
\end{itemize}

$$
1 - \frac{1}{-1} = \frac{2}{(-1)^2} \rightarrow 1 - (-1) = 2 \rightarrow 1 + 1 = 2
$$

\begin{itemize}
    \item For $x = 2$:
\end{itemize}

$$
1 - \frac{1}{2} = \frac{2}{2^2} \rightarrow 1 - \frac{1}{2} = \frac{2}{4} \rightarrow \frac{1}{2} = \frac{2}{4}
$$

Thus, $x = -1$ and $x = 2$ are solutions to our original rational equation.

\paragraph{Resources}

\begin{itemize}
    \item \href{https://openstax.org/books/intermediate-algebra-2e/pages/7-4-solve-rational-equations}{Solve Rational Equations, OpenStax}
    \item \href{https://www.youtube.com/watch?v=1fR_9ke5-n8}{Solving Rational Equations (Example), The Organic Chemistry Tutor}
    \item \href{https://www.youtube.com/watch?v=MFjx52NW5Bw}{Solving Rational Equations with Different Denominators (Example), The Organic Chemistry Tutor}
\end{itemize}

\subsection{Learning Objectives}

\begin{enumerate}
    \item \textbf{Identify and Address Undefined Values:} Students will identify values that make any denominator in a rational equation zero and understand why these values are excluded from solutions.
    \item \textbf{Solve Rational Equations:} Students will solve rational equations by finding the least common denominator, clearing fractions, solving the resulting equation, and checking solutions for validity.
    \item \textbf{Verify Solutions:} Students will check solutions to rational equations to ensure they do not make any denominator zero and confirm that they satisfy the original equation.
\end{enumerate}

\subsection{Assessment Instrument}

\begin{enumerate}
    \item \textbf{Instructions:}
    \item For each problem, show all steps involved in solving the equation.
    \item Clearly indicate any values that make denominators zero and exclude them from your solutions.
    \item Verify each solution by substituting it back into the original equation.
    \item \textbf{Problem 1:} Solve the following rational equation and check your solution: $$ \frac{1}{x-2} + \frac{3}{x} = \frac{4}{x(x-2)} $$
    \item \textbf{Problem 2:} Determine if the following value is a solution to the rational equation: $$ \frac{2x}{x+1} = \frac{3}{x-1} $$ where \( x = -1 \).
    \item \textbf{Problem 3:} Solve the rational equation: $$ \frac{2}{x^2-1} = \frac{3}{x+1} $$ and verify that your solutions do not make any denominator zero.
\end{enumerate}

%%===============================================================
\section{Models and Applications}
\label{sec:models-and-applications}
%%===============================================================

\paragraph{Models and Applications}

Many real-world applications can be modeled by linear equations. For example, a cell phone package may include a monthly service fee plus a charge per minute of talk-time; it costs a widget manufacturer a certain amount to produce $x$ widgets per month plus monthly operating charges; a car rental company charges a daily fee plus an amount per mile driven. These are examples of applications we come across every day that are modeled by linear equations. In this section, we will set up and use linear equations to solve such problems.

\paragraph{Writing a Linear Equation to Solve an Application}

To set up or model a linear equation to fit a real-world application, we must first determine the known quantities and define the unknown quantity as a variable. Then, we begin to interpret the words as mathematical expressions using mathematical symbols. Let us use the car rental example above. In this case, a known cost, such as $0.10/\text{mi}$, is multiplied by an unknown quantity, the number of miles driven. Therefore, we can write $0.10x$. This expression represents a variable cost because it changes according to the number of miles driven.

If a quantity is independent of a variable, we usually just add or subtract it according to the problem. As these amounts do not change, we call them fixed costs. Consider a car rental agency that charges $0.10/mi.$ plus a daily fee of ```$```50. We can use these quantities to model an equation that can be used to find the daily car rental cost $C$.

$$
C = 0.10x + 50
$$

When dealing with real-world applications, there are certain expressions that we can translate directly into math. The list below features some common verbal expressions and their equivalent mathematical expressions.

One number exceeds another by $a$ is expressed as $x, x + a$ 
                    
Twice a number is expressed as $2x$

One number is $a$ more than another number is expressed as $x, x + a$
                     
One number is $a$ less than twice another number is expressed as $x, 2x - a$ 
 
The product of a number and $a$, decreased by $b$ is expressed as $ax - b$  
                     
The quotient of a number and the number plus $a$ is three times the number is expressed as $\frac{x}{x+a} = 3x$

The product of three times a number and the number decreased by $b$ is $c$ is expressed as  $3x(x-b) = c$ 

\paragraph{How To: Given a real-world problem, model a linear equation to fit it}

1. Identify known quantities.
2. Assign a variable to represent the unknown quantity.
3. If there is more than one unknown quantity, find a way to write the second unknown in terms of the first.
4. Write an equation interpreting the words as mathematical operations.
5. Solve the equation. Be sure the solution can be explained in words including the units of measure.

\paragraph{Example: Modeling a Linear Equation to Solve an Unknown Number Problem}

Find a linear equation to solve for the following unknown quantities: One number exceeds another number by 17 and their sum is 31. Find the two numbers.

Let $x$ equal the first number. Then, as the second number exceeds the first by 17, we can write the second number as $x + 17$. The sum of the two numbers is 31. We usually interpret the word "is" as an equal sign.

$$
x + (x + 17) = 31
$$
$$
2x + 17 = 31
$$
Simplify and solve.
$$
2x = 14
$$
$$
x = 7
$$
$$
x + 17 = 7 + 17 = 24
$$

The two numbers are 7 and 24.

\paragraph{Example: Setting Up a Linear Equation to Solve a Real-World Application}

There are two cell phone companies that offer different packages. Company A charges a monthly service fee of ```$```34 plus ```$```0.05/min$ talk-time. Company B charges a monthly service fee of ```$```40 plus ```$```0.04/min talk-time.

\begin{itemize}
    \item Write a linear equation that models the packages offered by both companies.
    \item If the average number of minutes used each month is 1,160, which company offers the better plan?
    \item If the average number of minutes used each month is 420, which company offers the better plan?
    \item How many minutes of talk-time would yield equal monthly statements from both companies?
\end{itemize}

The model for Company A can be written as $A = 0.05x + 34$. This includes the variable cost of $0.05x$ plus the monthly service charge of ```$```34. Company B’s package charges a higher monthly fee of ```$```40, but a lower variable cost of $0.04x$. Company B’s model can be written as $B = 0.04x + 40$.

If the average number of minutes used each month is 1,160, we have the following:
$$
\text{Company A} = 0.05(1,160) + 34 = 58 + 34 = 92
$$
$$
\text{Company B} = 0.04(1,160) + 40 = 46.4 + 40 = 86.4
$$

So, Company B offers the lower monthly cost of ```$```86.40 as compared with the ```$```92 monthly cost offered by Company A when the average number of minutes used each month is 1,160.

If the average number of minutes used each month is 420, we have the following:
$$
\text{Company A} = 0.05(420) + 34 = 21 + 34 = 55
$$
$$
\text{Company B} = 0.04(420) + 40 = 16.8 + 40 = 56.8
$$

If the average number of minutes used each month is 420, then Company A offers a lower monthly cost of ```$```55 compared to Company B’s monthly cost of ```$```56.80.

To answer the question of how many talk-time minutes would yield the same bill from both companies, we should think about the problem in terms of $(x, y)$ coordinates: At what point are both the $x$-value and the $y$-value equal? We can find this point by setting the equations equal to each other and solving for $x$.
$$
0.05x + 34 = 0.04x + 40
$$
$$
0.01x = 6
$$
$$
x = 600
$$

Check the $x$-value in each equation.
$$
0.05(600) + 34 = 64
$$
$$
0.04(600) + 40 = 64
$$

Therefore, a monthly average of 600 talk-time minutes renders the plans equal.

\paragraph{Example: Solving an Area Problem}

The perimeter of a tablet of graph paper is 48 in². The length is 6 in. more than the width. Find the area of the graph paper.

The standard formula for area is $A = LW$; however, we will solve the problem using the perimeter formula. The reason we use the perimeter formula is because we know enough information about the perimeter that the formula will allow us to solve for one of the unknowns. As both perimeter and area use length and width as dimensions, they are often used together to solve a problem such as this one.

We know that the length is 6 in. more than the width, so we can write length as $L = W + 6$. Substitute the value of the perimeter and the expression for length into the perimeter formula and find the length.
$$
P = 2L + 2W
$$
$$
48 = 2(W + 6) + 2W
$$
$$
48 = 2W + 12 + 2W
$$
$$
48 = 4W + 12
$$
$$
36 = 4W
$$
$$
9 = W
$$
$$
(9 + 6) = L
$$
$$
15 = L
$$

Now, we find the area given the dimensions of $L = 15$ in. and $W = 9$ in.
$$
A = LW
$$
$$
A = 15(9) = 135 \text{ in}^2
$$

The area is 135 in².

\paragraph{Licensing}

Content obtained and/or adapted from:

\begin{itemize}
    \item \href{https://courses.lumenlearning.com/wmopen-collegealgebra/chapter/introduction-models-and-applications/}{Models and Applications, Lumen Learning College Algebra} under a CC BY license
\end{itemize}

\subsection{Learning Objectives}

\begin{enumerate}
    \item Student Learning Objectives (SLOs) 1.\textbf{Understand and Model Linear Equations}
    \item \textbf{Objective}: Students will be able to identify known and unknown quantities in a real-world problem and model it using a linear equation.
    \item \textbf{Performance Indicator}: Given a problem scenario, students can accurately set up a linear equation representing the situation. 2.\textbf{Solve Real-World Linear Equation Problems}
    \item \textbf{Objective}: Students will solve linear equations derived from real-world problems.
    \item \textbf{Performance Indicator}: Students will correctly solve linear equations and interpret the results in the context of the original problem. 3.\textbf{Translate Verbal Descriptions to Mathematical Expressions}
    \item \textbf{Objective}: Students will translate common verbal expressions into equivalent mathematical expressions and operations.
    \item \textbf{Performance Indicator}: Students can convert phrases like "twice a number" or "one number exceeds another by a" into appropriate mathematical expressions.
\end{enumerate}

\subsection{Assessment Instrument}

\begin{enumerate}
    \item Problem 1: Modeling and Solving a Real-World Linear Equation \textbf{Question}: A car rental company charges a daily fee of ```$```50 plus ```$```0.10 per mile driven. Write a linear equation representing the total cost $C$ for renting a car for one day and driving $x$ miles. Then, find the total cost if a customer drives 150 miles in one day. Problem 2: Solving a Real-World Application \textbf{Question}: Company A charges a monthly service fee of ```$```34 plus ```$```0.05 per minute of talk-time. Company B charges a monthly service fee of ```$```40 plus ```$```0.04 per minute of talk-time. Write linear equations to model the cost $A$ for Company A and the cost $B$ for Company B. Determine the number of minutes for which both companies charge the same amount. Problem 3: Translating Verbal Descriptions to Mathematical Expressions \textbf{Question}: Translate the following verbal expressions into mathematical expressions:
    \item Twice a number decreased by 5.
    \item The quotient of a number and the number plus 4 equals three times the number.
    \item One number is 6 more than another number, and their sum is 20.
\end{enumerate}

%%===============================================================
\section{Complex Numbers}
\label{sec:complex-numbers}
%%===============================================================

\paragraph{Complex Numbers}

In mathematics, a complex number is a number that can be expressed in the form $a + bi$, where $a$ and $b$ are real numbers, and $i$ is a symbol, called the imaginary unit, that satisfies the equation $i^2 = -1$. Because no real number satisfies this equation, $i$ was called an imaginary number by René Descartes. For the complex number $a + bi$, $a$ is called the real part and $b$ is called the imaginary part. The set of complex numbers is denoted by either of the symbols $\mathbb{C}$ or $C$. Despite the historical nomenclature "imaginary", complex numbers are regarded in the mathematical sciences as just as "real" as the real numbers and are fundamental in many aspects of the scientific description of the natural world.

Complex numbers allow solutions to all polynomial equations, even those that have no solutions in real numbers. More precisely, the fundamental theorem of algebra asserts that every polynomial equation with real or complex coefficients has a solution which is a complex number. For example, the equation $(x + 1)^2 = -9$ has no real solution, since the square of a real number cannot be negative, but has the two nonreal complex solutions $-1 + 3i$ and $-1 - 3i$.

A complex number can be visually represented as a pair of numbers $(a,?b)$ forming a vector on a diagram called an Argand diagram, representing the complex plane. $\text{Re}$ is the real axis, $\text{Im}$ is the imaginary axis, and $i$ is the "imaginary unit" that satisfies $i^2 = -1$.
% [Image: graph]

Addition, subtraction, and multiplication of complex numbers can be naturally defined by using the rule $i^2 = -1$ combined with the associative, commutative, and distributive laws. Every nonzero complex number has a multiplicative inverse. This makes the complex numbers a field that has the real numbers as a subfield. The complex numbers form also a real vector space of dimension two, with $\{1, i\}$ as a standard basis.

This standard basis makes the complex numbers a Cartesian plane, called the complex plane. This allows a geometric interpretation of the complex numbers and their operations, and conversely expressing in terms of complex numbers some geometric properties and constructions. For example, the real numbers form the real line which is identified to the horizontal axis of the complex plane. The complex numbers of absolute value one form the unit circle. The addition of a complex number is a translation in the complex plane, and the multiplication by a complex number is a similarity centered at the origin. The complex conjugation is the reflection symmetry with respect to the real axis. The complex absolute value is a Euclidean norm.

In summary, the complex numbers form a rich structure that is simultaneously an algebraically closed field, a commutative algebra over the reals, and a Euclidean vector space of dimension two.

\paragraph{Definition}

An illustration of the complex number $z = x + iy$ on the complex plane. 
% [Image: graph]

The real part is $x$, and its imaginary part is $y$. A complex number is a number of the form $a + bi$, where $a$ and $b$ are real numbers, and $i$ is an indeterminate satisfying $i^2 = -1$. For example, $2 + 3i$ is a complex number.

This way, a complex number is defined as a polynomial with real coefficients in the single indeterminate $i$, for which the relation $i^2 = -1$ is imposed. Based on this definition, complex numbers can be added and multiplied, using the addition and multiplication for polynomials. The relation $i^2 = -1$ induces the equalities $i^{2k} = (-1)^k$ and $i^{2k+1} = (-1)^ki$ which hold for all integers $k$; these allow the reduction of any polynomial that results from the addition and multiplication of complex numbers to a linear polynomial in $i$, again of the form $a + bi$ with real coefficients $a$, $b$.

The real number $a$ is called the real part of the complex number $a + bi$; the real number $b$ is called its imaginary part. To emphasize, the imaginary part does not include a factor $i$; that is, the imaginary part is $b$, not $bi$.

Formally, the complex numbers are defined as the quotient ring of the polynomial ring in the indeterminate $i$, by the ideal generated by the polynomial $i^2 + 1$.

\paragraph{Notation}

A real number $a$ can be regarded as a complex number $a + 0i$, whose imaginary part is $0$. A purely imaginary number $bi$ is a complex number $0 + bi$, whose real part is zero. As with polynomials, it is common to write $a$ for $a + 0i$ and $bi$ for $0 + bi$. Moreover, when the imaginary part is negative, that is, $b = -|b| < 0$, it is common to write $a - |b|i$ instead of $a + (-|b|)i$; for example, for $b = -4$, $3 - 4i$ can be written instead of $3 + (-4)i$.

Since the multiplication of the indeterminate $i$ and a real is commutative in polynomials with real coefficients, the polynomial $a + bi$ may be written as $a + ib$. This is often expedient for imaginary parts denoted by expressions, for example, when $b$ is a radical.

The real part of a complex number $z$ is denoted by $\text{Re}(z)$, $\mathcal{Re}(z)$, or $\mathfrak{R}(z)$; the imaginary part of a complex number $z$ is denoted by $\text{Im}(z)$, $\mathcal{Im}(z)$, or $\mathfrak{I}(z)$. For example,
$$
\text{Re}(2 + 3i) = 2 \quad \text{and} \quad \text{Im}(2 + 3i) = 3.
$$
The set of all complex numbers is denoted by $\mathbb{C}$ (blackboard bold) or $C$ (upright bold).

In some disciplines, particularly in electromagnetism and electrical engineering, $j$ is used instead of $i$ as $i$ is frequently used to represent electric current.

\paragraph{Visualization}

A complex number $z$, as a point (black) and its position vector (blue)
% [Image: graph]

A complex number $z$ can thus be identified with an ordered pair $\left( \Re(z), \Im(z) \right)$ of real numbers, which in turn may be interpreted as coordinates of a point in a two-dimensional space. The most immediate space is the Euclidean plane with suitable coordinates, which is then called the complex plane or Argand diagram, named after Jean-Robert Argand. Another prominent space on which the coordinates may be projected is the two-dimensional surface of a sphere, which is then called the Riemann sphere.

\paragraph{Cartesian Complex Plane}

The definition of the complex numbers involving two arbitrary real values immediately suggests the use of Cartesian coordinates in the complex plane. The horizontal (real) axis is generally used to display the real part, with increasing values to the right, and the imaginary part marks the vertical (imaginary) axis, with increasing values upwards.

A charted number may be viewed either as the coordinatized point or as a position vector from the origin to this point. The coordinate values of a complex number $z$ can hence be expressed in its Cartesian, rectangular, or algebraic form.

Notably, the operations of addition and multiplication take on a very natural geometric character, when complex numbers are viewed as position vectors: addition corresponds to vector addition, while multiplication (see below) corresponds to multiplying their magnitudes and adding the angles they make with the real axis. Viewed in this way, the multiplication of a complex number by $i$ corresponds to rotating the position vector counterclockwise by a quarter turn (90°) about the origin—a fact which can be expressed algebraically as follows:
$$
(a + bi) \cdot i = ai + b(i^2) = -b + ai.
$$

Argument $\varphi$ and modulus $r$ locate a point in the complex plane.
% [Image: graph]

\paragraph{Polar Complex Plane}

\paragraph{Modulus and Argument}

An alternative option for coordinates in the complex plane is the polar coordinate system that uses the distance of the point $z$ from the origin (O), and the angle subtended between the positive real axis and the line segment $Oz$ in a counterclockwise sense. This leads to the polar form of complex numbers.

The absolute value (or modulus or magnitude) of a complex number $z = x + yi$ is
$$
r = |z| = \sqrt{x^2 + y^2}.
$$
If $z$ is a real number (that is, if $y = 0$), then $r = |x|$. That is, the absolute value of a real number equals its absolute value as a complex number. By Pythagoras'' theorem, the absolute value of a complex number is the distance to the origin of the point representing the complex number in the complex plane.

The argument of $z$ (in many applications referred to as the "phase" $\varphi$) is the angle of the radius $Oz$ with the positive real axis, and is written as $\arg z$. As with the modulus, the argument can be found from the rectangular form $x + yi$ by applying the inverse tangent to the quotient of imaginary-by-real parts. By using a half-angle identity, a single branch of the arctan suffices to cover the range of the $\arg$-function, $(-\pi, \pi]$, and avoids a more subtle case-by-case analysis
$$
\varphi = \arg(x + yi) = 
\begin{cases} 
2\arctan\left(\dfrac{y}{\sqrt{x^2 + y^2} + x}\right) & \text{if } x > 0 \text{ or } y \neq 0, \\pi & \text{if } x < 0 \text{ and } y = 0, \\text{undefined} & \text{if } x = 0 \text{ and } y = 0. 
\end{cases}
$$
Normally, as given above, the principal value in the interval $(-\pi, \pi]$ is chosen. If the $\arg$ value is negative, values in the range $(-\pi, \pi]$ or $[0, 2\pi)$ can be obtained by adding $2\pi$. The value of $\varphi$ is expressed in radians in this article. It can increase by any integer multiple of $2\pi$ and still give the same angle, viewed as subtended by the rays of the positive real axis and from the origin through $z$. Hence, the $\arg$ function is sometimes considered as multivalued. The polar angle for the complex number $0$ is indeterminate, but arbitrary choice of the polar angle $0$ is common.

The value of $\varphi$ equals the result of $\text{atan2}$:
$$
\varphi = \operatorname{atan2} \left(\operatorname{Im}(z), \operatorname{Re}(z)\right).
$$
Together, $r$ and $\varphi$ give another way of representing complex numbers, the polar form, as the combination of modulus and argument fully specify the position of a point on the plane. Recovering the original rectangular coordinates from the polar form is done by the formula called trigonometric form
$$
z = r(\cos \varphi + i\sin \varphi).
$$
Using Euler''s formula, this can be written as
$$
z = re^{i\varphi} \text{ or } z = r\exp i\varphi.
$$
Using the cis function, this is sometimes abbreviated to
$$
z = r \operatorname{cis} \varphi.
$$
In angle notation, often used in electronics to represent a phasor with amplitude $r$ and phase $\varphi$, it is written as
$$
z = r \angle \varphi.
$$

\paragraph{Complex Graphs}

A color wheel graph of the expression $\frac{(z^2 - 1)(z - 2 - i)^2}{z^2 + 2 + 2i}$
% [Image: graph]

When visualizing complex functions, both a complex input and output are needed. Because each complex number is represented in two dimensions, visually graphing a complex function would require the perception of a four-dimensional space, which is possible only in projections. Because of this, other ways of visualizing complex functions have been designed.

In domain coloring, the output dimensions are represented by color and brightness, respectively. Each point in the complex plane as domain is ornated, typically with color representing the argument of the complex number, and brightness representing the magnitude. Dark spots mark moduli near zero, brighter spots are farther away from the origin, the gradation may be discontinuous, but is assumed as monotonous. The colors often vary in steps of $\frac{\pi}{3}$ for $0$ to $2\pi$ from red, yellow, green, cyan, blue, to magenta. These plots are called color wheel graphs. This provides a simple way to visualize the functions without losing information. The picture shows zeros for $\pm1$, $(2 + i)$ and poles at $\sqrt{-2 - 2i}$.

Riemann surfaces are another way to visualize complex functions. Riemann surfaces can be thought of as deformations of the complex plane; while the horizontal axes represent the real and imaginary inputs, the single vertical axis only represents either the real or imaginary output. However, Riemann surfaces are built in such a way that rotating them 180 degrees shows the imaginary output, and vice versa. Unlike domain coloring, Riemann surfaces can represent multivalued functions like $\sqrt{z}$.

\paragraph{Relations and Operations}

\paragraph{Equality}

Complex numbers have a similar definition of equality to real numbers: two complex numbers $a\_1 + b\_1i$ and $a\_2 + b\_2i$ are equal if and only if both their real and imaginary parts are equal, that is, if $a\_1 = a\_2$ and $b\_1 = b\_2$. Nonzero complex numbers written in polar form are equal if and only if they have the same magnitude and their arguments differ by an integer multiple of $2\pi$.

\paragraph{Ordering}

Unlike the real numbers, there is no natural ordering of the complex numbers. In particular, there is no linear ordering on the complex numbers that is compatible with addition and multiplication – the complex numbers cannot have the structure of an ordered field. This is because every non-trivial sum of squares in an ordered field is $\neq 0$, and $i^2 + 1^2 = 0$ is a non-trivial sum of squares. Thus, complex numbers are naturally thought of as existing on a two-dimensional plane.

\paragraph{Conjugate}

Geometric representation of $z$ and its conjugate $\overline{z}$ in the complex plane
% [Image: graph]

The complex conjugate of the complex number $z = x + yi$ is given by $x - yi$. It is denoted by either $\overline{z}$ or $z^*$. This unary operation on complex numbers cannot be expressed by applying only their basic operations addition, subtraction, multiplication, and division.

Geometrically, $\overline{z}$ is the "reflection" of $z$ about the real axis. Conjugating twice gives the original complex number:

$$\overline{\overline{z}} = z,$$

which makes this operation an involution. The reflection leaves both the real part and the magnitude of $z$ unchanged, that is:

$$\operatorname{Re}(\overline{z}) = \operatorname{Re}(z)$$

and

$$|\overline{z}| = |z|.$$

The imaginary part and the argument of a complex number $z$ change their sign under conjugation:

$$\operatorname{Im}(\overline{z}) = -\operatorname{Im}(z) \quad \text{and} \quad \operatorname{arg}(\overline{z}) \equiv -\operatorname{arg}(z) \pmod{2\pi}.$$

For details on argument and magnitude, see the section on Polar form.

The product of a complex number $z$ and its conjugate is known as the absolute square. It is always a non-negative real number and equals the square of the magnitude of each:

$$z \cdot \overline{z} = x^2 + y^2 = |z|^2 = |\overline{z}|^2.$$

This property can be used to convert a fraction with a complex denominator to an equivalent fraction with a real denominator by expanding both numerator and denominator of the fraction by the conjugate of the given denominator. This process is sometimes called "rationalization" of the denominator (although the denominator in the final expression might be an irrational real number), because it resembles the method to remove roots from simple expressions in a denominator.

The real and imaginary parts of a complex number $z$ can be extracted using the conjugation:

$$\operatorname{Re}(z) = \frac{z + \overline{z}}{2},$$

and

$$\operatorname{Im}(z) = \frac{z - \overline{z}}{2i}.$$

Moreover, a complex number is real if and only if it equals its own conjugate.

Conjugation distributes over the basic complex arithmetic operations:

$$\overline{z \pm w} = \overline{z} \pm \overline{w},$$

$$\overline{z \cdot w} = \overline{z} \cdot \overline{w},$$

$$\overline{\frac{z}{w}} = \frac{\overline{z}}{\overline{w}}.$$

Conjugation is also employed in inversive geometry, a branch of geometry studying reflections more general than ones about a line. In the network analysis of electrical circuits, the complex conjugate is used in finding the equivalent impedance when the maximum power transfer theorem is considered.

\paragraph{Addition and Subtraction}

Addition of two complex numbers can be done geometrically by constructing a parallelogram. 

% [Image: graph]

Two complex numbers $a$ and $b$ are most easily added by separately adding their real and imaginary parts of the summands. That is to say:

$$a + b = (x + yi) + (u + vi) = (x + u) + (y + v)i.$$

Similarly, subtraction can be performed as:

$$a - b = (x + yi) - (u + vi) = (x - u) + (y - v)i.$$

Using the visualization of complex numbers in the complex plane, addition has the following geometric interpretation: the sum of two complex numbers $a$ and $b$, interpreted as points in the complex plane, is the point obtained by building a parallelogram from the three vertices $O$, and the points of the arrows labeled $a$ and $b$ (provided that they are not on a line). Equivalently, calling these points $A$ and $B$, respectively, and the fourth point of the parallelogram $X$, the triangles $OAB$ and $XBA$ are congruent. A visualization of the subtraction can be achieved by considering the addition of the negative subtrahend.

\paragraph{Multiplication and Square}

The rules of the distributive property, the commutative properties (of addition and multiplication), and the defining property $i^2 = -1$ apply to complex numbers. It follows that:

$$(x + yi)(u + vi) = (xu - yv) + (xv + yu)i.$$

In particular:

$$(x + yi)^2 = x^2 - y^2 + 2xyi.$$

\paragraph{Reciprocal and Division}

Using the conjugation, the reciprocal of a nonzero complex number $z = x + yi$ can always be broken down to:

$$\frac{1}{z} = \frac{\overline{z}}{z \overline{z}} = \frac{\overline{z}}{|z|^2} = \frac{\overline{z}}{x^2 + y^2} = \frac{x}{x^2 + y^2} - \frac{y}{x^2 + y^2}i,$$

since non-zero implies that $x^2 + y^2$ is greater than zero.

This can be used to express a division of an arbitrary complex number $w$ by a non-zero complex number $z$ as:

$$\frac{w}{z} = w \cdot \frac{1}{z} = (u + vi) \cdot \left(\frac{x}{x^2 + y^2} - \frac{y}{x^2 + y^2}i\right) = \frac{(ux + vy) + (vx - uy)i}{x^2 + y^2}.$$

\paragraph{Multiplication and Division in Polar Form}

Formulas for multiplication, division, and exponentiation are simpler in polar form than the corresponding formulas in Cartesian coordinates. Given two complex numbers $z\_1 = r\_1(\cos \varphi\_1 + i \sin \varphi\_1)$ and $z\_2 = r\_2(\cos \varphi\_2 + i \sin \varphi\_2)$, because of the trigonometric identities

$$
\begin{aligned}
\cos a \cos b - \sin a \sin b &= \cos(a + b) \\cos a \sin b + \sin a \cos b &= \sin(a + b),
\end{aligned}
$$

we may derive

$$
z\_1 z\_2 = r\_1 r\_2 (\cos (\varphi\_1 + \varphi\_2) + i \sin (\varphi\_1 + \varphi\_2)).
$$

In other words, the absolute values are multiplied and the arguments are added to yield the polar form of the product. For example, multiplying by $i$ corresponds to a quarter-turn counter-clockwise, which gives back $i^2 = -1$. The picture below illustrates the multiplication of

$$
(2 + i)(3 + i) = 5 + 5i.
$$

% [Image: graph]

Multiplication of $2 + i$ (blue triangle) and $3 + i$ (red triangle). The red triangle is rotated to match the vertex of the blue one and stretched by $\sqrt{5}$, the length of the hypotenuse of the blue triangle.

Since the real and imaginary parts of $5 + 5i$ are equal, the argument of that number is $45^\circ$, or $\frac{\pi}{4}$ (in radians). On the other hand, it is also the sum of the angles at the origin of the red and blue triangles, which are $\arctan\left(\frac{1}{3}\right)$ and $\arctan\left(\frac{1}{2}\right)$, respectively. Thus, the formula

$$
\frac{\pi}{4} = \arctan \left(\frac{1}{2}\right) + \arctan \left(\frac{1}{3}\right)
$$

holds. As the $\arctan$ function can be approximated highly efficiently, formulas like this – known as Machin-like formulas – are used for high-precision approximations of $\pi$.

Similarly, division is given by

$$
\frac{z\_1}{z\_2} = \frac{r\_1}{r\_2} \left(\cos (\varphi\_1 - \varphi\_2) + i \sin (\varphi\_1 - \varphi\_2)\right).
$$

\paragraph{Square Root}

The square roots of $a + bi$ (with $b \ne 0$) are $\pm (\gamma + \delta i)$, where

$$
\gamma = \sqrt{\frac{a + \sqrt{a^2 + b^2}}{2}}
$$

and

$$
\delta = (\operatorname{sgn} b) \sqrt{\frac{-a + \sqrt{a^2 + b^2}}{2}},
$$

where $\operatorname{sgn}$ is the signum function. This can be seen by squaring $\pm (\gamma + \delta i)$ to obtain $a + bi$.

\paragraph{Exponential Function}

The exponential function $\exp: \mathbb{C} \to \mathbb{C};\ z \mapsto \exp z$ can be defined for every complex number $z$ by the power series

$$
\exp z = \sum_{n=0}^{\infty} \frac{z^n}{n!},
$$

which has an infinite radius of convergence.

The value at 1 of the exponential function is Euler''s number:

$$
e = \exp 1 = \sum_{n=0}^{\infty} \frac{1}{n!} \approx 2.71828.
$$

If $z$ is real, one has $\exp z = e^z$. Analytic continuation allows extending this equality for every complex value of $z$, and thus to define the complex exponentiation with base $e$ as

$$
e^z = \exp z.
$$

\paragraph{Functional Equation}

The exponential function satisfies the functional equation

$$
e^{z + t} = e^z e^t.
$$

This can be proved either by comparing the power series expansion of both members or by applying analytic continuation from the restriction of the equation to real arguments.

\paragraph{Euler''s Formula}

Euler''s formula states that, for any real number $y$:

$$
e^{iy} = \cos y + i \sin y.
$$

The functional equation implies thus that, if $x$ and $y$ are real, one has

$$
e^{x + iy} = e^x (\cos y + i \sin y) = e^x \cos y + i e^x \sin y,
$$

which is the decomposition of the exponential function into its real and imaginary parts.

\paragraph{Complex Logarithm}

In the real case, the natural logarithm can be defined as the inverse $\ln: \mathbb{R}^+ \to \mathbb{R};\ x \mapsto \ln x$ of the exponential function. For extending this to the complex domain, one can start from Euler''s formula. It implies that, if a complex number $z \in \mathbb{C}^\times$ is written in polar form

$$
z = r (\cos \varphi + i \sin \varphi)
$$

with $r, \varphi \in \mathbb{R}$, then with

$$
\ln z = \ln r + i \varphi
$$

as complex logarithm one has a proper inverse:

$$
\exp (\ln z) = \exp (\ln r + i \varphi) = r \exp (i \varphi) = r (\cos \varphi + i \sin \varphi) = z.
$$

However, because cosine and sine are periodic functions, the addition of an integer multiple of $2\pi$ to $\varphi$ does not change $z$. For example, $e^{i\pi} = e^{3i\pi} = -1$, so both $i\pi$ and $3i\pi$ are possible values for the natural logarithm of $-1$.

Therefore, if the complex logarithm is not to be defined as a multivalued function

$$
\ln z = \ln r + i (\varphi + 2\pi k) \mid k \in \mathbb{Z},
$$

one has to use a branch cut and to restrict the codomain, resulting in the bijective function

$$
\ln \colon \mathbb{C}^\times \to \mathbb{R}^+ + i (-\pi, \pi].
$$

If $z \in \mathbb{C} \setminus (-\mathbb{R}_{\geq 0})$ is not a non-positive real number (a positive or a non-real number), the resulting principal value of the complex logarithm is obtained with $-\pi < \varphi < \pi$. It is an analytic function outside the negative real numbers, but it cannot be prolonged to a function that is continuous at any negative real number $z \in -\mathbb{R}^+$, where the principal value is $\ln z = \ln(-z) + i\pi$.

\paragraph{Exponentiation}

If $x > 0$ is real and $z$ is complex, the exponentiation is defined as

$$
x^z = e^{z \ln x},
$$

where $\ln$ denotes the natural logarithm. 

It seems natural to extend this formula to complex values of $x$, but there are some difficulties resulting from the fact that the complex logarithm is not really a function, but a multivalued function.

It follows that if $z$ is as above, and if $t$ is another complex number, then the exponentiation is the multivalued function

$$
z^t =  e^{t \ln r}(\cos(\varphi t + 2\pi k) + i \sin(\varphi t + 2\pi k)) \mid k \in \mathbb{Z}$$.

\paragraph{Integer and Fractional Exponents}

If, in the preceding formula, $t$ is an integer, then the sine and the cosine are independent of $k$. Thus, if the exponent $n$ is an integer, then $z^n$ is well defined, and the exponentiation formula simplifies to de Moivre''s formula:

$$
z^n = \left(r (\cos \varphi + i \sin \varphi)\right)^n = r^n (\cos n \varphi + i \sin n \varphi).
$$

The $n$th roots of a complex number $z$ are given by

$$
z^{1/n} = \sqrt[n]{r} \left(\cos \left(\frac{\varphi + 2k\pi}{n}\right) + i \sin \left(\frac{\varphi + 2k\pi}{n}\right)\right)
$$

for $0 \leq k \leq n - 1. \ (\text{Here } \sqrt[n]{r} \text{ is the usual (positive) } n\text{th root of the positive real number } r.)$ Because sine and cosine are periodic, other integer values of $k$ do not give other values.

While the $n$th root of a positive real number $r$ is chosen to be the positive real number $c$ satisfying $c^n = r$, there is no natural way of distinguishing one particular complex $n$th root of a complex number. Therefore, the $n$th root is a $n$-valued function of $z$. This implies that, contrary to the case of positive real numbers, one has

$$
(z^n)^{1/n} \ne z,
$$

since the left-hand side consists of $n$ values, and the right-hand side is a single value.

\paragraph{Licensing}

Content obtained and/or adapted from:

\begin{itemize}
    \item \href{https://en.wikipedia.org/wiki/Complex_number}{Complex number, Wikipedia} under a CC BY-SA license
\end{itemize}

\subsection{Learning Objectives}

\begin{enumerate}
    \item SLO 1: Complex Number Representation \textbf{Objective:} Students will be able to express complex numbers in both rectangular and polar forms and convert between these forms. \textbf{Assessment Criteria:}
    \item Given a complex number in rectangular form, convert it to polar form.
    \item Given a complex number in polar form, convert it to rectangular form. SLO 2: Operations with Complex Numbers \textbf{Objective:} Students will be able to perform addition, subtraction, multiplication, and division of complex numbers, including using the polar form for these operations. \textbf{Assessment Criteria:}
    \item Solve problems involving addition, subtraction, multiplication, and division of complex numbers.
    \item Apply polar form to perform multiplication and division, and verify results by converting back to rectangular form. SLO 3: Complex Number Visualization and Properties \textbf{Objective:} Students will be able to visualize complex numbers on the complex plane and describe their geometric properties, including modulus and argument. \textbf{Assessment Criteria:}
    \item Plot complex numbers on the complex plane and identify their modulus and argument.
    \item Describe the geometric interpretation of addition, multiplication, and conjugation of complex numbers
\end{enumerate}

\subsection{Assessment Instrument}

\begin{enumerate}
    \item Part 1: Conversion Between Forms
    \item Convert the complex number $3 + 4i$ to polar form.
    \item Convert the complex number $5 \operatorname{cis} \frac{\pi}{6}$ to rectangular form. Part 2: Operations with Complex Numbers
    \item Compute $(2 + 3i) + (1 - 4i)$ and express the result in rectangular form.
    \item Compute $(2 + 3i) \cdot (1 - 2i)$ and verify the result by converting to polar form and back. Part 3: Visualization and Properties
    \item Plot the complex number $2 - 3i$ on the complex plane and calculate its modulus and argument.
    \item Given the complex number $1 + i$, determine its conjugate and describe the geometric effect of conjugation on the complex plane. Rubric
    \item \textbf{Conversion Between Forms:} Correct conversion and appropriate use of polar and rectangular forms.
    \item \textbf{Operations with Complex Numbers:} Accurate calculations and correct verification of results.
    \item \textbf{Visualization and Properties:} Correct plotting, modulus, argument calculation, and understanding of geometric interpretations.
\end{enumerate}

%%===============================================================
\section{Quadratic Equations}
\label{sec:quadratic-equations}
%%===============================================================

\paragraph{Quadratic Equations}

\paragraph{Quadratic Formula}

The quadratic formula for the roots of the general quadratic equation is:

$$
x = \frac{-b \pm \sqrt{b^2 - 4ac}}{2a}
$$

\paragraph{Definition}

In algebra, a quadratic equation (from the Latin quadratus for "square") is any equation that can be rearranged in standard form as

$$
ax^2 + bx + c = 0
$$

where $x$ represents an unknown, and $a$, $b$, and $c$ represent known numbers, where $a \neq 0$. If $a = 0$, then the equation is linear, not quadratic, as there is no $ax^2$ term. The numbers $a$, $b$, and $c$ are the coefficients of the equation and may be distinguished by calling them, respectively, the quadratic coefficient, the linear coefficient and the constant or free term.

The values of $x$ that satisfy the equation are called solutions of the equation, and roots or zeros of the expression on its left-hand side. A quadratic equation has at most two solutions. If there is only one solution, one says that it is a double root. If all the coefficients are real numbers, there are either two real solutions, or a single real double root, or two complex solutions. A quadratic equation always has two roots, if complex roots are included and a double root is counted for two.

A quadratic equation can be factored into an equivalent equation

$$
ax^2 + bx + c = a(x - r)(x - s) = 0
$$

where $r$ and $s$ are the solutions for $x$. Completing the square on a quadratic equation in standard form results in the quadratic formula, which expresses the solutions in terms of $a$, $b$, and $c$. Solutions to problems that can be expressed in terms of quadratic equations were known as early as 2000 BC.

Because the quadratic equation involves only one unknown, it is called "univariate". The quadratic equation contains only powers of $x$ that are non-negative integers, and therefore it is a polynomial equation. In particular, it is a second-degree polynomial equation, since the greatest power is two.

\paragraph{Solving the Quadratic Equation}

A quadratic equation with real or complex coefficients has two solutions, called roots. These two solutions may or may not be distinct, and they may or may not be real.

\paragraph{Factoring by Inspection}

It may be possible to express a quadratic equation $ax^2 + bx + c = 0$ as a product $(px + q)(rx + s) = 0$. In some cases, it is possible, by simple inspection, to determine values of $p$, $q$, $r$, and $s$ that make the two forms equivalent to one another. If the quadratic equation is written in the second form, then the "Zero Factor Property" states that the quadratic equation is satisfied if $px + q = 0$ or $rx + s = 0$. Solving these two linear equations provides the roots of the quadratic.

For most students, factoring by inspection is the first method of solving quadratic equations to which they are exposed. If one is given a quadratic equation in the form $x^2 + bx + c = 0$, the sought factorization has the form $(x + q)(x + s)$, and one has to find two numbers $q$ and $s$ that add up to $b$ and whose product is $c$ (this is sometimes called "Vieta''''s rule" and is related to Vieta''''s formulas). As an example, $x^2 + 5x + 6$ factors as $(x + 3)(x + 2)$. The more general case where $a$ does not equal 1 can require a considerable effort in trial and error guess-and-check, assuming that it can be factored at all by inspection.

Except for special cases such as where $b = 0$ or $c = 0$, factoring by inspection only works for quadratic equations that have rational roots. This means that the great majority of quadratic equations that arise in practical applications cannot be solved by factoring by inspection.

\paragraph{Completing the Square}

The process of completing the square makes use of the algebraic identity

$$
x^2 + 2hx + h^2 = (x + h)^2,
$$

which represents a well-defined algorithm that can be used to solve any quadratic equation. Starting with a quadratic equation in standard form, $ax^2 + bx + c = 0$:

1. Divide each side by $a$, the coefficient of the squared term.
2. Subtract the constant term $c/a$ from both sides.
3. Add the square of one-half of $b/a$, the coefficient of $x$, to both sides. This "completes the square", converting the left side into a perfect square.
4. Write the left side as a square and simplify the right side if necessary.
5. Produce two linear equations by equating the square root of the left side with the positive and negative square roots of the right side.
6. Solve each of the two linear equations.

We illustrate use of this algorithm by solving $2x^2 + 4x - 4 = 0$:

1. $x^2 + 2x - 2 = 0$
2. $x^2 + 2x = 2$
3. $x^2 + 2x + 1 = 2 + 1$
4. $(x + 1)^2 = 3$
5. $x + 1 = \pm \sqrt{3}$
6. $x = -1 \pm \sqrt{3}$

The plus–minus symbol "$\pm$" indicates that both $x = -1 + \sqrt{3}$ and $x = -1 - \sqrt{3}$ are solutions of the quadratic equation.

\paragraph{Quadratic Formula and Its Derivation}

Completing the square can be used to derive a general formula for solving quadratic equations, called the quadratic formula. The mathematical proof will now be briefly summarized. It can easily be seen, by polynomial expansion, that the following equation is equivalent to the quadratic equation:

$$
\left(x + \frac{b}{2a}\right)^2 = \frac{b^2 - 4ac}{4a^2}.
$$

Taking the square root of both sides, and isolating $x$, gives:

$$
x = \frac{-b \pm \sqrt{b^2 - 4ac}}{2a}.
$$

Some sources, particularly older ones, use alternative parameterizations of the quadratic equation such as $ax^2 + 2bx + c = 0$ or $ax^2 - 2bx + c = 0$, where $b$ has a magnitude one half of the more common one, possibly with opposite sign. These result in slightly different forms for the solution, but are otherwise equivalent.

A number of alternative derivations can be found in the literature. These proofs are simpler than the standard completing the square method, represent interesting applications of other frequently used techniques in algebra, or offer insight into other areas of mathematics.

A lesser known quadratic formula, as used in Muller''''s method, provides the same roots via the equation:

$$
x = \frac{2c}{-b \pm \sqrt{b^2 - 4ac}}.
$$

This can be deduced from the standard quadratic formula by Vieta''''s formulas, which assert that the product of the roots is $c/a$.

One property of this form is that it yields one valid root when $a = 0$, while the other root contains division by zero, because when $a = 0$, the quadratic equation becomes a linear equation, which has one root. By contrast, in this case, the more common formula has a division by zero for one root and an indeterminate form $0/0$ for the other root. On the other hand, when $c = 0$, the more common formula yields two correct roots whereas this form yields the zero root and an indeterminate form $0/0$.

\paragraph{Reduced Quadratic Equation}

It is sometimes convenient to reduce a quadratic equation so that its leading coefficient is one. This is done by dividing both sides by $a$, which is always possible since $a$ is non-zero. This produces the reduced quadratic equation:

$$
x^2 + px + q = 0,
$$

where $p = b/a$ and $q = c/a$. This monic polynomial equation has the same solutions as the original.

The quadratic formula for the solutions of the reduced quadratic equation, written in terms of its coefficients, is:

$$
x = \frac{1}{2} \left(-p \pm \sqrt{p^2 - 4q}\right),
$$

or equivalently:

$$
x = -\frac{p}{2} \pm \sqrt{\left(\frac{p}{2}\right)^2 - q}.
$$

\paragraph{Discriminant}

This figure plots three quadratic functions on a single Cartesian plane graph to illustrate the effects of discriminant values. When the discriminant, delta, is positive, the parabola intersects the x-axis at two points. When delta is zero, the vertex of the parabola touches the x-axis at a single point. When delta is negative, the parabola does not intersect the x-axis at all.

% [Image: graph]

In the quadratic formula, the expression underneath the square root sign is called the discriminant of the quadratic equation, and is often represented using an upper case D or an upper case Greek delta:

$$\Delta = b^2 - 4ac.$$

A quadratic equation with real coefficients can have either one or two distinct real roots, or two distinct complex roots. In this case the discriminant determines the number and nature of the roots. There are three cases:

1. If the discriminant is positive, then there are two distinct roots
   $$\frac{-b + \sqrt{\Delta}}{2a} \text{ and } \frac{-b - \sqrt{\Delta}}{2a},$$
   both of which are real numbers. For quadratic equations with rational coefficients, if the discriminant is a square number, then the roots are rational—in other cases they may be quadratic irrationals.

2. If the discriminant is zero, then there is exactly one real root
   $$-\frac{b}{2a},$$
   sometimes called a repeated or double root.

3. If the discriminant is negative, then there are no real roots. Rather, there are two distinct (non-real) complex roots
   $$-\frac{b}{2a} + i \frac{\sqrt{-\Delta}}{2a} \text{ and } -\frac{b}{2a} - i \frac{\sqrt{-\Delta}}{2a},$$
   which are complex conjugates of each other. In these expressions, $i$ is the imaginary unit.

Thus the roots are distinct if and only if the discriminant is non-zero, and the roots are real if and only if the discriminant is non-negative.

\paragraph{Geometric Interpretation}

Graph of 
$$y = ax^2 + bx + c,$$ 
where $a$ and the discriminant 
$$b^2 - 4ac$$ 
are positive, with:
\begin{itemize}
    \item Roots and y-intercept in red
    \item Vertex and axis of symmetry in blue
    \item Focus and directrix in pink
\end{itemize}

% [Image: graph]

The function $f(x) = ax^2 + bx + c$ is a quadratic function. The graph of any quadratic function has the same general shape, which is called a parabola. The location and size of the parabola, and how it opens, depend on the values of $a$, $b$, and $c$. As shown in Figure 1, if $a > 0$, the parabola has a minimum point and opens upward. If $a < 0$, the parabola has a maximum point and opens downward. The extreme point of the parabola, whether minimum or maximum, corresponds to its vertex. The x-coordinate of the vertex will be located at 
$$x = -\frac{b}{2a},$$ 
and the y-coordinate of the vertex may be found by substituting this x-value into the function. The y-intercept is located at the point $(0, c)$.

The solutions of the quadratic equation $ax^2 + bx + c = 0$ correspond to the roots of the function $f(x) = ax^2 + bx + c$, since they are the values of $x$ for which $f(x) = 0$. As shown in Figure 2, if $a$, $b$, and $c$ are real numbers and the domain of $f$ is the set of real numbers, then the roots of $f$ are exactly the x-coordinates of the points where the graph touches the x-axis. As shown in Figure 3, if the discriminant is positive, the graph touches the x-axis at two points; if zero, the graph touches at one point; and if negative, the graph does not touch the x-axis.

\paragraph{Quadratic Factorization}

The term 
$$x - r$$ 
is a factor of the polynomial 
$$ax^2 + bx + c$$ 
if and only if $r$ is a root of the quadratic equation 
$$ax^2 + bx + c = 0.$$ 
It follows from the quadratic formula that 
$$ax^2 + bx + c = a \left(x - \frac{-b + \sqrt{b^2 - 4ac}}{2a}\right) \left(x - \frac{-b - \sqrt{b^2 - 4ac}}{2a}\right).$$

In the special case $b^2 = 4ac$ where the quadratic has only one distinct root (i.e., the discriminant is zero), the quadratic polynomial can be factored as 
$$ax^2 + bx + c = a \left(x + \frac{b}{2a}\right)^2.$$

\paragraph{Graphical Solution}

A quadratic function without real root: $$y = (x - 5)^2 + 9.$$ The "3" is the imaginary part of the x-intercept. The real part is the x-coordinate of the vertex. Thus the roots are $$5 \pm 3i.$$
% [Image: graph]

The solutions of the quadratic equation
$$ax^2 + bx + c = 0$$
may be deduced from the graph of the quadratic function
$$y = ax^2 + bx + c,$$
which is a parabola.

\begin{itemize}
    \item If the parabola intersects the x-axis in two points, there are two real roots, which are the x-coordinates of these two points (also called x-intercept).
    \item If the parabola is tangent to the x-axis, there is a double root, which is the x-coordinate of the contact point between the graph and parabola.
    \item If the parabola does not intersect the x-axis, there are two complex conjugate roots. Although these roots cannot be visualized on the graph, their real and imaginary parts can be.
\end{itemize}

Let $h$ and $k$ be respectively the x-coordinate and the y-coordinate of the vertex of the parabola (that is the point with maximal or minimal y-coordinate). The quadratic function may be rewritten
$$y = a(x - h)^2 + k.$$

Let $d$ be the distance between the point of y-coordinate $2k$ on the axis of the parabola, and a point on the parabola with the same y-coordinate (see the figure; there are two such points, which give the same distance, because of the symmetry of the parabola). Then the real part of the roots is $h$, and their imaginary parts are $\pm d$. That is, the roots are
$$h + id \quad \text{and} \quad h - id,$$
or in the case of the example of the figure
$$5 + 3i \quad \text{and} \quad 5 - 3i.$$

\paragraph{Avoiding Loss of Significance}

Although the quadratic formula provides an exact solution, the result is not exact if real numbers are approximated during the computation, as usual in numerical analysis, where real numbers are approximated by floating point numbers (called "reals" in many programming languages). In this context, the quadratic formula is not completely stable.

This occurs when the roots have different orders of magnitude, or, equivalently, when $b^2$ and $b^2 - 4ac$ are close in magnitude. In this case, the subtraction of two nearly equal numbers will cause loss of significance or catastrophic cancellation in the smaller root. To avoid this, the root that is smaller in magnitude, $r$, can be computed as 
$$\frac{c / a}{R}$$ 
where $R$ is the root that is bigger in magnitude.

A second form of cancellation can occur between the terms $b^2$ and $4ac$ of the discriminant, that is when the two roots are very close. This can lead to loss of up to half of correct significant figures in the roots.

\paragraph{Examples and Applications}

% [Image: figure]

The trajectory of the cliff jumper is parabolic because horizontal displacement is a linear function of time
$$x = v\_x t,$$
while vertical displacement is a quadratic function of time
$$y = \frac{1}{2} a t^2 + v\_y t + h.$$ 
As a result, the path follows the quadratic equation
$$y = \frac{a}{2v\_x^2} x^2 + \frac{v\_y}{v\_x} x + h,$$
where $v\_x$ and $v\_y$ are the horizontal and vertical components of the original velocity, $a$ is gravitational acceleration, and $h$ is the original height. The $a$ value should be considered negative here, as its direction (downwards) is opposite to the height measurement (upwards).

The golden ratio is found as the positive solution of the quadratic equation
$$x^2 - x - 1 = 0.$$

The equations of the circle and the other conic sections—ellipses, parabolas, and hyperbolas—are quadratic equations in two variables.

Given the cosine or sine of an angle, finding the cosine or sine of the angle that is half as large involves solving a quadratic equation.

The process of simplifying expressions involving the square root of an expression involving the square root of another expression involves finding the two solutions of a quadratic equation.

Descartes'''' theorem states that for every four kissing (mutually tangent) circles, their radii satisfy a particular quadratic equation.

The equation given by Fuss'''' theorem, giving the relation among the radius of a bicentric quadrilateral''''s inscribed circle, the radius of its circumscribed circle, and the distance between the centers of those circles, can be expressed as a quadratic equation for which the distance between the two circles'''' centers in terms of their radii is one of the solutions. The other solution of the same equation in terms of the relevant radii gives the distance between the circumscribed circle''''s center and the center of the excircle of an ex-tangential quadrilateral.

Critical points of a cubic function and inflection points of a quartic function are found by solving a quadratic equation.

\paragraph{Resources}

\begin{itemize}
    \item \href{https://mathresearch.utsa.edu/wikiFiles/MAT1053/Quadratic_Functions/MAT1053_M2.2Quadratic_Functions.pdf}{Quadratic Functions} Book Chapters
\end{itemize}

\begin{itemize}
    \item \href{https://mathresearch.utsa.edu/wikiFiles/MAT1053/Quadratic_Functions/MAT1053_M2.2Quadratic_FunctionsGN.pdf}{Guided Notes}
\end{itemize}

\paragraph{Licensing}

Content obtained and/or adapted from:

\begin{itemize}
    \item \href{https://en.wikipedia.org/wiki/Quadratic_equation}{Quadratic equation, Wikipedia} under a CC BY-SA license
\end{itemize}

\subsection{Learning Objectives}

\begin{enumerate}
    \item 1.\textbf{Understanding and Applying the Quadratic Formula:}
    \item Students will be able to derive and apply the quadratic formula $$ x = \frac{-b \pm \sqrt{b^2 - 4ac}}{2a} $$ to solve any quadratic equation of the form $ax^2 + bx + c = 0$, identifying the nature of the roots based on the discriminant ($\Delta = b^2 - 4ac$). 2.\textbf{Factoring Quadratic Equations:}
    \item Students will be able to factor quadratic equations of the form $ax^2 + bx + c = 0$ into the product of two binomials $(px + q)(rx + s) = 0$, using methods such as factoring by inspection and Vieta's formulas when applicable. 3.\textbf{Solving Quadratic Equations by Completing the Square:}
    \item Students will be able to solve quadratic equations by completing the square, transforming equations into the form $$ (x + h)^2 = k $$ and subsequently solving for $x$ to find the roots of the equation.
\end{enumerate}

\subsection{Assessment Instrument}

\begin{enumerate}
    \item Instructions Solve the following quadratic equations using the methods specified. Show all steps clearly and indicate whether the roots are real or complex. 1.\textbf{Solve by Using the Quadratic Formula:}
    \item $( 3x^2 - 5x + 2 = 0 )$ 2.\textbf{Factor by Inspection:}
    \item $( x^2 + 7x + 10 = 0 )$ 3.\textbf{Solve by Completing the Square:}
    \item $( x^2 - 6x + 5 = 0 )$ Rubric 1.\textbf{Quadratic Formula:}
    \item Correct application of the quadratic formula.
    \item Accurate calculation of the discriminant.
    \item Correct identification of the roots.
    \item Proper interpretation of complex roots, if applicable. 2.\textbf{Factoring by Inspection:}
    \item Accurate identification of factors.
    \item Proper verification by expanding the factored form.
    \item Correct roots based on the factored form. 3.\textbf{Completing the Square:}
    \item Accurate completion of the square process.
    \item Proper simplification and solving of the resulting equation.
    \item Correct roots of the quadratic equation.
\end{enumerate}

%%===============================================================
\section{Other Types of Equations}
\label{sec:other-types-of-equations}
%%===============================================================

\paragraph{Other Types of Equations}

\paragraph{Square Root}

The graph of the function 
$$
f(x) = \sqrt{x}
$$
is made up of half a parabola with a vertical directrix.

% [Image: graph]

In mathematics, a square root of a number $x$ is a number $y$ such that $y^2 = x$; in other words, a number $y$ whose square (the result of multiplying the number by itself, or $y \cdot y$) is $x$. For example, 4 and -4 are square roots of 16, because $4^2 = (-4)^2 = 16$. Every nonnegative real number $x$ has a unique nonnegative square root, called the principal square root, which is denoted by 
$$
\sqrt{x}
$$
where the symbol 
$$
\sqrt{}
$$
is called the radical sign or radix. For example, the principal square root of 9 is 3, which is denoted by 
$$
\sqrt{9} = 3
$$
because $3^2 = 3 \cdot 3 = 9$ and 3 is nonnegative. The term (or number) whose square root is being considered is known as the radicand. The radicand is the number or expression underneath the radical sign, in this case 9.

Every positive number $x$ has two square roots: 
$$
\sqrt{x}
$$
which is positive, and 
$$
-\sqrt{x}
$$
which is negative. Together, these two roots are denoted as 
$$
\pm \sqrt{x}
$$
(see $\pm$ shorthand). Although the principal square root of a positive number is only one of its two square roots, the designation "the square root" is often used to refer to the principal square root. For positive $x$, the principal square root can also be written in exponent notation, as $x^{1/2}$.

Square roots of negative numbers can be discussed within the framework of complex numbers. More generally, square roots can be considered in any context in which a notion of the "square" of a mathematical object is defined. These include function spaces and square matrices, among other mathematical structures.

\paragraph{Cube Root}

Plot of 
$$
y = \sqrt[3]{x}
$$
The plot is symmetric with respect to the origin, as it is an odd function. At $x = 0$ this graph has a vertical tangent.

% [Image: graph]

In mathematics, a cube root of a number $x$ is a number $y$ such that $y^3 = x$. All nonzero real numbers have exactly one real cube root and a pair of complex conjugate cube roots, and all nonzero complex numbers have three distinct complex cube roots. For example, the real cube root of 8, denoted 
$$
\sqrt[3]{8}
$$
is 2, because $2^3 = 8$, while the other cube roots of 8 are 
$$
-1 + i\sqrt{3}
$$
and 
$$
-1 - i\sqrt{3}.
$$
The three cube roots of $-27i$ are
$$
3i,\quad \frac{3\sqrt{3}}{2} - \frac{3}{2}i,\quad \text{and}\quad -\frac{3\sqrt{3}}{2} - \frac{3}{2}i.
$$

In some contexts, particularly when the number whose cube root is to be taken is a real number, one of the cube roots (in this particular case the real one) is referred to as the principal cube root, denoted with the radical sign 
$$
\sqrt[3]{~}
$$
The cube root is the inverse function of the cube function if considering only real numbers, but not if considering also complex numbers: although one has always
$$
\left(\sqrt[3]{x}\right)^3 = x,
$$
the cube root of the cube of a number is not always this number. For example,
$$
-1 + i\sqrt{3}
$$
is a cube root of 8, (that is,
$$
(-1 + i\sqrt{3})^3 = 8
$$
), but
$$
-1 + i\sqrt{3} \neq 2 = \sqrt[3]{(-1 + i\sqrt{3})^3}.
$$

\paragraph{Rational Equations}

In mathematics, a rational function is any function that can be defined by a rational fraction, which is an algebraic fraction such that both the numerator and the denominator are polynomials. The coefficients of the polynomials need not be rational numbers; they may be taken in any field $K$. In this case, one speaks of a rational function and a rational fraction over $K$. The values of the variables may be taken in any field $L$ containing $K$. Then the domain of the function is the set of the values of the variables for which the denominator is not zero, and the codomain is $L$.

The set of rational functions over a field $K$ is a field, the field of fractions of the ring of the polynomial functions over $K$.

A function 
$$
f(x)
$$
is called a rational function if and only if it can be written in the form
$$
f(x) = \frac{P(x)}{Q(x)}
$$
where $P$ and $Q$ are polynomial functions of $x$ and $Q$ is not the zero function. The domain of $f$ is the set of all values of $x$ for which the denominator $Q(x)$ is not zero.

However, if $P$ and $Q$ have a non-constant polynomial greatest common divisor $R$, then setting 
$$
P = P\_{1}R
$$
and 
$$
Q = Q\_{1}R
$$
produces a rational function
$$
f\_{1}(x) = \frac{P\_{1}(x)}{Q\_{1}(x)},
$$
which may have a larger domain than $f(x)$, and is equal to $f(x)$ on the domain of $f(x)$. It is a common usage to identify $f(x)$ and $f\_{1}(x)$, that is to extend "by continuity" the domain of $f(x)$ to that of $f\_{1}(x)$. Indeed, one can define a rational fraction as an equivalence class of fractions of polynomials, where two fractions
$$
\frac{A(x)}{B(x)}
$$
and
$$
\frac{C(x)}{D(x)}
$$
are considered equivalent if
$$
A(x)D(x) = B(x)C(x).
$$
In this case, 
$$
\frac{P(x)}{Q(x)}
$$
is equivalent to
$$
\frac{P\_{1}(x)}{Q\_{1}(x)}.
$$

A proper rational function is a rational function in which the degree of $P(x)$ is less than the degree of $Q(x)$ and both are real polynomials, named by analogy to a proper fraction in $\mathbb{Q}$.

\paragraph{Degree}

There are several non-equivalent definitions of the degree of a rational function.

Most commonly, the degree of a rational function is the maximum of the degrees of its constituent polynomials $P$ and $Q$, when the fraction is reduced to lowest terms. If the degree of $f$ is $d$, then the equation
$$
f(z) = w
$$
has $d$ distinct solutions in $z$ except for certain values of $w$, called critical values, where two or more solutions coincide or where some solution is rejected at infinity (that is, when the degree of the equation decreases after having cleared the denominator).

In the case of complex coefficients, a rational function with degree one is a Möbius transformation.

The degree of the graph of a rational function is not the degree as defined above: it is the maximum of the degree of the numerator and one plus the degree of the denominator.

In some contexts, such as in asymptotic analysis, the degree of a rational function is the difference between the degrees of the numerator and the denominator.

In network synthesis and network analysis, a rational function of degree two (that is, the ratio of two polynomials of degree at most two) is often called a biquadratic function.

\paragraph{Examples}

The rational function
$$
f(x) = \frac{x^{3} - 2x}{2(x^{2} - 5)}
$$
is not defined at
$$
x^{2} = 5 \Leftrightarrow x = \pm \sqrt{5}.
$$
It is asymptotic to
$$
\frac{x}{2}
$$
as $x \to \infty$.

The rational function
$$
f(x) = \frac{x^{2} + 2}{x^{2} + 1}
$$
is defined for all real numbers, but not for all complex numbers, since if $x$ were a square root of $-1$ (i.e., the imaginary unit or its negative), then formal evaluation would lead to division by zero:
$$
f(i) = \frac{i^{2} + 2}{i^{2} + 1} = \frac{-1 + 2}{-1 + 1} = \frac{1}{0},
$$
which is undefined.

A constant function such as $f(x) = \pi$ is a rational function since constants are polynomials. The function itself is rational, even though the value of $f(x)$ is irrational for all $x$.

Every polynomial function
$$
f(x) = P(x)
$$
is a rational function with
$$
Q(x) = 1.
$$
A function that cannot be written in this form, such as
$$
f(x) = \sin(x),
$$
is not a rational function. However, the adjective "irrational" is not generally used for functions.

The rational function
$$
f(x) = \frac{x}{x}
$$
is equal to 1 for all $x$ except 0, where there is a removable singularity. The sum, product, or quotient (excepting division by the zero polynomial) of two rational functions is itself a rational function. However, the process of reduction to standard form may inadvertently result in the removal of such singularities unless care is taken. Using the definition of rational functions as equivalence classes gets around this, since $\frac{x}{x}$ is equivalent to $\frac{1}{1}$.

\paragraph{Absolute Value Equations}

In mathematics, the absolute value or modulus of a real number $x$, denoted $|x|$, is the non-negative value of $x$ without regard to its sign. Namely, $|x| = x$ if $x$ is positive, and $|x| = -x$ if $x$ is negative (in which case $-x$ is positive), and $|0| = 0$. For example, the absolute value of 3 is 3, and the absolute value of -3 is also 3. The absolute value of a number may be thought of as its distance from zero.

Generalizations of the absolute value for real numbers occur in a wide variety of mathematical settings. For example, an absolute value is also defined for the complex numbers, the quaternions, ordered rings, fields, and vector spaces. The absolute value is closely related to the notions of magnitude, distance, and norm in various mathematical and physical contexts.

The real absolute value function is continuous everywhere. It is differentiable everywhere except for $x = 0$. It is monotonically decreasing on the interval $(-\infty, 0]$ and monotonically increasing on the interval $[0, +\infty)$. Since a real number and its opposite have the same absolute value, it is an even function, and is hence not invertible. The real absolute value function is a piecewise linear, convex function.

Both the real and complex functions are idempotent.

\paragraph{Polynomial Equations}

A polynomial function is a function that can be defined by evaluating a polynomial. More precisely, a function $f$ of one argument from a given domain is a polynomial function if there exists a polynomial
$$
a\_{n}x^{n} + a\_{n-1}x^{n-1} + \cdots + a\_{2}x^{2} + a\_{1}x + a\_{0}
$$
that evaluates to $f(x)$ for all $x$ in the domain of $f$ (here, $n$ is a non-negative integer and $a\_{0}, a\_{1}, a\_{2}, \ldots, a\_{n}$ are constant coefficients). Generally, unless otherwise specified, polynomial functions have complex coefficients, arguments, and values. In particular, a polynomial, restricted to have real coefficients, defines a function from the complex numbers to the complex numbers. If the domain of this function is also restricted to the reals, the resulting function is a real function that maps reals to reals.

For example, the function $f$, defined by
$$
f(x) = x^{3} - x,
$$
is a polynomial function of one variable. Polynomial functions of several variables are similarly defined, using polynomials in more than one indeterminate, as in
$$
f(x, y) = 2x^{3} + 4x^{2}y + xy^{5} + y^{2} - 7.
$$
According to the definition of polynomial functions, there may be expressions that obviously are not polynomials but nevertheless define polynomial functions. An example is the expression
$$
\left(\sqrt{1 - x^{2}}\right)^{2},
$$
which takes the same values as the polynomial $1 - x^{2}$ on the interval $[-1, 1]$, and thus both expressions define the same polynomial function on this interval.

Every polynomial function is continuous, smooth, and entire.

\paragraph{Graphs}

A polynomial function in one real variable can be represented by a graph.

\begin{itemize}
    \item The graph of the zero polynomial $f(x) = 0$ is the $x$-axis.
    \item The graph of a degree 0 polynomial $f(x) = a_{0}$, where $a_{0} \neq 0$, is a horizontal line with $y$-intercept $a\_{0}$.
\end{itemize}

% [Image: graph]

\begin{itemize}
    \item The graph of a degree 1 polynomial (or linear function) $f(x) = a\_{0} + a\_{1}x$, where $a\_{1} \neq 0$, is an oblique line with $y$-intercept $a\_{0}$ and slope $a\_{1}$.
\end{itemize}

% [Image: graph]

\begin{itemize}
    \item The graph of a degree 2 polynomial $f(x) = a\_{0} + a\_{1}x + a\_{2}x^{2}$, where $a\_{2} \neq 0$, is a parabola.
\end{itemize}

% [Image: graph]

\begin{itemize}
    \item The graph of a degree 3 polynomial $f(x) = a\_{0} + a\_{1}x + a\_{2}x^{2} + a\_{3}x^{3}$, where $a\_{3} \neq 0$, is a cubic curve.
\end{itemize}

% [Image: graph]

\begin{itemize}
    \item The graph of any polynomial with degree 2 or greater $f(x) = a\_{0} + a\_{1}x + a\_{2}x^{2} + \cdots + a\_{n}x^{n}$, where $a\_{n} \neq 0$ and $n \geq 2$, is a continuous non-linear curve.
\end{itemize}

A non-constant polynomial function tends to infinity when the variable increases indefinitely (in absolute value). If the degree is higher than one, the graph does not have any asymptote. It has two parabolic branches with vertical direction (one branch for positive $x$ and one for negative $x$).

Polynomial graphs are analyzed in calculus using intercepts, slopes, concavity, and end behavior.

\paragraph{Resources}

\begin{itemize}
    \item \href{https://courses.lumenlearning.com/waymakercollegealgebra/chapter/writing-formulas-for-polynomial-functions/}{Polynomial Equations, Lumen Learning}
    \item \href{https://courses.lumenlearning.com/suny-beginalgebra/chapter/7-3-1-solving-radical-equations/}{Radical Equations, Lumen Learning}
    \item \href{https://courses.lumenlearning.com/ivytech-collegealgebra/chapter/solving-equations-involving-rational-exponents/}{Rational Exponent Equations, Lumen Learning}
    \item \href{https://courses.lumenlearning.com/ivytech-collegealgebra/chapter/solve-an-absolute-value-equation/}{Absolute Value Equation, Lumen Learning}
\end{itemize}

\paragraph{Licensing}

Content obtained and/or adapted from:

\begin{itemize}
    \item \href{https://en.wikipedia.org/wiki/Polynomial}{Polynomial, Wikipedia} under a CC BY-SA license
    \item \href{https://en.wikipedia.org/wiki/Square_root}{Square root, Wikipedia} under a CC BY-SA license
    \item \href{https://en.wikipedia.org/wiki/Cube_root}{Cube root, Wikipedia} under a CC BY-SA license
    \item \href{https://en.wikipedia.org/wiki/Rational_function}{Rational function, Wikipedia} under a CC BY-SA license
    \item \href{https://en.wikipedia.org/wiki/Absolute_value}{Absolute value, Wikipedia} under a CC BY-SA license
\end{itemize}

\subsection{Learning Objectives}

\begin{enumerate}
    \item \textbf{Understand and Identify Square Roots} Students will be able to determine the principal square root of a nonnegative real number $x$, as well as identify both the positive and negative square roots of $x$. They should also be able to express the square root of a number in radical notation and exponential notation. \textbf{Example Objective:} Given a number $x$, find $\sqrt{x}$ and $-\sqrt{x}$, and express $\sqrt{x}$ as $x^{1/2}$.
    \item \textbf{Evaluate and Apply Cube Roots} Students will be able to find the real cube root of a number $x$, identify the cube roots of complex numbers, and understand the properties of the cube root function. They will also be able to recognize that a cube root function is the inverse of the cube function. \textbf{Example Objective:} Find $\sqrt[3]{27}$, and list all three cube roots of $-27i$.
    \item \textbf{Analyze Rational Functions} Students will be able to analyze and simplify rational functions, including identifying their domains, asymptotic behavior, and critical values. They should also understand the concepts of proper rational functions and degrees of rational functions. \textbf{Example Objective:} For the rational function $f(x) = \frac{x^{3} - 2x}{2(x^{2} - 5)}$, determine the values where the function is undefined and describe its asymptotic behavior as $x \to \infty$.
\end{enumerate}

\subsection{Assessment Instrument}

\begin{enumerate}
    \item \textbf{Instructions:} Answer the following questions and show all work for full credit.
    \item \textbf{Square Roots} a. Compute the principal square root of 16 and express it in both radical and exponential notation. b. Find both the positive and negative square roots of 25. c. Solve the equation $\sqrt{x} = 4$ for $x$.
    \item \textbf{Cube Roots} a. Determine the real cube root of 64 and express it in radical notation. b. Find all three cube roots of $-8$. c. For the complex number $-1 + i\sqrt{3}$, verify that it is a cube root of 8 by computing $(-1 + i\sqrt{3})^3$.
    \item \textbf{Rational Functions} a. Consider the rational function $f(x) = \frac{x^{2} - 1}{x - 1}$. Simplify the function and state its domain. b. Find the vertical and horizontal asymptotes of the function $g(x) = \frac{3x^{2} + 2}{x^{2} - 4}$. c. Given the function $h(x) = \frac{x^{2} - 4}{x^{2} - 1}$, determine if there are any removable singularities.
    \item \textbf{Absolute Value Equations} a. Solve the absolute value equation $|x - 5| = 7$ for $x$. b. Graph the function $f(x) = |x + 2|$ and describe its key features.
    \item \textbf{Polynomial Equations} a. Find the roots of the polynomial $p(x) = x^{2} - 9$. b. Sketch the graph of the polynomial function $q(x) = x^{3} - 3x + 2$ and identify its key features such as intercepts and turning points. c. Determine the degree of the polynomial $r(x) = 2x^{4} - x^{3} + 3x^{2} - x + 1$.
\end{enumerate}

%%===============================================================
\section{Linear and Absolute Value Inequalities}
\label{sec:linear-and-absolute-value-inequalities}
%%===============================================================

\paragraph{Systems of Inequalities}



In mathematics, a linear inequality is an inequality that involves a linear function. A linear inequality contains one of the symbols of inequality: It shows the data which is not equal in graph form.



\begin{itemize}
    \item $<$ less than
    \item $>$ greater than
    \item $\leq$ less than or equal to
    \item $\geq$ greater than or equal to
    \item $\neq$ not equal to
    \item $=$ equal to
\end{itemize}

A linear inequality looks exactly like a linear equation, with the inequality sign replacing the equality sign.



\paragraph{Contents}





\paragraph{Linear inequalities of real numbers}



\paragraph{Two-dimensional linear inequalities}







Two-dimensional linear inequalities are expressions in two variables of the form:



$$

a x + b y < c \text{ and } a x + b y \geq c

$$



where the inequalities may either be strict or not. The solution set of such an inequality can be graphically represented by a half-plane (all the points on one "side" of a fixed line) in the Euclidean plane. Technically, for this statement to be correct both $a$ and $b$ cannot simultaneously be zero. In that situation, the solution set is either empty or the entire plane. The line that determines the half-planes ($a x + b y = c$) is not included in the solution set when the inequality is strict. A simple procedure to determine which half-plane is in the solution set is to calculate the value of $a x + b y$ at a point $(x\_0, y\_0)$ which is not on the line and observe whether or not the inequality is satisfied.



For example, to draw the solution set of $x + 3y < 9$, one first draws the line with equation $x + 3y = 9$ as a dotted line, to indicate that the line is not included in the solution set since the inequality is strict. Then, pick a convenient point not on the line, such as $(0,0)$. Since $0 + 3(0) = 0 < 9$, this point is in the solution set, so the half-plane containing this point (the half-plane "below" the line) is the solution set of this linear inequality.



Graph of linear inequality: $x + 3y < 9$



% [Image: graph]



\paragraph{Linear inequalities in general dimensions}



In $\mathbb{R}^n$, linear inequalities are the expressions that may be written in the form:



$$

f(\bar{x}) < b \text{ or } f(\bar{x}) \leq b

$$



where $f$ is a linear form (also called a linear functional), $\bar{x} = (x\_1, x\_2, \ldots, x\_n)$, and $b$ is a constant real number.



More concretely, this may be written out as:



$$

a\_1 x\_1 + a\_2 x\_2 + \cdots + a\_n x\_n < b

$$



or



$$

a\_1 x\_1 + a\_2 x\_2 + \cdots + a\_n x\_n \leq b

$$



Here $x\_1, x\_2, \ldots, x\_n$ are called the unknowns, and $a\_1, a\_2, \ldots, a\_n$ are called the coefficients.



Alternatively, these may be written as:



$$

g(x) < 0

$$



or



$$

g(x) \leq 0

$$



where $g$ is an affine function.



That is:



$$

a\_0 + a\_1 x\_1 + a\_2 x\_2 + \cdots + a\_n x\_n < 0

$$



or



$$

a\_0 + a\_1 x\_1 + a\_2 x\_2 + \cdots + a\_n x\_n \leq 0

$$



Note that any inequality containing a "greater than" or a "greater than or equal" sign can be rewritten with a "less than" or "less than or equal" sign, so there is no need to define linear inequalities using those signs.



\paragraph{Systems of linear inequalities}



A system of linear inequalities is a set of linear inequalities in the same variables:



$ a\_{11} x\_1 + a\_{12} x\_2  + \cdots + a\_{1n} x\_n \leq b\_1 $



$ a\_{21} x\_1 + a\_{22} x\_2 + \cdots + a\_{2n} x\_n \leq b\_2 $ 



$             \vdots                         \vdots                          \vdots                          \vdots                          \vdots $



$ a\_{m1} x\_1 + a\_{m2} x\_2 + \cdots + a\_{mn} x\_n \leq b\_m $ 





Here $x\_1, x\_2, \ldots, x\_n$ are the unknowns, $a\_{11}, a\_{12}, \ldots, a\_{mn}$ are the coefficients of the system, and $b\_1, b\_2, \ldots, b\_m$ are the constant terms.



This can be concisely written as the matrix inequality:



$$

Ax \leq b

$$



where $A$ is an $m \times n$ matrix, $x$ is an $n \times 1$ column vector of variables, and $b$ is an $m \times 1$ column vector of constants.



In the above systems, both strict and non-strict inequalities may be used.



Not all systems of linear inequalities have solutions. Variables can be eliminated from systems of linear inequalities using Fourier–Motzkin elimination.



\paragraph{Applications}



\paragraph{Polyhedra}



The set of solutions of a real linear inequality constitutes a half-space of the $n$-dimensional real space, one of the two defined by the corresponding linear equation.



The set of solutions of a system of linear inequalities corresponds to the intersection of the half-spaces defined by individual inequalities. It is a convex set, since the half-spaces are convex sets, and the intersection of a set of convex sets is also convex. In the non-degenerate cases, this convex set is a convex polyhedron (possibly unbounded, e.g., a half-space, a slab between two parallel half-spaces, or a polyhedral cone). It may also be empty or a convex polyhedron of lower dimension confined to an affine subspace of the $n$-dimensional space $\mathbb{R}^n$.



\paragraph{Linear programming}



A linear programming problem seeks to optimize (find a maximum or minimum value) a function (called the objective function) subject to a number of constraints on the variables which, in general, are linear inequalities. The list of constraints is a system of linear inequalities.

\subsection{Learning Objectives}

\begin{enumerate}
    \item \textbf{SLO 1:} Students will be able to define and identify linear inequalities and their symbols ($<$, $>$, $\leq$, $\geq$, $\neq$, $=$).
    \item \textbf{SLO 2:} Students will be able to graphically represent two-dimensional linear inequalities and determine the solution set of such inequalities.
    \item \textbf{SLO 3:} Students will be able to solve systems of linear inequalities and express the solution sets using matrix notation.
\end{enumerate}

\subsection{Assessment Instrument}

\begin{enumerate}
    \item Multiple-Choice Questions
    \item \textbf{Question 1:} Which of the following represents a linear inequality?
    \item A) $2x + 3 = 7$
    \item B) $4y - 5 > 2$
    \item C) $x^2 + y < 5$
    \item D) $7x \geq 4x - 9$
    \item \textbf{Question 2:} What is the solution set for the inequality $x + 3y < 9$ in the coordinate plane?
    \item A) The line $x + 3y = 9$
    \item B) The half-plane below the line $x + 3y = 9$
    \item C) The half-plane above the line $x + 3y = 9$
    \item D) The entire coordinate plane Short Answer Questions
    \item \textbf{Question 3:} Explain how you can determine which half-plane is included in the solution set of a linear inequality.
    \item \textbf{Question 4:} Write the system of linear inequalities in matrix form: $$ 2x + 3y \leq 5 $$ $$ -x + 4y > 3$$ Problem Solving
    \item \textbf{Question 5:} Solve the system of linear inequalities and graph the solution set: $$ x + y \leq 4 $$ $$ x - y \geq 1 $$
\end{enumerate}

%%===============================================================
\section{Functions}
\label{sec:functions}
%%===============================================================

\paragraph{Functions}

\paragraph{Introduction}

Formally, a function $f$ from a set $X$ to a set $Y$ is defined by a set $G$ of ordered pairs $(x, y)$ such that $x \in X$, $y \in Y$, and every element of $X$ is the first component of exactly one ordered pair in $G$. In other words, for every $x$ in $X$, there is exactly one element $y$ such that the ordered pair $(x, y)$ belongs to the set of pairs defining the function $f$. The set $G$ is called the graph of the function. Formally speaking, it may be identified with the function, but this hides the usual interpretation of a function as a process. Therefore, in common usage, the function is generally distinguished from its graph. Whenever one quantity uniquely determines the value of another quantity, we have a function. That is, the set $X$ uniquely determines the set $Y$. You can think of a function as a kind of machine. You feed the machine raw materials, and the machine changes the raw materials into a finished product.

\paragraph{Functions in Everyday Life}

Think about dropping a ball from a bridge. At each moment in time, the ball is a height above the ground. The height of the ball is a function of time. It was the job of physicists to come up with a formula for this function. This type of function is called real-valued since the "finished product" is a number (or, more specifically, a real number).

Think about a wind storm. At different places, the wind can be blowing in different directions with different intensities. The direction and intensity of the wind can be thought of as a function of position. This is a function of two real variables (a location is described by two values - an $x$ and a $y$) which results in a vector (which is something that can be used to hold a direction and an intensity). These functions are studied in multivariable calculus (which is usually studied after a one-year college-level calculus course). This is a vector-valued function of two real variables.

We will be looking at real-valued functions until studying multivariable calculus. Think of a real-valued function as an input-output machine; you give the function an input, and it gives you an output which is a number (more specifically, a real number). For example, the squaring function takes the input 4 and gives the output value 16. The same squaring function takes the input -1 and gives the output value 1.

This is an intuitive way to understand functions: a machine that makes the input $x$ go through a transformation $f$ into the output $f(x)$.

% [Image: image]

\paragraph{Notation}

Functions are used so much that there is a special notation for them. The notation is somewhat ambiguous, so familiarity with it is important in order to understand the intention of an equation or formula.

Though there are no strict rules for naming a function, it is standard practice to use the letters $f$, $g$, and $h$ to denote functions, and the variable $x$ to denote an independent variable. $y$ is used for both dependent and independent variables.

When discussing or working with a function $f$, it''s important to know not only the function but also its independent variable $x$. Thus, when referring to a function $f$, you usually do not write $f$, but instead $f(x)$. The function is now referred to as " $f$ of $x$ ". The name of the function is adjacent to the independent variable (in parentheses). This is useful for indicating the value of the function at a particular value of the independent variable. For instance, if 

$$
f(x) = 7x + 1
$$

and if we want to use the value of $f$ for $x$ equal to 2, then we would substitute 2 for $x$ on both sides of the definition above and write

$$
f(2) = 7(2) + 1 = 14 + 1 = 15
$$

This notation is more informative than leaving off the independent variable and writing simply '' $f$ '', but can be ambiguous since the parentheses next to $f$ can be misinterpreted as multiplication, $2f$. To make sure nobody is too confused, follow this procedure:

1. Define the function $f$ by equating it to some expression.
2. Give a sentence like the following: "At $x = c$, the function $f$ is..."
3. Evaluate the function.

\paragraph{Description}

There are many ways people describe functions. In the examples above, a verbal description is given (the height of the ball above the earth as a function of time). Here is a list of ways to describe functions. The top three listed approaches to describing functions are the most popular.

1. A function is given a name (such as $f$) and a formula for the function is also given. For example, $f(x) = 3x + 2$ describes a function. We refer to the input as the argument of the function (or the independent variable), and to the output as the value of the function at the given argument.
2. A function is described using an equation and two variables. One variable is for the input of the function and one is for the output of the function. The variable for the input is called the independent variable. The variable for the output is called the dependent variable. For example, $y = 3x + 2$ describes a function. The dependent variable appears by itself on the left-hand side of the equal sign.
3. A verbal description of the function.

When a function is given a name (like in number 1 above), the name of the function is usually a single letter of the alphabet (such as $f$ or $g$). Some functions whose names are multiple letters (like the sine function $y = \sin(x)$).

\paragraph{Plugging a Value into a Function}

If we write $f(x) = 3x + 2$, then we know that:

\begin{itemize}
    \item The function $f$ is a function of $x$.
    \item To evaluate the function at a certain number, replace the $x$ with that number.
    \item Replacing $x$ with that number in the right side of the function will produce the function''s output for that certain input.
\end{itemize}

In English, the definition of $f$ is interpreted, "Given a number, $f$ will return two more than the triple of that number."

How would we know the value of the function $f$ at 3? We would have the following three thoughts:

$$
f(3) = 3(3) + 2
$$

$$
3(3) + 2 = 9 + 2
$$

$$
9 + 2 = 11
$$

and we would write

$$
f(3) = 3(3) + 2 = 9 + 2 = 11
$$

The value of $f$ at 3 is 11.

Note that $f(3)$ means the value of the dependent variable when $x$ takes on the value of 3. So we see that the number 11 is the output of the function when we give the number 3 as the input. People often summarize the work above by writing "the value of $f$ at three is eleven", or simply "$f$ of three equals eleven".

\paragraph{Basic Concepts of Functions}

The formal definition of a function states that a function is actually a mapping that associates the elements of one set called the domain of the function, $A$, with the elements of another set called the range of the function, $B$. For each value we select from the domain of the function, there exists exactly one corresponding element in the range of the function. The definition of the function tells us which element in the range corresponds to the element we picked from the domain. An example is given below.

\paragraph{Example}

Let function

$$
f(x) = 5x + \frac{5x}{2}
$$

for all $x$. For what value of $x$ gives 

$$
f(x) = 225 \frac{1}{2}
$$

In mathematics, it is important to notice any repetition. If something seems to repeat, keep a note of that in the back of your mind somewhere.

Here, notice that 

$$
f(x) = 5x + \frac{5x}{2}
$$

and 

$$
f(x) = 225 \frac{1}{2}
$$

Because $f(x)$ is equal to two different things, it must be the case that the other side of the equal sign to $f(x)$ is equal to the other. This property is known as the transitive property and could thus make the following equation true:

$$
225 \frac{1}{2} = 5x + \frac{5x}{2}
$$

Next, simplify — make your life easier rather than harder. In this instance, since $5x + \frac{5x}{2}$ has $x$ as a like-term, and the two terms are fractions added to the other, put it over a common denominator and simplify. Similarly, since $225 \frac{1}{2}$ is a mixed fraction, $225 \frac{1}{2} = 225 + \frac{1}{2}$.

$$
225 \frac{1}{2} = 5x + \frac{5x}{2}
$$

$$
\Rightarrow \frac{225}{1} + \frac{1}{2} = \frac{5x}{1} + \frac{5x}{2}
$$

$$
\Rightarrow \frac{450}{2} + \frac{1}{2} = \frac{10x}{2} + \frac{5x}{2}
$$

$$
\Rightarrow \frac{451}{2} = \frac{15x}{2}
$$

Multiply both sides by the reciprocal of the coefficient of $x$, $\frac{2}{15}$ from both sides by multiplying by it.

$$
\frac{2}{15} \cdot \frac{451}{2} = x
$$

$$
\frac{451}{15} = x
$$

or 

$$
x = \frac{451}{15}
$$

The value of $x$ that makes 

$$
f(x) = 225 \frac{1}{2}
$$

is 

$$
x = \frac{451}{15}
$$

Classically, the element picked from the domain is pictured as something that is fed into the function and the corresponding element in the range is pictured as the output. Since we "pick" the element in the domain whose corresponding element in the range we want to find, we have control over what element we pick and hence this element is also known as the "independent variable". The element mapped in the range is beyond our control and is "mapped to" by the function. This element is hence also known as the "dependent variable", for it depends on which independent variable we pick. Since the elementary idea of functions is better understood from the classical viewpoint, we shall use it hereafter. However, it is still important to remember the correct definition of functions at all times.

\paragraph{Basic Types of Transformation}

To make it simple, for the function $f(x)$, all of the possible $x$ values constitute the domain, and all of the values $f(x)$ ($y$ on the $x$-$y$ plane) constitute the range. To put it in more formal terms, a function $f$ is a mapping of some element $a \in A$, called the domain, to exactly one element $b \in B$, called the range, such that $f:A \to B$. The image below should help explain the modern definition of a function:

The image demonstrates a mapping of some element $a$ (the circle) in $A$, the domain, to exactly one element $b$ in $B$, the range. $A$ is the domain of the function while $B$ is the range. This transformation from set $A$ to $B$ is an example of a one-to-one function.

% [Image: image]

A function is considered one-to-one if an element $a \in A$ from domain $A$ of function $f$ leads to exactly one element $b \in B$ from range $B$ of the function. By definition, since only one element $b$ is mapped by function $f$ from some element $a$, $f:A \to B$ implies that there exists only one element $b$ from the mapping. Therefore, there exists a one-to-one function because it complies with the definition of a function. 

A function is considered many-to-one if some elements $a\_1, a\_2, \ldots a\_n \in A$ from domain $A$ of function $f$ maps to exactly one element $b \in B$ from range $B$ of the function. Since some elements $(a\_1, a\_2, \ldots a\_n) \in A$ must map onto exactly one element $b \in B$, $f:A \to B$ must be compliant with the definition of a function. Therefore, there exists a many-to-one function.

A function is considered one-to-many if exactly one element $a \in A$ from domain $A$ of function $f$ maps to some elements $(b\_1, b\_2, \ldots b\_n) \in B$ from range $B$ of the function. If some element $a$ maps onto many distinct elements $b$, then $f:A \to B$ is non-functional since there exist many distinct elements $b$. Given many-to-one is non-compliant with the definition of a function, there exists no function that is one-to-many.

The modern definition describes a function sufficiently such that it helps us determine whether some new type of "function" is indeed one. Now that each case is defined above, you can now prove whether functions are one of these function cases. Here is an example problem:

\paragraph{One-to-One Example}

Given: $f(x) = ax + b$, where $a$ and $b$ are constant and $a \neq 0$.
Prove: function $f$ is one-to-one.

Notice that the only changing element in the function $f$ is $x$. To prove a function is one-to-one by applying the definition may be impossible because although two random elements of domain set $A$ may not be many-to-one, there may be some elements $(a\_1, a\_2, \ldots a\_n) \in A$ that may make the function many-to-one. To account for this possibility, we must prove that it is impossible to have some result like that.

Assume there exist two distinct elements $x\_1 \neq x\_2$ that will result in $f(x\_1) = f(x\_2)$. This would make the function many-to-one. In consequence,
$$
ax\_1 + b = ax\_2 + b
$$
Subtract $b$ from both sides of the equation.
$$
\Rightarrow ax\_1 = ax\_2
$$
Subtract $ax\_2$ from both sides of the equation.
$$
\Rightarrow ax\_1 - ax\_2 = 0
$$
Factor $a$ from both terms to the left of the equation.
$$
\Rightarrow a(x\_1 - x\_2) = 0
$$
Multiply $\frac{1}{a}$ to both sides of the equation.
$$
\Rightarrow (x\_1 - x\_2) = 0
$$
Add $x\_2$ to both sides of the equation.
$$
\Rightarrow x\_1 = x\_2
$$
Notice that $x\_1 = x\_2$. However, this is impossible because $x\_1$ and $x\_2$ are distinct. Ergo, $f(x\_1) \neq f(x\_2)$. No two distinct inputs can give the same output. Therefore, the function must be one-to-one.

\paragraph{Basic Concepts}

The domain is all the elements in set $A$ that can be mapped to the elements in set $B$. The range is those elements in set $B$ which are mapped with the domain. The codomain is all the elements in set $B$.

% [Image: image]

There are a few more important ideas to learn from the modern definition of the function, and it comes from knowing the difference between domain, range, and codomain. We already discussed some of these, yet knowing about sets adds a new definition for each of the following ideas. Therefore, let us discuss these based on these new ideas. Let $A$ and $B$ be a set. If we were to combine these two sets, it would be defined as the cartesian cross product $A \times B$. The subset of this product is the function. The below definitions are a little confusing. The best way to interpret these is to draw an image. To the right of these definitions is the image that corresponds to it.

\paragraph{Definition of Domain, Range, and Codomain of a Function}

The domain is defined to be a set $A$ with all elements $a \in A$ mapping to at least one unique $b \in B$.

The set of elements in $B$ is the range of the function mapping $f:A \to B$ in the cartesian cross product, whereby the set of all elements $a \in A$ maps to some element $b \in B$.

The codomain is the set $B$, where it is not necessarily the case that all elements $b \in B$ were mapped by some $a \in A$.

Note that the codomain is not as important as the other two concepts.

Take $f(x) = \sqrt{1 - x^2}$ for example:

  Because of the square root, the content in the square root should be strictly not smaller than 0.
  $$
  1 - x^2 \geq 0
  $$
  $$
  \Leftrightarrow -1 \leq x \leq 1
  $$
  Thus the domain is
  $$
  x \in [-1, 1]
  $$
\begin{itemize}
    \item \textbf{The domain of the function is the interval from -1 to 1}:
\end{itemize}
% [Image: domain]

  Correspondingly, the range will be
  $$
  f(x) \in [0, 1]
  $$

\begin{itemize}
    \item \textbf{The range of the function is the interval from 0 to 1}:
\end{itemize}
% [Image: range]

\paragraph{Other Types of Transformation}

There is one more final aspect that needs to be defined. We already have a good idea of what makes a mapping a function (e.g. it cannot be one-to-many). However, three more definitions that you will often hear will be necessary to talk about: injective, surjective, bijective.

% [Image: mapping]

\begin{itemize}
    \item The function mapping on the left is an example of an injective function because it is one-to-one. However, it is not surjective because the range and the codomain are not the same.
    \item A function is \textbf{injective} if it is one-to-one.
    \item A function is \textbf{surjective} if it is "onto." That is, the mapping $f:A \to B$ has $B$ as the range of the function, where the codomain and the range of the function are the same.
    \item A function is \textbf{bijective} if it is both surjective and injective.
\end{itemize}

Again, the above definitions are often very confusing. Another image is provided to the right of the bullet points. Along with that, another example is also provided. Let us analyze the function $g(x) = ax^2$.

Given: $g(x) = ax^2$, where $a$ is constant and $a \neq 0$.
Prove: function $g$ is non-surjective and non-injective.

Notice that the only changing element in the function $g$ is $x$. Let us check to see the conditions of this function.
Is it injective? Let $g(x) = y$, and solve for $x$. First, divide by $a$.

$$
y = ax^2
$$

$$
\Rightarrow \frac{y}{a} = x^2
$$

Then take the square root of $x^2$. By definition, 

$$
\sqrt{x^2} = \pm x
$$

so

$$
\frac{y}{a} = x^2 \Rightarrow \pm \sqrt{\frac{y}{a}} = \pm x
$$

Notice, however, what we learned from the above manipulation. As a result of solving for $x$, we found that there are two solutions for $x$. However, this resulted in two different values from $\frac{y}{a}$. This means that for some individual $x$ that gives $y$, there are two different inputs that result in the same value. Because $f(x_1) = f(x_2)$ when $x_1 \neq x_2$, this is by definition non-injective.

Is it surjective? As a natural consequence of what we learned about As a result of solving for $x$, we found that there are two solutions for $x$. However, this resulted in two different values from $\frac{y}{a}$. This means that for some individual $x$ that gives $y$, there are two different inputs that result in the same value. Because $f(x_1) = f(x_2)$ when $x_1 \neq x_2$, this is by definition non-injective.

Is it surjective? As a natural consequence of what we learned about inputs, let us determine the range of the function. After all, the only way to determine if something is surjective is to see if the range applies to all real numbers. A good way to determine this is by finding a pattern involving our domains. Let us say we input a negative number for the domain:

$$
f(-x) = a(-x)^2 = a \cdot (-x) \cdot (-x) = ax^2
$$

This seemingly trivial exercise tells us that negative numbers give us non-negative numbers for our range. This is huge information! For $x \in \mathbb{R}$, the function always results $y > 0$ for our range. For the set $B$, we have elements in that set that have no mappings from the set $A$. As such, $y \leq 0$ is the codomain of set $B$. Therefore, this function is non-surjective!

\paragraph{Tests for Equations}

\paragraph{The Vertical Line Test}

The vertical line test is a systematic test to find out if an equation involving $x$ and $y$ can serve as a function (with $x$ the independent variable and $y$ the dependent variable). Simply graph the equation and draw a vertical line through each point of the $x$-axis. If any vertical line ever touches the graph at more than one point, then the equation is not a function; if the line always touches at most one point of the graph, then the equation is a function.

The circle, below, is not a function because the vertical line intercepts two points on the graph when $x = 1$.

% [Image: vertical line test]

\paragraph{The Horizontal Line and the Algebraic 1-1 Test}

Similarly, the horizontal line test, though it does not test if an equation is a function, tests if a function is injective (one-to-one). If any horizontal line ever touches the graph at more than one point, then the function is not one-to-one; if the line always touches at most one point on the graph, then the function is one-to-one.

The algebraic 1-1 test is the non-geometric way to see if a function is one-to-one. The basic concept is that:

Assume there is a function $f$. If:

$$f(a) = f(b) \text{ and } a = b,$$

then

function $f$ is one-to-one.

Here is an example: prove that $f(x) = \frac{1 - 2x}{1 + x}$ is injective.

Since the notation is the notation for a function, the equation is a function. So we only need to prove that it is injective. Let $a$ and $b$ be the inputs of the function and that $f(a) = f(b)$. Thus,

$$
\frac{1 - 2a}{1 + a} = \frac{1 - 2b}{1 + b}
\Leftrightarrow (1 + b)(1 - 2a) = (1 + a)(1 - 2b)
\Leftrightarrow 1 - 2a + b - 2ab = 1 - 2b + a - 2ab
\Leftrightarrow 1 - 2a + b = 1 - 2b + a
\Leftrightarrow 1 - 2a + 3b = 1 + a
\Leftrightarrow 1 + 3b = 1 + 3a
\Leftrightarrow a = b
$$

So, the result is $a = b$, proving that the function is injective.

Another example is proving that $g(x) = x^2$ is not injective.

Using the same method, one can find that $a = \pm b$, which is not $a = b$. So, the function is not injective.

\paragraph{Remarks}

The following arise as a direct consequence of the definition of functions:

\begin{itemize}
    \item By definition, for each "input" a function returns only one "output", corresponding to that input. While the same output may correspond to more than one input, one input cannot correspond to more than one output. This is expressed graphically as the vertical line test: a line drawn parallel to the axis of the dependent variable (normally vertical) will intersect the graph of a function only once. However, a line drawn parallel to the axis of the independent variable (normally horizontal) may intersect the graph of a function as many times as it likes. Equivalently, this has an algebraic (or formula-based) interpretation. We can always say if $a = b$, then $f(a) = f(b)$, but if we only know that $f(a) = f(b)$ then we can''t be sure that $a = b$.
    \item Each function has a set of values, the function''s domain, which it can accept as input. Perhaps this set is all positive real numbers; perhaps it is the set {pork, mutton, beef}. This set must be implicitly/explicitly defined in the definition of the function. You cannot feed the function an element that isn''t in the domain, as the function is not defined for that input element.
    \item Each function has a set of values, the function''s range, which it can output. This may be the set of real numbers. It may be the set of positive integers or even the set {0, 1}. This set, too, must be implicitly/explicitly defined in the definition of the function.
\end{itemize}

Functions are an important foundation of mathematics. This modern interpretation is a result of the hard work of the mathematicians of history. It was not until recently that the definition of the relation was introduced by René Descartes in \textit{Geometry} (1637). Nearly a century later, the standard notation ($f(x) = y$) was first introduced by Leonhard Euler in \textit{Introductio in Analysin Infinitorum} and \textit{Institutiones Calculi Differentialis}. The term function was also a new innovation during Euler''s time as well, being introduced by Gottfried Wilhelm Leibniz in a 1673 letter "to describe a quantity related to points of a curve, such as a coordinate or curve''s slope." Finally, the modern definition of the function being the relation among sets was first introduced in 1908 by Godfrey Harold Hardy where there is a relation between two variables $x$ and $y$ such that "to some values of $x$ at any rate correspond values of $y$."

\paragraph{Important Functions}

\paragraph{Polynomials}

Polynomial functions are the most common and most convenient functions in the world of calculus. Their frequent appearances and their applications on computer calculations have made them important. A polynomial in a single indeterminate $x$ can always be written (or rewritten) in the form:

$$
f(x) = a\_n x^n + a\_{n-1} x^{n-1} + a\_{n-2} x^{n-2} + \cdots + a\_2 x^2 + a\_1 x + a\_0
$$

To be more concise, it can also be written in the summation form:

$$
\sum\_{k=0}^n a\_k x^k
$$

\paragraph{Constant}

When $n = 0$, the polynomial can be rewritten into the following function:

$$
f(x) = C,
$$

where $C$ is a constant.

The graph of this function is a horizontal line passing through the point $(0, C)$.

\paragraph{Linear}

Two linear functions are shown in the image. One is $f(x) = \frac{1}{2}x + 2$ and the other is $g(x) = -2x + 5$. 

% [Image: constant]

When $n = 1$, the polynomial can be rewritten into

$$
f(x) = mx + b,
$$

where $m$ and $b$ are constants.

The graph of this function is a straight line passing through the points $(0, b)$ and $\left(-\frac{b}{m}, 0\right)$, and the slope of this function is $m$.

\paragraph{Quadratic}

When $n = 2$, the polynomial can be rewritten into

$$
f(x) = ax^2 + bx + c,
$$

where $a$, $b$, $c$ are constants.

The graph of this function is a parabola, like the trajectory of a basketball thrown into the hoop.

% [Image: graph]

If there are questions about the quadratic formula and other basic polynomial concepts, please review the respective chapters in Algebra.

\paragraph{Trigonometric}

Trigonometric functions are also very important because they can connect algebra and geometry. Trigonometric functions are not explained in detail here due to their importance and difficulty.

\paragraph{Exponential and Logarithmic}

Exponential and logarithmic functions have a unique identity when calculating the derivatives, so this is a great time to review those functions. The exponential function is defined as:

$$
f(x) = a^x,
$$

where $a$ is a constant.

The logarithmic function is defined as:

$$
f(x) = \log_b x,
$$

where $b$ is a constant.

A special number will be frequently seen in these functions: Euler''s constant, also known as the base of the natural logarithm. Notated as $e$, it is defined as

$$
e = \sum_{n=0}^{\infty} \frac{1}{n!}.
$$

The curve on the left is an exponential function while the curve on the right is a logarithmic one:

% [Image: log]

\paragraph{Signum}

The Signum (sign) function is defined as:

\begin{itemize}
    \item \(sgn(x) = -1\) if \(x < 0\)
    \item \(sgn(x) = 0\) if \(x = 0\)
    \item \(sgn(x) = 1\) if \(x > 0\)
\end{itemize}

\paragraph{Properties of Functions}

Sometimes, a lot of function manipulations cannot be achieved without some help from basic properties of functions.

\paragraph{Domain and Range}

This concept is discussed above.

\paragraph{Growth}

Although it seems obvious to spot a function increasing or decreasing, without the help of graphing software, we need a mathematical way to spot whether the function is increasing or decreasing, or else our human minds cannot immediately comprehend the huge amount of information.

Assume a function $f(x)$ with inputs $x\_1$, $x\_2$, and that $x\_1 \in [a,b]$, $x\_2 \in [a,b]$, and $x\_2 > x\_1$ at all times.

If for all $x\_1$ and $x\_2$, $f(x\_2) - f(x\_1) > 0$, then

$f(x)$ is increasing in $[a,b]$

If for all $x\_1$ and $x\_2$, $f(x\_2) - f(x\_1) < 0$, then

$f(x)$ is decreasing in $[a,b]$

\paragraph{Example}

In which intervals is $f(x) = \frac{1}{1 - x^2}$ increasing?

Firstly, the domain is important. Because the denominator cannot be 0, the domain for this function is

$x \in (-\infty, -1) \cup (-1, 1) \cup (1, \infty)$

In $(-\infty, -1)$, the growth of the function is:

Let $x\_1, x\_2 \in [-\infty, -1]$ and $x\_2 > x\_1$. Thus,

$$
f(x\_2) - f(x\_1) = \frac{1}{1 - x\_2^2} - \frac{1}{1 - x\_1^2} = \frac{x\_2^2 - x\_1^2}{(1 - x\_2^2)(1 - x\_1^2)}
$$

$\because$ both $x\_1, x\_2 \in [-\infty, -1]$

$\therefore (1 - x\_2^2)(1 - x\_1^2) > 0$

$\because x\_2 > x\_1$ and $x\_1 < x\_2 < 0$

$\therefore x\_2^2 - x\_1^2 < 0$

So,

$$
f(x\_2) - f(x\_1) = \frac{x\_2^2 - x\_1^2}{(1 - x\_2^2)(1 - x\_1^2)} < 0
$$

$f(x)$ is decreasing in $(-\infty, -1)$

In $(-1, 1)$

Let $x\_1, x\_2 \in [-1, 1]$ and $x\_2 > x\_1$. Thus,

$$
f(x\_2) - f(x\_1) = \frac{1}{1 - x\_2^2} - \frac{1}{1 - x\_1^2} = \frac{x\_2^2 - x\_1^2}{(1 - x\_2^2)(1 - x\_1^2)}
$$

$\because$ both $x\_1, x\_2 \in [-1, 1]$

$\therefore (1 - x\_2^2)(1 - x\_1^2) > 0$

However, the sign of $x\_2^2 - x\_1^2$ in $(-1, 1)$ cannot be determined. It can only be determined in $(-1, 0)$ and $(0, 1)$.

In $(-1, 0)$

$\because x\_2 > x\_1$ and $x\_1 < x\_2 < 0$

$\therefore x\_2^2 - x\_1^2 < 0$

In $(0, 1)$

$\because x\_2 > x\_1$ and $0 < x\_1 < x\_2$

$\therefore x\_2^2 - x\_1^2 > 0$

As a result, $f(x)$ is decreasing in $(-1, 0)$ and increasing in $(0, 1)$.

In $(1, \infty)$

Let $x\_1, x\_2 \in [1, \infty]$ and $x\_2 > x\_1$. Thus,

$$
f(x\_2) - f(x\_1) = \frac{1}{1 - x\_2^2} - \frac{1}{1 - x\_1^2} = \frac{x\_2^2 - x\_1^2}{(1 - x\_2^2)(1 - x\_1^2)}
$$

$\because$ both $x\_1, x\_2 \in [1, \infty]$

$\therefore (1 - x\_2^2)(1 - x\_1^2) > 0$

$\because x\_2 > x\_1$ and $0 < x\_1 < x\_2$

$\therefore x\_2^2 - x\_1^2 > 0$

So,

$$
f(x\_2) - f(x\_1) = \frac{x\_2^2 - x\_1^2}{(1 - x\_2^2)(1 - x\_1^2)} > 0
$$

$f(x)$ is increasing in $(1, \infty)$.

Therefore, the intervals in which the function is increasing are $(0, 1) \cup (1, \infty)$.

$\blacksquare$

After learning derivatives, there will be more ways to determine the growth of a function.

\paragraph{Parity}

The properties odd and even are associated with symmetry. While even functions have a symmetry about the $y$-axis, odd functions are symmetric about the origin. In mathematical terms:

A function is even when $f(-x) = f(x)$ 

A function is odd when $f(-x) = -f(x)$

\paragraph{Example}

Prove that $f(x) = \frac{1}{1 - x^2}$ is an even function.

$\because f(-x) = \frac{1}{1 - (-x)^2} = \frac{1}{1 - x^2} = f(x)$

$\therefore f(x)$ is an even function

\paragraph{Manipulating Functions}

\paragraph{Addition, Subtraction, Multiplication, and Division of Functions}

For two real-valued functions, we can add the functions, multiply the functions, raised to a power, etc. For example:

\begin{itemize}
    \item If we add the functions $y = 3x + 2$ and $y = x^2$, we obtain $y = x^2 + 3x + 2$.
    \item If we subtract $y = 3x + 2$ from $y = x^2$, we obtain $y = x^2 - (3x + 2)$. We can also write this as $y = x^2 - 3x - 2$.
    \item If we multiply the function $y = 3x + 2$ and the function $y = x^2$, we obtain $y = (3x + 2)x^2$. We can also write this as $y = 3x^3 + 2x^2$.
    \item If we divide the function $y = 3x + 2$ by the function $y = x^2$, we obtain $y = \frac{3x + 2}{x^2}$.
\end{itemize}

If a math problem wants you to add two functions $f$ and $g$, there are two ways that the problem will likely be worded:

\begin{itemize}
    \item If you are told that $f(x) = 3x + 2$, that $g(x) = x^2$, that $h(x) = f(x) + g(x)$ and asked about $h$, then you are being asked to add two functions. Your answer would be $h(x) = x^2 + 3x + 2$.
    \item If you are told that $f(x) = 3x + 2$, that $g(x) = x^2$ and you are asked about $f + g$, then you are being asked to add two functions. The addition of $f$ and $g$ is called $f + g$. Your answer would be $(f + g)(x) = x^2 + 3x + 2$.
\end{itemize}

Similar statements can be made for subtraction, multiplication, and division.

Let $f(x) = 3x + 2$ and $g(x) = x^2$. Let''s add, subtract, multiply, and divide.

$$(f + g)(x) = f(x) + g(x) = (3x + 2) + (x^2) = x^2 + 3x + 2$$

$$(f - g)(x) = f(x) - g(x) = (3x + 2) - (x^2) = -x^2 + 3x + 2$$

$$(f \times g)(x) = f(x) \times g(x) = (3x + 2) \times (x^2) = 3x^3 + 2x^2$$

$$\left(\frac{f}{g}\right)(x) = \frac{f(x)}{g(x)} = \frac{3x + 2}{x^2}$$

\paragraph{Composition of Functions}

We begin with a fun (and not too complicated) application of composition of functions before we talk about what composition of functions is.

\paragraph{Example: Dropping a Ball}

If we drop a ball from a bridge which is 20 meters above the ground, then the height of our ball above the earth is a function of time. The physicists tell us that if we measure time in seconds and distance in meters, then the formula for height in terms of time is 

$$
h = -4.9t^2 + 20
$$

Suppose we are tracking the ball with a camera and always want the ball to be in the center of our picture. Suppose we have 

$$
\theta = f(h)
$$

The angle will depend upon the height of the ball above the ground and the height above the ground depends upon time. So the angle will depend upon time. This can be written as 

$$
\theta = f(-4.9t^2 + 20)
$$

We replace $h$ with what it is equal to. This is the essence of composition.

\paragraph{Composition of Functions}

Composition of functions is another way to combine functions which is different from addition, subtraction, multiplication, or division.

The value of a function $f$ depends upon the value of another variable $x$; however, that variable could be equal to another function $g$, so its value depends on the value of a third variable. If this is the case, then the first variable is a function $h$ of the third variable; this function ($h$) is called the composition of the other two functions ($f$ and $g$).

\paragraph{Example: Composing Two Functions}

Let 

$$
f(x) = 3x + 2
$$

and 

$$
g(x) = x^2
$$

The composition of $f$ with $g$ is read as either "f composed with g" or "f of g of x."

Let

$$
h(x) = f(g(x))
$$

Then

$$
h(x) = f(g(x)) = f(x^2) = 3(x^2) + 2 = 3x^2 + 2
$$

Sometimes a math problem asks you to compute $(f \circ g)(x)$ when they want you to compute $f(g(x))$.

Here, $h$ is the composition of $f$ and $g$ and we write $h = f \circ g$. Note that composition is not commutative:

$$
f(g(x)) = 3x^2 + 2
$$

and

$$
g(f(x)) = g(3x + 2) = (3x + 2)^2 = 9x^2 + 12x + 4
$$

so 

$$
f(g(x)) \neq g(f(x))
$$

Composition of functions is very common, mainly because functions themselves are common. For instance, squaring and sine are both functions:

$$
\text{square}(x) = x^2
$$

$$
\text{sine}(x) = \sin(x)
$$

Thus, the expression 

$$
\sin^2(x)
$$ 

is a composition of functions:

$$
\sin^2(x) = \text{square}(\sin(x)) = \text{square}(\text{sine}(x))
$$

(Note that this is not the same as 

$$
\text{sine}(\text{square}(x)) = \sin(x^2)
$$

Since the function sine equals $\frac{1}{2}$ if $x = \frac{\pi}{6}$,

$$
\text{square}(\sin(\frac{\pi}{6})) = \text{square}(\frac{1}{2})
$$

Since the function square equals $\frac{1}{4}$ if $x = \frac{\pi}{6}$,

$$
\sin^2(\frac{\pi}{6}) = \text{square}(\sin(\frac{\pi}{6})) = \text{square}(\frac{1}{2}) = \frac{1}{4}
$$

\paragraph{Transformations}

Transformations are a type of function manipulation that are very common. They consist of multiplying, dividing, adding, or subtracting constants to either the input or the output. Multiplying by a constant is called dilation and adding a constant is called translation. Here are a few examples:

\begin{itemize}
    \item $ f(2 \times x) $ — Dilation
    \item $ f(x + 2) $ — Translation
    \item $ 2 \times f(x) $ — Dilation
    \item $ 2 + f(x) $ — Translation
\end{itemize}

\paragraph{Examples of Horizontal and Vertical Translations}

Examples of horizontal and vertical translations can be seen below. 

% [Image: translation]

The red graphs represent functions in their ''original'' state, the solid blue graphs have been translated (shifted) horizontally, and the dashed graphs have been translated vertically.

\paragraph{Examples of Horizontal and Vertical Dilations}

% [Image: dilations]

Dilations are demonstrated in a similar fashion. The function

$$
f(2 \times x)
$$

has had its input doubled. One way to think about this is that now any change in the input will be doubled. If I add one to $x$, I add two to the input of $f$, so it will now change twice as quickly. Thus, this is a horizontal dilation by $\frac{1}{2}$ because the distance to the $y$-axis has been halved. A vertical dilation, such as

$$
2 \times f(x)
$$

is slightly more straightforward. In this case, you double the output of the function. The output represents the distance from the $x$-axis, so in effect, you have made the graph of the function ''taller''. Here are a few basic examples where $a$ is any positive constant:

\paragraph{Transformations}

\begin{itemize}
    \item \textbf{Original graph}: $f(x)$
    \item \textbf{Rotation about origin}: $-f(-x)$
    \item \textbf{Horizontal translation by $a$ units left}: $f(x + a)$
    \item \textbf{Horizontal translation by $a$ units right}: $f(x - a)$
    \item \textbf{Horizontal dilation by a factor of $a$}: $f\left(x \times \frac{1}{a}\right)$
    \item \textbf{Vertical dilation by a factor of $a$}: $a \times f(x)$
    \item \textbf{Vertical translation by $a$ units down}: $f(x) - a$
    \item \textbf{Vertical translation by $a$ units up}: $f(x) + a$
    \item \textbf{Reflection about the $x$-axis}: $-f(x)$
    \item \textbf{Reflection about the $y$-axis}: $f(-x)$
\end{itemize}

\paragraph{Inverse Functions}

We call $g(x)$ the inverse function of $f(x)$ if, for all $x$:

$$
g(f(x)) = f(g(x)) = x
$$

A function $f(x)$ has an inverse function if and only if $f(x)$ is one-to-one. For example, the inverse of $f(x) = x + 2$ is $g(x) = x - 2$. The function $f(x) = \sqrt{1 - x^2}$ has no inverse because it is not injective. Similarly, the inverse functions of trigonometric functions have to undergo transformations and limitations to be considered valid functions.

\paragraph{Notation}

The inverse function of $f$ is denoted as $f^{-1}(x)$. Thus, $f^{-1}(x)$ is defined as the function that follows this rule:

$$
f(f^{-1}(x)) = f^{-1}(f(x)) = x
$$

To determine $f^{-1}(x)$ when given a function $f$, substitute $f^{-1}(x)$ for $x$ and substitute $x$ for $f(x)$. Then solve for $f^{-1}(x)$, provided that it is also a function.

\paragraph{Example}

Given $f(x) = 2x - 7$, find $f^{-1}(x)$.

Substitute $f^{-1}(x)$ for $x$ and $x$ for $f(x)$. Then solve for $f^{-1}(x)$:

$$
f(x) = 2x - 7 \Leftrightarrow x = 2[f^{-1}(x)] - 7 \Leftrightarrow x + 7 = 2[f^{-1}(x)] \Leftrightarrow \frac{x + 7}{2} = f^{-1}(x)
$$

To check your work, confirm that $f^{-1}(f(x)) = x$:

$$
f^{-1}(f(x)) = f^{-1}(2x - 7) = \frac{(2x - 7) + 7}{2} = \frac{2x}{2} = x
$$

If $f$ isn''t one-to-one, then, as we said before, it doesn''t have an inverse. Then this method will fail.

\paragraph{Example}

Given $f(x) = x^2$, find $f^{-1}(x)$.

Substitute $f^{-1}(x)$ for $x$ and $x$ for $f(x)$. Then solve for $f^{-1}(x)$:

$$
f(x) = x^2 \Leftrightarrow x = (f^{-1}(x))^2 \Leftrightarrow f^{-1}(x) = \pm \sqrt{x}
$$

Since there are two possibilities for $f^{-1}(x)$, it''s not a function. Thus $f(x) = x^2$ doesn''t have an inverse. Of course, we could also have found this out from the graph by applying the Horizontal Line Test. It''s useful to have lots of ways to solve a problem, since in a specific case some of them might be very difficult while others might be easy. For example, we might only know an algebraic expression for $f(x)$ but not a graph.

\paragraph{Resources}

\begin{itemize}
    \item \href{https://www.khanacademy.org/math/algebra/x2f8bb11595b61c86:functions/x2f8bb11595b61c86:evaluating-functions/v/what-is-a-function}{What is a Function?, Khan Academy}
    \item \href{https://www.khanacademy.org/math/algebra2/functions_and_graphs}{Function and Their Graphs, Khan Academy}
    \item \href{https://www.cliffsnotes.com/study-guides/calculus/precalculus/functions/relations-vs-functions}{Relations vs. Functions, Cliff''s Notes}
    \item \href{https://courses.lumenlearning.com/ivytech-collegealgebra/chapter/determine-whether-a-relation-represents-a-function/}{Determining if a Relation is a Function, Lumen Learning}
    \item \href{https://courses.lumenlearning.com/waymakercollegealgebra/chapter/identify-functions-using-graphs/}{Identifying Functions Using Graphs, Lumen Learning}
\end{itemize}

\paragraph{Licensing}

Content obtained and/or adapted from:

\begin{itemize}
    \item \href{https://en.wikibooks.org/wiki/Calculus/Functions}{Functions, Wikibooks: Calculus} under a CC BY-SA license
\end{itemize}

\subsection{Learning Objectives}

\begin{enumerate}
    \item Define and differentiate between domain, range, and codomain of a function.
    \item Apply vertical and horizontal line tests to determine if a relation represents a function or if it is injective.
    \item Evaluate functions using their formulas and solve problems involving function mappings, including real-valued and vector-valued functions. Quadratic Functions 1.\textbf{Identify and Describe Quadratic Functions:}
    \item Students will be able to identify the general form of a quadratic function $f(x) = ax^2 + bx + c$, describe its properties (such as vertex, axis of symmetry, and direction of opening), and explain how changes in $a$, $b$, and $c$ affect the graph of the function. 2.\textbf{Solve Quadratic Equations Using the Quadratic Formula:}
    \item Students will be able to apply the quadratic formula $ x = \frac{-b \pm \sqrt{b^2 - 4ac}}{2a} $ to solve quadratic equations and interpret the solutions in the context of real-world problems. 3.\textbf{Graph Quadratic Functions and Analyze Their Properties:}
    \item Students will be able to graph quadratic functions, identify key features such as the vertex, axis of symmetry, and x-intercepts, and analyze how transformations (translations, dilations) affect the graph. Exponential and Logarithmic Functions 1.\textbf{Understand and Apply Exponential Functions:}
    \item Students will be able to define exponential functions $ f(x) = a^x $, understand the impact of different bases and coefficients on the function's growth, and solve problems involving exponential growth and decay. 2.\textbf{Understand and Apply Logarithmic Functions:}
    \item Students will be able to define logarithmic functions $ f(x) = \log_b x $, apply properties of logarithms to simplify expressions, and solve logarithmic equations. 3.\textbf{Utilize Euler’s Constant $ e $ in Exponential and Logarithmic Functions:}
    \item Students will be able to use Euler's constant $ e $ to solve problems involving natural exponential functions $ f(x) = e^x $ and natural logarithms $ f(x) = \ln(x) $, and understand its significance in continuous growth and decay processes. Signum Function 1.\textbf{Evaluate the Signum Function:}
    \item Students will be able to evaluate the signum function $ \text{sgn}(x) $ for given values of $ x $ and interpret its meaning in different contexts. 2.\textbf{Apply the Signum Function in Real-World Scenarios:}
    \item Students will be able to apply the signum function to model and solve problems involving different types of inputs and outputs, including scenarios involving directional or categorical data. Properties of Functions 1.\textbf{Determine the Domain and Range of Functions:}
    \item Students will be able to determine the domain and range of various functions, including rational, polynomial, and piecewise functions, and explain how these properties affect the behavior of the functions. 2.\textbf{Analyze and Describe Function Growth:}
    \item Students will be able to analyze whether a function is increasing or decreasing over specified intervals and use this information to describe the function's behavior in real-world contexts. 3.\textbf{Identify and Apply Function Symmetry:}
    \item Students will be able to determine if a function is even, odd, or neither by analyzing its symmetry about the y-axis and the origin, and apply these properties to solve related problems.
\end{enumerate}

\subsection{Assessment Instrument}

\begin{enumerate}
    \item 1.\textbf{Define and Explain}: Write a brief definition for domain, range, and codomain. Provide an example function and identify its domain, range, and codomain. 2.\textbf{Function Testing}: Given the function $f(x) = 2x + 3$, perform the vertical line test and the horizontal line test. Explain if the function is one-to-one and justify your answer. 3.\textbf{Function Evaluation}: Evaluate the function $f(x) = \frac{2x^2 - 1}{x + 1}$ for $x = 4$. Solve for $x$ if $f(x) = 7$. Quadratic Functions 1.\textbf{Identify Quadratic Functions (Multiple Choice):}
    \item Given a list of functions, select the one that is quadratic.
    \item $a) \quad f(x) = 2x^2 + 3x - 5$
    \item $b) \quad g(x) = \sin(x)$
    \item $c) \quad h(x) = \log(x)$ 2.\textbf{Solve Quadratic Equations (Short Answer):}
    \item Solve the quadratic equation $ 2x^2 - 4x - 6 = 0 $ using the quadratic formula and provide both solutions. 3.\textbf{Graph Quadratic Functions (Graphing):}
    \item Graph the quadratic function $ f(x) = -x^2 + 4x + 1 $, and identify the vertex, axis of symmetry, and x-intercepts. Exponential and Logarithmic Functions 1.\textbf{Evaluate Exponential Functions (Short Answer):}
    \item Calculate $ f(x) $ for the exponential function $ f(x) = 3^x $ at $ x = 2 $ and $ x = -1 $. 2.\textbf{Solve Logarithmic Equations (Short Answer):}
    \item Solve $ \log_2 (x) = 5 $ for $ x $, and explain the steps taken. 3.\textbf{Apply Euler’s Constant (Multiple Choice):}
    \item Identify which function represents continuous growth: $ a) \quad e^x $ $ b) \quad 2^x $ $ c) \quad 3x $ Signum Function 1.\textbf{Evaluate the Signum Function (Short Answer):}
    \item Determine the value of $ \text{sgn}(-7) $, $ \text{sgn}(0) $, and $ \text{sgn}(5) $. 2.\textbf{Application Problem (Word Problem):}
    \item Given a temperature function that changes sign, describe how the signum function can help categorize temperatures as positive, negative, or zero. General Questions 1.\textbf{Determine Domain and Range (Multiple Choice):}
    \item Find the domain of $ f(x) = \frac{1}{x - 3} $:
    \item $a) \quad (-\infty, 3) \cup (3, \infty)$
    \item $b) \quad (-\infty, \infty)$
    \item $c) \quad [3, \infty)$ 2.\textbf{Analyze Function Growth (Short Answer):}
    \item Determine if the function $ f(x) = \frac{1}{x^2} $ is increasing or decreasing in the interval $ (0, \infty) $ and explain why. 3.\textbf{Function Symmetry (True/False):}
    \item Determine if the function $ f(x) = x^3 - 3x $ is odd or even, and justify your answer.
\end{enumerate}

%%===============================================================
\section{Graphs}
\label{sec:graphs}
%%===============================================================

\paragraph{Graphs}

It is sometimes difficult to understand the behavior of a function given only its definition; a visual representation or graph can be very helpful. A graph is a set of points in the Cartesian plane, where each point $(x, y)$ indicates that $f(x) = y$. In other words, a graph uses the position of a point in one direction (the vertical-axis or $y$-axis) to indicate the value of $f$ for a position of the point in the other direction (the horizontal-axis or $x$-axis).

Functions may be graphed by finding the value of $f$ for various $x$ and plotting the points $(x, f(x))$ in a Cartesian plane. For the functions that you will deal with, the parts of the function between the points can generally be approximated by drawing a line or curve between the points. Extending the function beyond the set of points is also possible but becomes increasingly inaccurate.

\paragraph{Linear functions}

Graphing linear functions are easy to understand and do. Because we know that two points can form a line, only two points are needed for us to graph a linear function if those two points are on the function. Oppositely, we can write down the equation of a linear function if we only know two points that are on the function.

The following section mainly talks about different forms of linear function notations so that you can easily identify or graph the function.

\paragraph{Introduction}

Plotting points like this is laborious. Fortunately, many functions'''' graphs fall into general patterns. For a simple case, consider functions of the form

$$
f(x) = 3x + 2
$$

The graph of $f$ is a single line, passing through the point $(0,2)$ with slope 3. Thus, after plotting the point, a straightedge may be used to draw the graph. This type of function is called linear and there are a few different ways to present a function of this type.

\paragraph{Slope}

The slope is the backbone of linear functions because it shows how much the output of a function changes when the input changes. For example, if the slope of a function is 2, then it means when the input of a function increases by 1 unit, the output of the function increases by 2 units. Now, let''''s look at a more mathematical example.

Consider this function: 

$$
f(x) = -5x + 8
$$

What does the number $-5$ mean?

It means that when $x$ increases by 1, $f(x)$ decreases by 5.

Using mathematical terms:

$$
f(x+1) - f(x) = (-5(x+1) + 8) - (-5x + 8) = -5
$$

It is easy to calculate the slope because the slope is like the speed of a vehicle. If we divide the change in distance and the corresponding change in time, we get the speed. Similarly, if we divide the change in $f(x)$ over the corresponding change in $x$, we get the slope. If given two points, $(x\_1, y\_1)$ and $(x\_2, y\_2)$, we may then compute the slope of the line that passes through these two points. Remember, the slope is determined as "rise over run." That is, the slope is the change in $y$-values divided by the change in $x$-values. In symbols:

$$
\text{slope} = \frac{\text{change in } y}{\text{change in } x} = \frac{\Delta y}{\Delta x} = \frac{y\_2 - y\_1}{x\_2 - x\_1}
$$

Slope in a linear function: If two points $(x\_1, y\_1)$ and $(x\_2, y\_2)$ are on a linear function, then the slope of the linear function is

$$
m = \frac{y\_2 - y\_1}{x\_2 - x\_1}
$$

Interestingly, there is a subtle relationship between the slope and the angle between the graph of the function and the positive $x$-axis, $\theta$. The relationship is:

$$
m = \tan \theta
$$

It is an obvious relationship, but it can be ignored relatively easily.

\paragraph{Slope-intercept form}

This is a linear function notated in the slope-intercept form. The slope here is $a$ instead of $m$. When we see a function presented as

$$
y = f(x) = mx + b
$$

we call this presentation the slope-intercept form. This is because, not surprisingly, this way of writing a linear function involves the slope, $m$, and the $y$-intercept, $b$.

\paragraph{Example 1:}

Graph the function 

$$
f(x) = 3x + 2
$$

The slope of the function is 3, and it intercepts the $y$-axis at point $(0,2)$. In order to graph the function, we need another point. Since the slope of the function is 3, then

$$
f(x+1) = f(x) + 3
$$

$$
f(1) = f(0+1) = f(0) + 3 = 5
$$

Knowing that the function goes through points $(0,2)$ and $(1,5)$, the function can be easily graphed.

\paragraph{Example 2:}

Now, consider another unknown linear function that goes through points $(-3, -4)$ and $(0, 5)$. What is the equation for this function?

The slope can be calculated with the formula mentioned above.

$$
m = \frac{y\_2 - y\_1}{x\_2 - x\_1} = \frac{5 - (-4)}{0 - (-3)} = 3
$$

And since the $y$-axis interception is $(0,5)$, we can know that

$$
b = 5
$$

Thus, the equation of this linear function should be

$$
g(x) = 3x + 5
$$

\paragraph{Point-slope form}

If someone walks up to you and gives you one point and a slope, you can draw one line and only one line that goes through that point and has that slope. Said differently, a point and a slope uniquely determine a line. So, if given a point $(x\_0, y\_0)$ and a slope $m$, we present the graph as

$$
y - y\_0 = m(x - x\_0)
$$

We call this presentation the point-slope form. The point-slope and slope-intercept form are essentially the same. In the point-slope form, we can use any point the graph passes through. Whereas, in the slope-intercept form, we use the $y$-intercept, that is the point $(0, b)$. The point-slope form is very important. Although it is not used as frequently as its counterpart, the slope-intercept form, the concept of knowing a point and drawing the line in the direction of the slope will be encountered when we go into vector equations for lines and planes in future chapters.

\paragraph{Example 1:}

If a linear function goes through points $(3, -4)$ and $(1, 5)$, what is the equation for this function?

The slope is:

$$
m = \frac{y\_2 - y\_1}{x\_2 - x\_1} = \frac{5 - (-4)}{1 - 3} = -\frac{9}{2}
$$

Since we know two points, the following answers are all correct:

$$
y + 4 = -\frac{9}{2}(x - 3) \text{ or } y - 5 = -\frac{9}{2}(x - 1)
$$

The two-point form is another form to write the equation for a linear function. It is similar to the point-slope form. Given points $(x\_1, y\_1)$ and $(x\_2, y\_2)$, we have the equation

$$
y - y\_1 = \frac{y\_2 - y\_1}{x\_2 - x\_1}(x - x\_1)
$$

This presentation is in the two-point form. It is essentially the same as the point-slope form except we substitute the expression $\frac{y\_2 - y\_1}{x\_2 - x\_1}$ for $m$. However, this expression is not widely used in mathematics because in most situations, $(x\_1, y\_1)$ and $(x\_2, y\_2)$ are known coordinates. It would be redundant to write down a bulky $\frac{y\_2 - y\_1}{x\_2 - x\_1}$ instead of a simple expression of the slope.

\paragraph{Intercept form}

The intercept form looks like this:

$$
\frac{x}{a} + \frac{y}{b} = 1
$$

By writing the function in the intercept form, we can quickly determine the $x$, $y$-axis intercepts.

$x$-axis intercept: $(a, 0)$

$y$-axis intercept: $(0, b)$

When we discuss planes in 3D space, this form will be quite useful to determine the $x$, $y$, $z$-axis intercepts.

\paragraph{Quadratic functions}

To graph a quadratic function, there is the simple but work-heavy way, and there is the complicated but clever way. The simple way is to substitute the independent variable $x$ with various numbers and calculate the output $f(x)$. After some substitutions, plot those $(x, f(x))$ and connect those points with a curve. The complicated way is to find special points, such as intercepts and the vertex, and plot it out. The following section is a guide to find those special points, which will be useful in later chapters.

\paragraph{Standard form}

Quadratic functions are functions that look like this:

$$
y = f(x) = ax^2 + bx + c
$$

where $a$, $b$, $c$ are constants

The constant $a$ determines the concavity of the function: if $a > 0$, $f(x)$ concaves up; if $a < 0$, $f(x)$ concaves down.

The constant $c$ is the $y$-coordinate of the $y$-axis interception. In other words, this function goes through point $(0, c)$.

\paragraph{Vertex Form}

The vertex form has its advantages over the standard form. While the standard form can determine the concavity and the $y$-axis interception, the vertex form can, as the name suggests, determine the vertex of the function. The vertex of a quadratic function is the highest/lowest point on the graph of a function, depending on the concavity. If $a > 0$, the vertex is the lowest point on the graph; if $a < 0$, the vertex is the highest point on the graph.

The vertex form looks like this:

$$
y = f(x) = a(x - h)^2 + k
$$

where $a$, $h$, $k$ are constants.

The vertex of this function is $(h, k)$ because when $x = h$, $f(h) = k$. If $a > 0$, $f(h) = k$ is the absolute minimum value that the function can achieve. If $a < 0$, $f(h) = k$ is the absolute maximum value that the function can achieve. Any standard form can be converted into the vertex form. The vertex form with constants $a$, $b$, $c$ looks like this:

$$
y = f(x) = a \left(x + \frac{b}{2a}\right)^2 + \left(c - \frac{b^2}{4a}\right)
$$

where $a$, $b$, $c$ are constants in the standard form.

\paragraph{Factored Form}

The factored form can determine the $x$-axis intercepts because the factored form looks like this:

$$
y = f(x) = a(x - x\_1)(x - x\_2)
$$

where $x\_1$, $x\_2$ are constants and are solutions for the equation $a(x - x\_1)(x - x\_2) = 0$.

Thus, it can be determined that the function passes through points $(x\_1, 0)$ and $(x\_2, 0)$.

However, only certain functions can be written in this form. If the quadratic function does not have $x$-axis intercepts, it is impossible to write it in the factored form.

\paragraph{Example 1}

What is the vertex of this function? 

$$
f(x) = x^2 + 2x + 3
$$

The equation can be transformed into the vertex form very easily:

$$
x^2 + 2x + 3 = (x^2 + 2x + 1) + 2 = (x + 1)^2 + 2
$$

Thus, the vertex is $(-1, 2)$.

\paragraph{Example 2}

This is an image of a quadratic function $y = ax^2 + bx + c$ with key values. 

% [Image: quadratic]

The image is a quadratic function. Describe the meaning of the colored texts, which are important properties of a quadratic function.

$$
y = ax^2 + bx + c
$$

is the equation for the quadratic function. In this case, $a > 0$, $c < 0$. Since there are two $x$-axis intercepts, we can find that $\Delta = b^2 - 4ac > 0$.

Points 

$$
\left(\frac{-b - \sqrt{b^2 - 4ac}}{2a}, 0\right) \text{ and } \left(\frac{-b + \sqrt{b^2 - 4ac}}{2a}, 0\right)
$$

These are the coordinates for the two $x$-axis intercepts.

Knowing the coordinates, the function can be written in its factored form:

$$
y = a \left(x - \frac{-b - \sqrt{b^2 - 4ac}}{2a}\right) \left(x - \frac{-b + \sqrt{b^2 - 4ac}}{2a}\right)
$$

Point 

$$
\left(-\frac{b}{2a}, -\frac{b^2 - 4ac}{4a}\right)
$$

is the vertex for the quadratic function. Because $a > 0$, the vertex is the lowest point on the graph. Since the vertex is known, we can write the function in the vertex form:

$$
y = a \left(x + \frac{b}{2a}\right)^2 - \frac{b^2 - 4ac}{4a}
$$

Although this does not look like the equation we''''ve just discussed earlier, note that 

$$
\- \frac{b^2 - 4ac}{4a} = c - \frac{b^2}{4a}
$$

Line 

$$
x = -\frac{b}{2a}
$$

The graph of the function is symmetric about this line. In other words,

$$
f\left(-\frac{b}{2a} + x\right) = f\left(-\frac{b}{2a} - x\right)
$$

If you can skillfully and quickly determine those special points, graphing quadratic functions will be less torturing.

\paragraph{Exponential and Logarithmic Functions}

Exponential and logarithmic functions are inverse functions of each other. Take the exponential function 

$$
f(x) = a^x
$$

for example. The inverse function of $f(x)$, $f^{-1}(x)$, is

$$
x = a^{f^{-1}(x)}
$$

$$
\Leftrightarrow f^{-1}(x) = \log_a x
$$

which is a logarithmic function.

Since geometrically, the graph of the inverse function is flipping the graph of the original function over the line $y = x$, we only need to know how to graph one of those functions.

\paragraph{Resources}

\paragraph{Basic Information}

\begin{itemize}
    \item \href{https://www.mathsisfun.com/data/cartesian-coordinates.html}{Introduction to Cartesian Coordinates and Graphs, mathisfun.com}
    \item \href{https://www.khanacademy.org/math/algebra/x2f8bb11595b61c86:foundation-algebra/x2f8bb11595b61c86:algebra-overview-history/v/descartes-and-cartesian-coordinates}{Introduction to the Coordinate Plane, Khan Academy}
    \item \href{https://courses.lumenlearning.com/ivytech-collegealgebra/chapter/plotting-ordered-pairs-in-the-cartesian-coordinate-system/}{Plotting Points on a Graph, Lumen Learning}
\end{itemize}

\paragraph{Graphing Linear Functions}

\begin{itemize}
    \item \href{https://courses.lumenlearning.com/waymakercollegealgebra/chapter/graph-linear-functions/}{Graphing Linear Functions, Lumen Learning}
    \item \href{https://www.khanacademy.org/math/algebra/x2f8bb11595b61c86:linear-equations-graphs}{Linear Equations and Graphs, Khan Academy}
    \item \href{https://www.palmbeachstate.edu/prepmathlw/Documents/graphing\%20linear\%20equations.pdf}{Graphing Linear Equations, Palm Beach State College}
    \item \href{https://www.youtube.com/watch?v=lkRyxnqOb8M}{Graphing a Line Using a Point and Slope, patrickJMT}
    \item \href{https://www.youtube.com/watch?v=mxBoni8N70Y}{Graphing a Line with X and Y Intercepts, patrickJMT}
\end{itemize}

\paragraph{Licensing}

Content obtained and/or adapted from:

\begin{itemize}
    \item \href{https://en.wikibooks.org/wiki/Calculus/Graphing_functions}{Graphing Functions, Wikibooks: Calculus} under a CC BY-SA license
\end{itemize}

\subsection{Learning Objectives}

\begin{enumerate}
    \item SLO 1: Graphing Linear Functions \textbf{Objective:} Students will be able to graph a linear function given its equation in slope-intercept form. \textbf{Criteria:}
    \item Identify the slope $m$ and $y$-intercept $b$ from the equation of the linear function $y = mx + b$.
    \item Plot the $y$-intercept on the Cartesian plane.
    \item Use the slope to determine a second point on the line.
    \item Draw the line through the two points. SLO 2: Calculating and Interpreting Slope \textbf{Objective:} Students will be able to calculate the slope of a line given two points on the line and interpret its meaning. \textbf{Criteria:}
    \item Use the formula $m = \frac{y\_2 - y\_1}{x\_2 - x\_1}$ to compute the slope of the line through points $(x\_1, y\_1)$ and $(x\_2, y\_2)$.
    \item Explain the meaning of the slope in the context of the linear function, including how it represents the rate of change.
    \item Relate the slope to the angle $\theta$ with the $x$-axis using the relationship $m = \tan \theta$. SLO 3: Converting Between Linear Forms \textbf{Objective:} Students will be able to convert the equation of a linear function from slope-intercept form to point-slope form and vice versa. \textbf{Criteria:}
    \item Given a linear function in slope-intercept form $y = mx + b$, rewrite it in point-slope form $y - y\_0 = m(x - x\_0)$ using a known point $(x\_0, y\_0)$ on the line.
    \item Given a linear function in point-slope form, convert it to slope-intercept form.
    \item Demonstrate understanding of how each form provides different information about the line, such as slope and intercepts.
\end{enumerate}

\subsection{Assessment Instrument}

\begin{enumerate}
    \item Assessment Instrument: Graphs and Functions Section 1: Linear Functions 1.1. Graph of $y = 2x$
    \item \textbf{Graphing Task:} Plot the graph of the function $y = 2x$ on a Cartesian plane. Indicate the slope and the $y$-intercept on the graph.
    \item \textbf{Question:} Describe the behavior of the function $y = 2x$ based on its graph. What is the slope of the line, and where does it intersect the $y$-axis? 1.2. Slope Calculation
    \item \textbf{Given Function:} $f(x) = -5x + 8$
    \item \textbf{Question:} What is the slope of this function? Explain what the slope indicates about the function's behavior.
    \item \textbf{Points Provided:} $(2, 3)$ and $(5, 7)$
    \item \textbf{Question:} Calculate the slope of the line passing through these points. Use the slope formula and show your work. 1.3. Slope-Intercept Form
    \item \textbf{Example 1:}
    \item \textbf{Function:} $f(x) = 3x + 2$
    \item \textbf{Question:} Graph the function and identify the slope and $y$-intercept.
    \item \textbf{Example 2:}
    \item \textbf{Points Provided:} $(-3, -4)$ and $(0, 5)$
    \item \textbf{Question:} Determine the equation of the linear function that passes through these points. 1.4. Point-Slope Form
    \item \textbf{Example 1:}
    \item \textbf{Points Provided:} $(3, -4)$ and $(1, 5)$
    \item \textbf{Question:} Write the equation of the line in point-slope form using either of the provided points.
    \item \textbf{Additional Form:}
    \item \textbf{Question:} Write the equation of the line in two-point form using the given points. 1.5. Intercept Form
    \item \textbf{Question:} Convert the linear function $2x + 3y = 6$ into intercept form and identify the $x$-axis and $y$-axis intercepts. Section 2: Quadratic Functions 2.1. Standard Form
    \item \textbf{Function Provided:} $y = 2x^2 - 4x + 1$
    \item \textbf{Question:} Identify the $y$-intercept of the function. What does the coefficient of $x^2$ tell you about the concavity of the parabola? 2.2. Vertex Form
    \item \textbf{Example 1:}
    \item \textbf{Function Provided:} $f(x) = x^2 + 2x + 3$
    \item \textbf{Question:} Convert the function into vertex form and identify the vertex.
    \item \textbf{Example 2:}
    \item \textbf{Graph Analysis:} Using the graph provided in the image, identify the $x$-intercepts and vertex of the quadratic function. Convert the function into vertex form. 2.3. Factored Form
    \item \textbf{Example:}
    \item \textbf{Function Provided:} $f(x) = (x - 2)(x + 3)$
    \item \textbf{Question:} Expand the factored form and identify the $x$-intercepts of the quadratic function. Section 3: Exponential and Logarithmic Functions 3.1. Exponential Functions
    \item \textbf{Function Provided:} $f(x) = 2^x$
    \item \textbf{Question:} Sketch the graph of $f(x) = 2^x$. What are the key features of the graph? 3.2. Logarithmic Functions
    \item \textbf{Function Provided:} $g(x) = \log_2 x$
    \item \textbf{Question:} Sketch the graph of $g(x) = \log_2 x$. How does it relate to the graph of the exponential function $f(x) = 2^x$?
    \item \textbf{Question:} Verify that the logarithmic function $g(x)$ is the inverse of the exponential function $f(x)$. Explain the geometric relationship between the two graphs. ---
\end{enumerate}

%%===============================================================
\section{Linear Functions}
\label{sec:linear-functions}
%%===============================================================

\paragraph{Linear Functions}

\paragraph{What are Linear Equations?}

In the functions section, we discuss how a function is like a box that takes an independent input value and uses a rule defined mathematically to create a unique output value. The value for the output is dependent on the value that is put in the box. We call the values that go into the box the independent values or the domain. The values coming out of the box are called the dependent values or the range.

Unless specified otherwise, on Cartesian graphs, the domain is the real numbers. In the Cartesian Plane section, we saw how running different values through a function identifies points on the Cartesian plane by picking the first (x) value of the point from the domain and the second (y) value from the range. To restate: by convention, the two variables for a function on the Cartesian Plane are $x$ for the domain (independent variable) and $y$ for the range (dependent variable). The variable $y$ is the same as writing $f(x)$. Mathematicians recognize this equivalence but generally prefer to write $y$ because it''''''''s shorter. Since a function definition has an input and an output, it must also contain an equal sign. The section on solving equations showed the various operations we can perform on both sides of an equal sign while maintaining the notion of equivalence. In this section, we plug different values into the independent variable and solve to find the associated dependent variable. For instance, if we start with the equation:

$$
y = x
$$

We can add $-x$ to both sides to get the equation

$$
y - x = x - x
$$

which simplifies to

$$
y - x = 0
$$

Or we can add $-y$ to both sides to get the equation

$$
y - y = x - y
$$

which simplifies to

$$
0 = x - y
$$

Since

$$
y - x = 0 \equiv 0 = x - y
$$

then:

$$
y - x = 0 = x - y
$$

Using the transitive property:

$$
y - x = x - y
$$

Adding $x + y$ to both sides gives us

$$
(y - x) + (y + x) = (x - y) + (y + x)
$$

Using the associative property, we change this to

$$
(y - y) + (x + x) = (x - x) + (y + y)
$$

which simplifies to

$$
2x = 2y
$$

And dividing both sides by 2:

$$
\frac{2x}{2} = \frac{2y}{2}
$$

To simply reverse the order and show

$$
x = y
$$

We have not really proved anything mathematically above, but these operations allow us to manipulate equations to get the dependent variable by itself on one side of the equals sign. We can then plug numbers into the independent variable to discover the function values for those numbers. We can plot these values as points on the Cartesian plane to get a feel for what the function would look like if we could see all the points defined by the function at once.

\paragraph{Horizontal Linear Equations}

Equations of the form $y = C\_2$ are linear functions of the general form $y = mx + b$ where the slope $m = 0$ and the constant $C\_2$ is the y-intercept $b$ (in the general form). The graph of this zero-slope function is a straight horizontal line, intercepting the y-axis at $C\_2$, including zero and extending infinitely in the positive and negative directions for all $R$ values of $x$.

% [Image: Horizontal line y=3]
$$
y = c, \text{ where } c = 3
$$

The domain for such functions is $R$ covering all real numbers (unless otherwise specified), but the range is just the set $\{ c \}$. The equation $y = 0$ is the x-axis.

\paragraph{Vertical Linear Equations}

Equation $x = C\_1$ where $x$ is one single value $C\_1$ and $y$, being unrestricted, covers every $R$ number. The graph of $x = C\_1$ is a straight vertical line where $x = C\_1$, covering all positive, negative, and zero values of $y$.

% [Image: Vertical Line, x=3]
$$
x = c, \text{ where } c = 3
$$

Its domain is set $\{ C\_1 \}$ and range is set $R$ (unless otherwise specified). $x = C\_1$ is technically not a function (there is more than one value of $y$ for each value of $x$), but it''''''''s a relation. Vertical lines have no slope ($m = \text{divide by zero, undefined, plus and minus infinity}$). These are the only types of linear equations of the general form shown previously which are not linear functions. The equation $x = 0$ is exactly the y-axis. Lines with steepness approaching vertical have very large-magnitude slopes but are still functions.

\paragraph{General Linear Equations}

We are going to start by looking at simple functions called linear equations. When none of the instances of $x$ and $y$ in the algebraic expression defining the function rule have exponents, all the instances of $x$ and $y$ can be combined into just two occurrences. The graph of the expression can be represented as a straight line. The equation that expresses the function is considered a linear equation with two variables. The following equation is a simple example of such a linear equation:

$$
y - x = 2
$$

Since $y$ is the dependent variable, it is standing in for the function. We can re-write the expression as $f(x) - x = 2$. If we add an $x$ to both sides, the equality property holds and we get the expression $f(x) - x + x = x + 2$. Simplifying, we get $f(x) = x + 2$. In the following table, we''''''''ll pick 3 values for $x$ and then calculate the dependent ($y$) values from $f(x)$.

When $x = -1$, $y = 1$  Coordinate (-1,1)
When $x = 0$, $y = 2$  Coordinate (0,2)
When $x = 1$, $y = 3$  Coordinate (1,3)

where $x$ and $y$ are variables to be plotted in a two-dimensional Cartesian coordinate graph as shown here:

% [Image: Graph]
% [Image: Graph2]

This function is equivalent to the previous example of a linear equation, $y - x = 2$. The arrows at each end of the line indicate that the line extends infinitely in both directions. All linear functions of a single input variable have or can be algebraically arranged to have the general form:

$$
y = f(x) = mx + b
$$

where $x$ and $y$ are variables, $f(x)$ is the function of $x$, $m$ is a constant called the slope of the line, and $b$ is a constant which is the ordinate of the y-intercept (i.e., the value of $y$ where the function line crosses the y-axis). The slope indicates the steepness of the line. In the previous example where $y = x + 2$, the slope $m = 1$ and the y-intercept ordinate $b = 2$. The $y = mx + b$ form of a linear function is called the slope-intercept form.

Unless a domain for $x$ is otherwise stated, the domain for linear functions will be assumed to be all real numbers and so the lines in graphs of all linear functions extend infinitely in both directions. Also, in linear functions with all real number domains, the range of a linear function will cover the entire set of real numbers for $y$, unless the slope $m = 0$ and the function equals a constant. In such cases, the range is simply the constant.

Conversely, if a function has the general form $y = mx + b$ or if it can be arranged to have that form, the function is linear. A linear equation with two variables has or can be algebraically rearranged to have the general form:

$$
Ax + By = C
$$

where $x$ and $y$ represent the linear variables, and the letters $A$, $B$, and $C$ can represent any real constants, either positive or negative. Conversely, an equation with two variables $x$ and $y$ having that general form, or being able to be arranged in that form, would be linear as long as $A$ and $B$ are not both equal to 0. In the preceding equation, capital letters are used to avoid confusion with other constants in this chapter and for consistency.

If one divides the preceding equation by $B$ (when $B \neq 0$) and solves for $y$, the following form can be obtained:

$$
y = \left(-\frac{A}{B}\right)x + \frac{C}{B}
$$

If one equates $-\frac{A}{B}$ to the slope $m$ and $\frac{C}{B}$ to the y-intercept ordinate $b$, it can be seen that the general form for a linear equation and the slope-intercept form for a linear function are practically interconvertible except for the fact that, in a linear function, the B constant in the linear equation form cannot equal 0.

\paragraph{X and Y Axis Intercepts}

An axis intercept point is a point where the graph of a function, relation, or equation intersects the X or Y axes. This section is about finding out how a particular set of functions—linear functions—cross the axes.

We know that the domains of most lines are infinite because they are defined at every value of $X$. The exception is lines that are defined as $X = c$ where $c$ is a number we choose when we write the function. By definition, this line is only defined for one value of $X$. Since the domain maps onto more than one value for the range, this is actually a relationship and not a function. We''''''''ve seen that the graph of these relationships is a vertical line that passes through the point $(c, Y)$.

Lines with the equation $X = C$ intersect the X axis once and the Y axis 0 or infinite times.

We can also restrict the range of a function by simply writing $Y = c$ where $c$ is again any number we choose. The graph of this line is a horizontal line that passes through the point $(X, c)$. When $Y = 0$, this line is the same as the X axis. When $c \neq 0$, then the line can never intercept the X axis.

Lines with the equation $Y = C$ intersect the Y axis once and the X axis 0 or infinite times.

When looking at Cartesian graphs and linear equations, we run into a mathematical axiom: "Two points determine a line." We will see how this axiom affects the slope-intercept definition of a line $y = f(x) = mx + b$ in the next section. When two lines intersect, they intersect at a point. If a line is not horizontal or perpendicular, it will have to intersect the X and Y axes once, but only once.

In this page, we are going to accept the statement: "At most one line can be drawn through any point not on a given line parallel to the given line in a plane."

We''''''''ve seen that for the equation $Y = mX + b$, the Y intercept will always be at $b$ because that is where $X = 0$.

Using Algebra, we can subtract $b$ from both sides:

$$
Y - b = mX
$$

and multiply by $\frac{1}{m}$:

$$
\frac{Y - b}{m} = X
$$

we can see that the X intercept is going to be $-\frac{b}{m}$.

An axis intercept may simply refer to the number value on the axis where the intersection occurs. For brevity, we may say the line has an X intercept of 1 and a Y intercept of 2. After graphing just a few lines, you will be able to tell that this line points down and runs through quadrants II, I, and IV. With a little more practice, you will be able to know that the equation for the line is $Y = -2x + 2$. We will see that by specifying the two points, we are actually implying the slope of the line. There is an exception to this rule. If we say a line crosses the axes at 0, we know that the line will pass through 2 quadrants instead of 3, but we won''''''''t know which quadrants or how steep the line is. When we look at slope in the next section, we will see why the equations above specify a point and a slope.

When you are trying to graph a linear equation, finding the axis intercepts is often the easiest way to go about doing it. To find the x-intercept, set $y = 0$ and solve for $x$. To find the y-intercept, set $x = 0$ and solve for $y$. For most examples, the intercepts are different points, and a line can be drawn through the two intercepts. If both intercepts are $(0,0)$, then another point must be determined to graph the line. If the equation is in the form $x = c$ or $y = c$, the horizontal or vertical lines are very simple to plot.

\paragraph{Slope}

\paragraph{Algebra/Slope}

Slope is the change in the vertical distance of a line on a coordinate plane over the change in horizontal distance. In other words, it is the “rise” over the “run” or the steepness of a line. Slope is usually represented by the symbol $m$, like in the equation $y = mx + b$, where $m$ is the coefficient of $x$ representing the slope of the line.

Slope is computed by measuring the change in vertical distance divided by the change in horizontal distance:

$$
m = \frac{\Delta y}{\Delta x} = \frac{\text{rise}}{\text{run}}
$$

The Greek uppercase letter $\Delta$ represents change, in this case, change in the $y$-coordinates divided by the change in $x$-coordinates.

\paragraph{Positive Slope/ Negative Slope}

If a line goes up from left to right, then the slope has to be positive. For example, a slope of $\frac{3}{4}$ would have a “rise” of 3 (or go up 3) and a “run” of 4 (or go right 4). Both numbers in the slope are either negative or positive in order to have a positive slope.

If a line goes down from left to right, then the slope has to be negative. For example, a slope of $-\frac{3}{4}$ would have a “rise” of -3 (or go down 3) and a “run” of 4 (or go right 4). Only one number in the slope can be negative for a line to have a negative slope.

\paragraph{Other Types of Slope}

There are two special circumstances: no slope and a slope of zero. A horizontal line has a slope of 0 and a vertical line has an undefined slope.

Horizontal lines have the form:

$$
y = a
$$

where $a$ is a constant, i.e., $a \in \mathbb{R}$.

Vertical lines have the form:

$$
x = a
$$

where $a$ is a constant, i.e., $a \in \mathbb{R}$.

\paragraph{Determining Slope}

To determine the slope, you need some information. This can include two (or more) coordinates, a parallel slope and a coordinate, a perpendicular slope and a coordinate, or the y-intercept and slope.

For completely horizontal lines, the difference in $y$-coordinates between any two points is 0, so the slope $m = 0$, indicating no steepness in the line at all. If the line extends between right-upper (+,+) and left-lower (-,-) directions, then the slope is positive. As the slope increases, the line becomes steeper until the line is almost vertical when the slope is very large. When the slope $m = 1$, the line is diagonal with an angle halfway between the $x$ and $y$ axes. If the line extends between left-upper (-,+) and right-lower (+,-) directions, then the slope is negative. As the slope changes from 0 to very negative numbers, the steepness in the opposite direction increases. Compare the slope ($m$) values in the following graph of functions: $y = 1$ (where $m = 0$), $y = \frac{1}{2}x + 1$, $y = x + 1$, $y = 2x$, $y = -\frac{1}{2}x + 1$, $y = -x + 1$, and $y = -2x + 1$. For all two-variable linear equations that can be converted to linear functions, the same calculation applies to slopes for those lines.

% [Image: Various Linear Slopes]

For the most part, finding the slope when given information is a simple matter. Simply take the slope equation $y = mx + b$ and replace the variable with whatever information you know, and solve.

\paragraph{Two Coordinates}

To find the slope with two coordinates, you must first find the slope using the standard equation:

$$
\frac{y\_2 - y\_1}{x\_2 - x\_1}
$$

Put that into the equation as $m$, and replace $x$ and $y$ with $x$ and $y$ from one of the coordinates. Solve for $b$. Put that into the equation, and you''''''''re done.

\textbf{Example:} $(1, 4)$ and $(4, 8)$

$$
m = \frac{y\_2 - y\_1}{x\_2 - x\_1}
$$

$$
m = \frac{8 - 4}{4 - 1}
$$

$$
m = \frac{4}{3}
$$

Plug that into the equation:

$$
y = \frac{4}{3}x + b
$$

To find $b$:

$$
4 = \frac{4}{3}(1) + b
$$

$$
4 = \frac{4}{3} + b
$$

$$
4 - \frac{4}{3} = b
$$

$$
b = \frac{8}{3}
$$

Put that in the equation:

$$
y = \frac{4}{3}x + \frac{8}{3}
$$

\paragraph{Parallel Lines}

If you need to find the slope of a line (let''''''''s say AB) which is parallel to another line (let''''''''s say XY), then using the coordinates of line XY, you can find the coordinates of the slope of line AB. The slope of a line parallel to another line is equal, i.e., Slope of AB = Slope of XY.

\textbf{Example:}

AB and XY are parallel lines. Find the slope of AB.

Let line XY have the coordinates:

$X(2, 4) = (x\_1, y\_1)$ and $Y(3, 6) = (x\_2, y\_2)$

Slope of XY:

$$
m = \frac{y\_2 - y\_1}{x\_2 - x\_1} = \frac{6 - 4}{3 - 2} = 2
$$

Slope of AB = Slope of XY, so Slope of AB = 2

\paragraph{Resources}

\begin{itemize}
    \item \href{https://courses.lumenlearning.com/beginalgebra/chapter/introduction-to-solving-linear-equations/}{Solving Linear Equations, Lumen Learning}
    \item \href{https://openstax.org/books/intermediate-algebra-2e/pages/2-1-use-a-general-strategy-to-solve-linear-equations}{General Strategy to Solving Linear Equations, OpenStax}
    \item \href{https://openstax.org/books/intermediate-algebra-2e/pages/2-3-solve-a-formula-for-a-specific-variable}{Solving Linear Equations for a Specific Variable, OpenStax}
    \item \href{https://www.youtube.com/watch?v=bsGyJ0GeRy8}{Solving Linear Equations: Example 1.} Produced by TA Catherine Sporer, UTSA
    \item \href{https://www.youtube.com/watch?v=nGZcfHQEdqo}{Solving Linear Equations: Example 2.} Produced by TA Catherine Sporer, UTSA
\end{itemize}

\paragraph{Licensing}

Content obtained and/or adapted from:

\begin{itemize}
    \item \href{https://en.wikibooks.org/wiki/Algebra/Linear_Equations_and_Functions}{Linear Equations and Functions, Wikibooks: Algebra} under a CC BY-SA license
    \item \href{https://en.wikibooks.org/wiki/Algebra/Intercepts}{Intercepts, Wikibooks: Algebra} under a CC BY-SA license
    \item \href{https://en.wikibooks.org/wiki/Algebra/Slope}{Slope, Wikibooks: Algebra} under a CC BY-SA license
\end{itemize}

\subsection{Learning Objectives}

\begin{enumerate}
    \item \textbf{Identify and Graph}: Students will identify and graph horizontal and vertical linear equations on the Cartesian plane.
    \item \textbf{Determine Slope}: Students will calculate the slope of a linear equation given two points and interpret the slope in terms of the line’s steepness.
    \item \textbf{Convert Forms}: Students will convert between different forms of linear equations (e.g., general form to slope-intercept form) and use this to find intercepts.
\end{enumerate}

\subsection{Assessment Instrument}

\begin{enumerate}
    \item SLO 1: Identify and Graph \textbf{Instructions:} For each of the following equations, identify whether it is a horizontal or vertical line, and then graph it on the Cartesian plane.
    \item $y = 4$
    \item $x = -2$ \textbf{Questions:}
    \item \textbf{Equation}: $y = 4$
    \item \textbf{Type}: Horizontal or Vertical?
    \item \textbf{Graph}: (Draw the graph in the space below or use a graphing tool)
    \item \textbf{Equation}: $x = -2$
    \item \textbf{Type}: Horizontal or Vertical?
    \item \textbf{Graph}: (Draw the graph in the space below or use a graphing tool) --- SLO 2: Determine Slope \textbf{Instructions:} Calculate the slope of the line passing through the given points. Interpret the slope in terms of the line’s steepness.
    \item Points: $(1, 3)$ and $(4, 11)$
    \item Points: $(-2, 5)$ and $(3, -1)$ \textbf{Questions:}
    \item \textbf{Points}: $(1, 3)$ and $(4, 11)$
    \item \textbf{Calculate the Slope}: $$ m = \frac{y\_2 - y\_1}{x\_2 - x\_1} $$
    \item \textbf{Interpret the Slope}:
    \item \textbf{Points}: $(-2, 5)$ and $(3, -1)$
    \item \textbf{Calculate the Slope}: $$ m = \frac{y\_2 - y\_1}{x\_2 - x\_1} $$
    \item \textbf{Interpret the Slope}: --- SLO 3: Convert Forms \textbf{Instructions:} Convert the given linear equations from general form to slope-intercept form. Use this to find the $x$-intercept and $y$-intercept.
    \item $2x - 3y = 6$
    \item $5x + 4y = -8$ \textbf{Questions:}
    \item \textbf{Equation}: $2x - 3y = 6$
    \item \textbf{Convert to Slope-Intercept Form}: $$ y = mx + b $$
    \item \textbf{Find $x$-intercept}:
    \item \textbf{Find $y$-intercept}:
    \item \textbf{Equation}: $5x + 4y = -8$
    \item \textbf{Convert to Slope-Intercept Form}: $$ y = mx + b $$
    \item \textbf{Find $x$-intercept}:
    \item \textbf{Find $y$-intercept}: ---
\end{enumerate}

%%===============================================================
\section{Slope}
\label{sec:slope}
%%===============================================================

\paragraph{Slope}

\paragraph{Slope Between Two Points}

Given two points $(x\_1, y\_1)$ and $(x\_2, y\_2)$, the slope between these two points is

$$
m = \frac{y\_2 - y\_1}{x\_2 - x\_1}
$$

That is, the slope between two points is the difference between the $y$-coordinates of the points, divided by the difference between the $x$-coordinates of the points. 

For example, the slope between the two points $(1, 3)$ and $(5, 6)$ is

$$
\frac{6 - 3}{5 - 1} = \frac{3}{4}
$$

The slope between $(-1, -1)$ and $(15, -21)$ is

$$
\frac{-21 - (-1)}{15 - (-1)} = \frac{-21 + 1}{15 + 1} = \frac{-20}{16} = -\frac{5}{4}
$$

---

\paragraph{Point-Slope Form of a Line}

The equation for a line with a slope of $m$ that goes through some point $(x\_1, y\_1)$, in point-slope form, is

$$
y - y\_1 = m (x - x\_1)
$$

For example, the equation of a line with a slope of 3 that goes through the point $(1, 4)$ is

$$
y - 4 = 3 (x - 1)
$$

The equation of a line with a slope of $-\frac{1}{2}$ that goes through the point $(-7, -7)$ is

$$
y + 7 = -\frac{1}{2} (x + 7)
$$

---

\paragraph{Slope-Intercept Form of a Line}

Another form of an equation of a line is slope-intercept form. The equation of a line with a $y$-intercept of $b$ (that is, a line that intersects the $y$-axis at the point $(0, b)$) and a slope of $m$ is

$$
y = mx + b
$$

For example, the equation of a line with a $y$-intercept of 5 and slope of 6 is

$$
y = 6x + 5
$$

Note that this equation is equivalent to point-slope form. A line with a $y$-intercept of 5 goes through the point $(0, 5)$, so the point-slope form of this same line is

$$
y - 5 = 6 (x - 0) = 6x
$$

By adding 5 to each side of the equation, we get the slope-intercept form of the line.

\paragraph{Resources}

\begin{itemize}
    \item \href{https://tasks.illustrativemathematics.org/content-standards/tasks/1537}{Slope Between Points, Illustrative Mathematics}
    \item \href{https://www.khanacademy.org/math/algebra/x2f8bb11595b61c86:forms-of-linear-equations/x2f8bb11595b61c86:point-slope-form/v/idea-behind-point-slope-form}{Point-Slope Form, Khan Academy}
    \item \href{https://www.khanacademy.org/math/algebra/x2f8bb11595b61c86:forms-of-linear-equations/x2f8bb11595b61c86:intro-to-slope-intercept-form/v/slope-intercept-form}{Slope-Intercept Form, Khan Academy}
    \item \href{https://www.khanacademy.org/math/algebra/x2f8bb11595b61c86:forms-of-linear-equations/x2f8bb11595b61c86:summary-forms-of-two-variable-linear-equations/a/forms-of-linear-equations-review}{Forms of Linear Equations, Khan Academy}
\end{itemize}

\subsection{Learning Objectives}

\begin{enumerate}
    \item 1. Calculate Slope Between Two Points Students will be able to calculate the slope $m$ between two points $(x\_1, y\_1)$ and $(x\_2, y\_2)$ on the Cartesian plane using the formula: $$ m = \frac{y\_2 - y\_1}{x\_2 - x\_1} $$ 2. Write the Equation of a Line in Point-Slope Form Students will be able to write the equation of a line in point-slope form given a slope $m$ and a point $(x\_1, y\_1)$ on the line, using the formula: $$ y - y\_1 = m (x - x\_1) $$ 3. Convert Between Point-Slope Form and Slope-Intercept Form Students will be able to convert the equation of a line between point-slope form and slope-intercept form. Specifically, they should be able to:
    \item Convert from point-slope form: $$ y - y\_1 = m (x - x\_1) $$ to slope-intercept form: $$ y = mx + b $$
    \item Identify and interpret the $y$-intercept $b$ from the slope-intercept form.
\end{enumerate}

\subsection{Assessment Instrument}

\begin{enumerate}
    \item 1. Calculate Slope Between Two Points \textbf{Objective:} Students will calculate the slope $m$ between two points $(x\_1, y\_1)$ and $(x\_2, y\_2)$. \textbf{Questions:}
    \item Given the points $(2, 3)$ and $(6, 7)$, calculate the slope $m$.
    \item Find the slope of the line passing through the points $(-3, 5)$ and $(4, -2)$. 2. Write the Equation of a Line in Point-Slope Form \textbf{Objective:} Students will write the equation of a line in point-slope form given a slope $m$ and a point $(x_1, y_1)$. \textbf{Questions:}
    \item Write the equation of the line with a slope of $2$ that passes through the point $(1, -3)$ in point-slope form.
    \item Given a slope of $-\frac{3}{4}$ and the point $(2, 5)$, write the equation of the line in point-slope form. 3. Convert Between Point-Slope Form and Slope-Intercept Form \textbf{Objective:} Students will convert the equation of a line between point-slope form and slope-intercept form. \textbf{Questions:}
    \item Convert the following point-slope form equation to slope-intercept form: $y - 2 = 4(x + 1)$.
    \item Convert the slope-intercept form equation $y = -3x + 7$ to point-slope form using the point $(1, 4)$.
\end{enumerate}

%%===============================================================
\section{Transformation of Functions}
\label{sec:transformation-of-functions}
%%===============================================================

\paragraph{Transformation of Functions}

\paragraph{Introduction}

\paragraph{Vertical Shift}

$f(x) = x^2$ (red) and $g(x) = x^2 + 4$ (blue)

% [Image: graph]

\paragraph{Horizontal Shift}

$f(x) = \sin(x)$ (red) and $g(x) = \sin(x + \frac{\pi}{2})$ (blue)

% \href{https://upload.wikimedia.org/wikipedia/commons/thumb/2/2b/Horizontal_shift.png/300px-Horizontal_shift.png}{Image: graph}

\paragraph{Vertical Reflection}

$f(x) = \sqrt{x}$ (red) and $g(x) = -\sqrt{x}$ (blue)

% [Image: graph]

\paragraph{Horizontal Reflection}

$f(x) = e^x$ (red) and $g(x) = e^{-x}$ (blue)

% [Image: graph]

\paragraph{Vertical Stretch/Compression}

$f(x) = x^2$ (red), $g(x) = \frac{1}{2}x^2$ (blue, vertical compression), $h(x) = 2x^2$ (green, vertical stretch), and $j(x) = -2x^2$ (black, vertical stretch and reflection)

% [Image: graph]

\paragraph{Horizontal Stretch/Compression}

$f(x) = \sqrt{x}$ (red), $g(x) = \sqrt{2x}$ (blue, horizontal compression), $h(x) = \sqrt{\frac{1}{2}x}$ (green, horizontal stretch), and $j(x) = \sqrt{-2x}$ (black, horizontal compression and reflection)

% [Image: graph]

\paragraph{Translations}

A vertical shift adds or subtracts a constant to the function. For example, $g(x) = x^2 + 4$ is $f(x) = x^2$ shifted up by 4 units, and $g(x) = \sin(x) - 7.7$ is $f(x) = \sin(x)$ shifted down by 7.7 units.

A horizontal shift modifies the function as $g(x) = f(x - h)$. For example, $g(x) = (x - 3)^2$ is $f(x) = x^2$ shifted 3 units to the right, and $g(x) = \sin(x + \frac{\pi}{2})$ is $f(x) = \sin(x)$ shifted $\frac{\pi}{2}$ units to the left.

\paragraph{Reflections}

A vertical reflection changes $f(x)$ to $g(x) = -f(x)$, flipping the graph over the x-axis. For example, $g(x) = -\sqrt{x}$ is a vertical reflection of $f(x) = \sqrt{x}$.

A horizontal reflection changes $f(x)$ to $g(x) = f(-x)$, flipping the graph over the y-axis. For example, $g(x) = e^{-x}$ is a horizontal reflection of $f(x) = e^x$.

\paragraph{Even and Odd Functions}

A function $f$ is even if $f(x) = f(-x)$. For example, $f(x) = x^2$ is even. 

A function $f$ is odd if $f(x) = -f(-x)$. For example, $f(x) = x^3$ is odd.

If a function is neither even nor odd, it does not satisfy either condition. For example, $f(x) = x^2 + x$ is neither even nor odd.

\paragraph{Compressions and Stretches}

A vertical stretch/compression is represented by $g(x) = af(x)$. If $a > 1$, the graph stretches. If $0 < a < 1$, it compresses. If $a < 0$, it stretches or compresses with a reflection.

A horizontal stretch/compression is represented by $g(x) = f(bx)$. If $b > 1$, the graph compresses horizontally. If $0 < b < 1$, it stretches. If $b < 0$, it stretches or compresses with a reflection.

\paragraph{Resources}

\begin{itemize}
    \item \href{https://courses.lumenlearning.com/wmopen-collegealgebra/chapter/introduction-transformations-of-functions/}{Intro to Transformations of Functions, Lumen Learning}
\end{itemize}

\paragraph{Licensing}

Content obtained and/or adapted from:

\begin{itemize}
    \item \href{https://courses.lumenlearning.com/wmopen-collegealgebra/chapter/introduction-transformations-of-functions/}{Transformations of Functions, Lumen Learning College Algebra under a CC BY-SA license}
\end{itemize}

\subsection{Learning Objectives}

\begin{enumerate}
    \item \textbf{Vertical and Horizontal Shifts}: Given a function $f(x)$, identify and describe the effect of vertical and horizontal shifts on the function. For example, explain how $g(x) = f(x - h)$ shifts $f(x)$ horizontally.
    \item \textbf{Reflections}: Given a function $f(x)$, determine the result of vertical and horizontal reflections on the function. For example, describe how $g(x) = -f(x)$ reflects $f(x)$ over the x-axis.
    \item \textbf{Stretching and Compression}: Given a function $f(x)$, analyze the impact of vertical and horizontal stretches or compressions. For example, explain the effect of $g(x) = af(x)$ and $g(x) = f(bx)$ on the graph of $f(x)$.
\end{enumerate}

\subsection{Assessment Instrument}

\begin{enumerate}
    \item Instructions For each question, use the provided functions and transformations to answer the questions below. Show all work and provide explanations where necessary. Question 1: Vertical and Horizontal Shifts
    \item \textbf{Vertical Shift}: Given the function $f(x) = x^2$, write the equation of the function that represents a vertical shift up by 5 units.
    \item \textbf{Horizontal Shift}: Given the function $f(x) = \sin(x)$, write the equation of the function that represents a horizontal shift 4 units to the left.
    \item \textbf{Describe the Shift}: For the function $g(x) = (x - 2)^2$, describe the effect of this transformation on the graph of $f(x) = x^2$. Question 2: Reflections
    \item \textbf{Vertical Reflection}: Given the function $f(x) = \sqrt{x}$, write the equation of the function that represents a vertical reflection over the x-axis.
    \item \textbf{Horizontal Reflection}: Given the function $f(x) = e^x$, write the equation of the function that represents a horizontal reflection over the y-axis.
    \item \textbf{Describe the Reflection}: For the function $g(x) = -\sin(x)$, describe the effect of this reflection on the graph of $f(x) = \sin(x)$. Question 3: Stretching and Compression
    \item \textbf{Vertical Stretch/Compression}: Given the function $f(x) = x^2$, write the equation of the function that represents a vertical stretch by a factor of 3.
    \item \textbf{Horizontal Stretch/Compression}: Given the function $f(x) = \sqrt{x}$, write the equation of the function that represents a horizontal compression by a factor of 2.
    \item \textbf{Describe the Transformation}: For the function $g(x) = \frac{1}{2}(x^2)$, describe the effect of this transformation on the graph of $f(x) = x^2$. --- \textbf{Note}: Use appropriate graphs or sketches to illustrate the transformations where necessary.
\end{enumerate}

%%===============================================================
\section{Combinations of Functions/Composite Functions}
\label{sec:combinations-of-functions-composite-functions}
%%===============================================================

\paragraph{Composition of Functions}

In mathematics, function composition is an operation that takes two functions $f$ and $g$ and produces a function $h$ such that $h(x) = g(f(x))$. In this operation, the function $g$ is applied to the result of applying the function $f$ to $x$. That is, the functions $f : X \to Y$ and $g : Y \to Z$ are composed to yield a function that maps $x$ in $X$ to $g(f(x))$ in $Z$.

Intuitively, if $z$ is a function of $y$, and $y$ is a function of $x$, then $z$ is a function of $x$. The resulting composite function is denoted $g \circ f : X \to Z$, defined by $(g \circ f)(x) = g(f(x))$ for all $x$ in $X$. The notation $g \circ f$ is read as "g circle f", "g round f", "g about f", "g composed with f", "g after f", "g following f", "g of f", "f then g", or "g on f", or "the composition of $g$ and $f$". Intuitively, composing functions is a chaining process in which the output of function $f$ feeds the input of function $g$.

The composition of functions is a special case of the composition of relations, sometimes also denoted by $\circ$. As a result, all properties of composition of relations are true of composition of functions, though the composition of functions has some additional properties.

Composition of functions is different from multiplication of functions and has quite different properties; in particular, composition of functions is not commutative.

\paragraph{Examples}

1. \textbf{Concrete example for the composition of two functions.}

% [Image: example]

\begin{itemize}
    \item Composition of functions on a finite set: If $f = \{(1, 1), (2, 3), (3, 1), (4, 2)\}$, and $g = \{(1, 2), (2, 3), (3, 1), (4, 2)\}$, then $g \circ f = \{(1, 2), (2, 1), (3, 2), (4, 3)\}$, as shown in the figure.
\end{itemize}

2. \textbf{Composition of functions on an infinite set:}
\begin{itemize}
    \item If $f : \mathbb{R} \to \mathbb{R}$ (where $\mathbb{R}$ is the set of all real numbers) is given by $f(x) = 2x + 4$ and $g : \mathbb{R} \to \mathbb{R}$ is given by $g(x) = x^3$, then:
    \item $(f \circ g)(x) = f(g(x)) = f(x^3) = 2x^3 + 4$, and
    \item $(g \circ f)(x) = g(f(x)) = g(2x + 4) = (2x + 4)^3$.
    \item If an airplane''s altitude at time $t$ is $a(t)$, and the air pressure at altitude $x$ is $p(x)$, then $(p \circ a)(t)$ is the pressure around the plane at time $t$.
\end{itemize}

\paragraph{Properties}

\begin{itemize}
    \item The composition of functions is always associative—a property inherited from the composition of relations. That is, if $f$, $g$, and $h$ are composable, then $f \circ (g \circ h) = (f \circ g) \circ h$. Since the parentheses do not change the result, they are generally omitted.
\end{itemize}

\begin{itemize}
    \item In a strict sense, the composition $g \circ f$ is only meaningful if the codomain of $f$ equals the domain of $g$; in a wider sense, it is sufficient that the former be a subset of the latter. Moreover, it is often convenient to tacitly restrict the domain of $f$, such that $f$ produces only values in the domain of $g$. For example, the composition $g \circ f$ of the functions $f : \mathbb{R} \to \mathbb{R}$ defined by $f(x) = 9 - x^2$ and $g : \mathbb{R} \to \mathbb{R}$ defined by $g(x) = \sqrt{x}$ can be defined on the interval $[-3, +3]$.
\end{itemize}

\begin{itemize}
    \item The functions $g$ and $f$ are said to commute with each other if $g \circ f = f \circ g$. Commutativity is a special property, attained only by particular functions, and often in special circumstances. For example, $|x| + 3 = |x + 3|$ only when $x \geq 0$. The picture shows another example.
\end{itemize}

% [Image: graph]

\begin{itemize}
    \item The composition of one-to-one (injective) functions is always one-to-one. Similarly, the composition of onto (surjective) functions is always onto. It follows that the composition of two bijections is also a bijection. The inverse function of a composition (assumed invertible) has the property that $(f \circ g)^{-1} = g^{-1} \circ f^{-1}$.
\end{itemize}

\paragraph{Resources}

\begin{itemize}
    \item \href{https://www.khanacademy.org/math/precalculus/x9e81a4f98389efdf:composite/x9e81a4f98389efdf:composing/v/function-composition}{Intro to Composing Functions, Khan Academy}
    \item \href{https://courses.lumenlearning.com/wmopen-collegealgebra/chapter/introduction-compositions-of-functions/}{Compositions of Functions, Lumen Learning}
    \item \href{https://www.mathsisfun.com/sets/functions-composition.html}{Composition of Functions, Math Is Fun}
\end{itemize}

\paragraph{Licensing}

Content obtained and/or adapted from:

\begin{itemize}
    \item \href{https://en.wikipedia.org/wiki/Function_composition}{Function composition, Wikipedia under a CC BY-SA license}
\end{itemize}

\subsection{Learning Objectives}

\begin{enumerate}
    \item \textbf{1. Understand and Apply Function Composition}:
    \item Students will be able to compose two functions $f$ and $g$ to form a new function $h$ such that $h(x) = g(f(x))$. \textbf{2. Identify and Interpret Composite Functions}:
    \item Students will be able to identify composite functions and interpret the chaining process of applying $f$ and $g$ in different orders. \textbf{3. Properties and Characteristics of Composite Functions}:
    \item Students will be able to explain the associative property of function composition, recognize non-commutativity, and understand the conditions under which functions commute.
\end{enumerate}

\subsection{Assessment Instrument}

\begin{enumerate}
    \item Instructions Answer the following questions about the composition of functions. Show all work and justify your answers where necessary. Questions 1. Understanding Function Composition Given the functions $f(x) = 3x + 2$ and $g(x) = x^2 - 1$: a. Find the composite function $(g \circ f)(x)$. b. Find the composite function $(f \circ g)(x)$. c. Are $(g \circ f)(x)$ and $(f \circ g)(x)$ equal? Explain your reasoning. 2. Real-World Application An airplane's altitude at time $t$ is given by $a(t) = 1000t + 500$, and the air pressure at altitude $x$ is given by $p(x) = 101.3 - 0.01x$: a. Express the air pressure around the plane as a function of time $t$ using composition. b. Calculate the air pressure at $t = 5$ hours. 3. Associativity of Function Composition Given three functions $f(x) = 2x$, $g(x) = x + 1$, and $h(x) = \sqrt{x}$: a. Verify the associative property by calculating $f \circ (g \circ h)(x)$ and $(f \circ g) \circ h(x)$. b. Show that the results are the same. 4. Domain and Codomain Given the functions $f : \mathbb{R} \to \mathbb{R}$ defined by $f(x) = \sqrt{x}$ and $g : \mathbb{R} \to \mathbb{R}$ defined by $g(x) = x^2$: a. Determine the domain and codomain of the composite function $(g \circ f)(x)$. b. Determine the domain and codomain of the composite function $(f \circ g)(x)$. 5. Commutativity of Functions Consider the functions $f(x) = x^3$ and $g(x) = x + 2$: a. Compute $(f \circ g)(x)$ and $(g \circ f)(x)$. b. Determine if $f$ and $g$ commute with each other. Justify your answer. Extra Credit Explain why the composition of one-to-one (injective) functions is always one-to-one, and why the composition of onto (surjective) functions is always onto.
\end{enumerate}

%%===============================================================
\section{Inverse Functions}
\label{sec:inverse-functions}
%%===============================================================

\paragraph{Inverse Functions}

In mathematics, an inverse function (or anti-function) is a function that "reverses" another function: if the function $f$ applied to an input $x$ gives a result of $y$, then applying its inverse function $g$ to $y$ gives the result $x$, i.e., $g(y) = x$ if and only if $f(x) = y$. The inverse function of $f$ is also denoted as $f^{-1}$.

As an example, consider the real-valued function of a real variable given by $f(x) = 5x - 7$. Thinking of this as a step-by-step procedure (namely, take a number $x$, multiply it by 5, then subtract 7 from the result), to reverse this and get $x$ back from some output value, say $y$, we would undo each step in reverse order. In this case, it means to add 7 to $y$, and then divide the result by 5. In functional notation, this inverse function would be given by,

$$
g(y) = \frac{y + 7}{5}.
$$

With $y = 5x - 7$ we have that $f(x) = y$ and $g(y) = x$.

Not all functions have inverse functions. Those that do are called invertible. For a function $f: X \to Y$ to have an inverse, it must have the property that for every $y$ in $Y$, there is exactly one $x$ in $X$ such that $f(x) = y$. This property ensures that a function $g: Y \to X$ exists with the necessary relationship with $f$.

\paragraph{Definitions}

If $f$ maps $X$ to $Y$, then $f^{-1}$ maps $Y$ back to $X$.
Let $f$ be a function whose domain is the set $X$, and whose codomain is the set $Y$. Then $f$ is invertible if there exists a function $g$ with domain $Y$ and codomain $X$, with the property:

$$
f(x) = y \Leftrightarrow g(y) = x.
$$

If $f$ is invertible, then the function $g$ is unique, which means that there is exactly one function $g$ satisfying this property. Moreover, it also follows that the ranges of $g$ and $f$ equal their respective codomains. The function $g$ is called the inverse of $f$, and is usually denoted as $f^{-1}$, a notation introduced by John Frederick William Herschel in 1813.

Stated otherwise, a function, considered as a binary relation, has an inverse if and only if the converse relation is a function on the codomain $Y$, in which case the converse relation is the inverse function. Not all functions have an inverse. For a function to have an inverse, each element $y \in Y$ must correspond to no more than one $x \in X$; a function $f$ with this property is called one-to-one or an injection. If $f^{-1}$ is to be a function on $Y$, then each element $y \in Y$ must correspond to some $x \in X$. Functions with this property are called surjections. This property is satisfied by definition if $Y$ is the image of $f$, but may not hold in a more general context. To be invertible, a function must be both an injection and a surjection. Such functions are called bijections. The inverse of an injection $f: X \to Y$ that is not a bijection (that is, not a surjection), is only a partial function on $Y$, which means that for some $y \in Y$, $f^{-1}(y)$ is undefined. If a function $f$ is invertible, then both it and its inverse function $f^{-1}$ are bijections.

% [Image: mapping]

Another convention is used in the definition of functions, referred to as the "set-theoretic" or "graph" definition using ordered pairs, which makes the codomain and image of the function the same. Under this convention, all functions are surjective, so bijectivity and injectivity are the same. Authors using this convention may use the phrasing that a function is invertible if and only if it is an injection. The two conventions need not cause confusion, as long as it is remembered that in this alternate convention, the codomain of a function is always taken to be the image of the function.

\paragraph{Properties}

Since a function is a special type of binary relation, many of the properties of an inverse function correspond to properties of converse relations.

\paragraph{Uniqueness}

If an inverse function exists for a given function $f$, then it is unique. This follows since the inverse function must be the converse relation, which is completely determined by $f$.

\paragraph{Symmetry}

There is a symmetry between a function and its inverse. Specifically, if $f$ is an invertible function with domain $X$ and codomain $Y$, then its inverse $f^{-1}$ has domain $Y$ and image $X$, and the inverse of $f^{-1}$ is the original function $f$. In symbols, for functions $f: X \to Y$ and $f^{-1}: Y \to X$,

$$
f^{-1} \circ f = \operatorname{id}\_X
$$

and 

$$
f \circ f^{-1} = \operatorname{id}\_Y.
$$

This statement is a consequence of the implication that for $f$ to be invertible it must be bijective. The involutory nature of the inverse can be concisely expressed by

$$
\left(f^{-1}\right)^{-1} = f.
$$

The inverse of a composition of functions is given by

$$
(g \circ f)^{-1} = f^{-1} \circ g^{-1}.
$$

% [Image: inverse]

Notice that the order of $g$ and $f$ have been reversed; to undo $f$ followed by $g$, we must first undo $g$, and then undo $f$.

For example, let $f(x) = 3x$ and let $g(x) = x + 5$. Then the composition $g \circ f$ is the function that first multiplies by three and then adds five,

$$
(g \circ f)(x) = 3x + 5.
$$

To reverse this process, we must first subtract five, and then divide by three,

$$
(g \circ f)^{-1}(x) = \frac{1}{3}(x - 5).
$$

This is the composition $(f^{-1} \circ g^{-1})(x)$.

\paragraph{Self-inverses}

If $X$ is a set, then the identity function on $X$ is its own inverse:

$$
\operatorname{id}\_X^{-1} = \operatorname{id}\_X.
$$

More generally, a function $f : X \to X$ is equal to its own inverse, if and only if the composition $f \circ f$ is equal to $\operatorname{id}\_X$. Such a function is called an involution.

\paragraph{Resources}

\begin{itemize}
    \item \href{https://courses.lumenlearning.com/collegealgebra2017/chapter/introduction-inverse-functions/}{Inverse Functions, Lumen Learning}
    \item \href{https://www.khanacademy.org/math/algebra/x2f8bb11595b61c86:functions/x2f8bb11595b61c86:inverse-functions-intro/v/introduction-to-function-inverses}{Introduction to Inverse Function Video, Khan Academy}
\end{itemize}

\paragraph{Licensing}

Content obtained and/or adapted from:

\begin{itemize}
    \item \href{https://en.wikipedia.org/wiki/Inverse_function}{Inverse function, Wikipedia under a CC BY-SA license}
\end{itemize}

\subsection{Learning Objectives}

\begin{enumerate}
    \item Identify and denote the inverse function $f^{-1}$ for a given function $f$ and understand the notation and concept of inverse functions.
    \item Calculate the inverse of a linear function by reversing the operations of the original function and express the inverse function in functional notation.
    \item Determine if a function is invertible by analyzing if it is bijective, ensuring it meets the criteria of being both one-to-one and onto.
\end{enumerate}

\subsection{Assessment Instrument}

\begin{enumerate}
    \item 1. Identifying Inverse Functions \textbf{Question:} Given the function $f(x) = 4x + 3$, find its inverse function $f^{-1}$. Express the inverse function in functional notation. 2. Verifying Inverse Properties \textbf{Question:} Consider the functions $f(x) = 2x - 5$ and $g(x) = \frac{x + 5}{2}$. Verify if $g$ is the inverse of $f$ by showing that $f(g(x)) = x$ and $g(f(x)) = x$. 3. Composition of Inverse Functions \textbf{Question:} If $h(x) = x - 7$ and $k(x) = x + 7$, find the composition of $h$ and $k$ and determine its inverse.
\end{enumerate}

%%===============================================================
\section{Quadratic Functions}
\label{sec:quadratic-functions}
%%===============================================================

\paragraph{Quadratic Functions}

In algebra, a quadratic function, a quadratic polynomial, or a polynomial of degree 2, is a polynomial function where the highest-degree term is of the second degree.

A quadratic polynomial with two real roots (crossings of the x-axis) and no complex roots can be written as:

$$
f(x) = ax^2 + bx + c, \quad a \neq 0
$$

% [Image: graph]

In this case, the graph of the function is a parabola with an axis of symmetry parallel to the y-axis. If the quadratic function is set to zero, it forms a quadratic equation, and the solutions are called the roots of the function.

For a bivariate case involving variables $x$ and $y$, the quadratic function has the form:

$$
f(x, y) = ax^2 + by^2 + cxy + dx + ey + f
$$

where at least one of $a$, $b$, or $c$ is non-zero. Setting this function equal to zero results in a conic section, which could be a circle, ellipse, parabola, or hyperbola.

In three variables $x$, $y$, and $z$, a quadratic function is expressed as:

$$
f(x, y, z) = ax^2 + by^2 + cz^2 + dxy + exz + fyz + gx + hy + iz + j
$$

where at least one of the coefficients $a$, $b$, $c$, $d$, $e$, or $f$ is non-zero. For more than three variables, the resulting surface is called a quadric, with the highest degree term being of degree 2.

\paragraph{Etymology}

The adjective \textit{quadratic} comes from the Latin word \textit{quadratum} ("square"). In algebra, a term like $x^2$ is called a square because it represents the area of a square with side $x$.

\paragraph{Terminology}

\paragraph{Coefficients}

The coefficients of a polynomial are often real or complex numbers, but a polynomial may be defined over any ring.

\paragraph{Degree}

When referring to a "quadratic polynomial," authors may mean either "having degree exactly 2" or "having degree at most 2." If the degree is less than 2, this is sometimes called a "degenerate case." The term "order" is also used to denote "degree," e.g., a second-order polynomial.

\paragraph{Variables}

A quadratic polynomial may involve a single variable $x$ (the univariate case) or multiple variables such as $x$, $y$, and $z$ (the multivariate case).

\paragraph{The One-Variable Case}

Any single-variable quadratic polynomial can be written as:

$$
f(x) = ax^2 + bx + c
$$

where $x$ is the variable, and $a$, $b$, and $c$ are the coefficients. In elementary algebra, such polynomials often appear in the form of a quadratic equation:

$$
ax^2 + bx + c = 0
$$

The solutions to this equation are called the roots of the quadratic polynomial and can be found through factorization, completing the square, graphing, Newton''s method, or using the quadratic formula. Each quadratic polynomial has an associated quadratic function, whose graph is a parabola.

\paragraph{Bivariate Case}

Any quadratic polynomial with two variables can be written as:

$$
f(x, y) = ax^2 + by^2 + cxy + dx + ey + f
$$

where $x$ and $y$ are the variables and $a$, $b$, $c$, $d$, $e$, and $f$ are the coefficients. Such polynomials are fundamental to the study of conic sections, characterized by setting $f(x, y)$ to zero. Quadratic polynomials with three or more variables correspond to quadric surfaces and hypersurfaces. In linear algebra, quadratic polynomials can be generalized to quadratic forms on vector spaces.

\paragraph{Forms of a Univariate Quadratic Function}

A univariate quadratic function can be expressed in three formats:

1. \textbf{Standard Form}:
   $$
   f(x) = ax^2 + bx + c
   $$

2. \textbf{Factored Form}:
   $$
   f(x) = a(x - r\_1)(x - r\_2)
   $$
   where $r\_1$ and $r\_2$ are the roots of the quadratic function.

3. \textbf{Vertex Form}:
   $$
   f(x) = a(x - h)^2 + k
   $$
   where $h$ and $k$ are the $x$ and $y$ coordinates of the vertex, respectively.

The coefficient $a$ is the same in all three forms. To convert the standard form to factored form, one uses the quadratic formula to determine the roots $r\_1$ and $r\_2$. Converting to vertex form involves completing the square. To convert from factored form (or vertex form) to standard form, one multiplies, expands, and/or distributes the factors.

\paragraph{Graph of the Univariate Function}

1. For $a = \{0.1, 0.3, 1, 3\}$:
   $$
   f(x) = ax^2
   $$

2. For $b = \{1, 2, 3, 4\}$:
   $$
   f(x) = x^2 + bx
   $$

3. For $b = \{-1, -2, -3, -4\}$:
   $$
   f(x) = x^2 + bx
   $$

Regardless of the format, the graph of a univariate quadratic function:

$$
f(x) = ax^2 + bx + c
$$

is a parabola. Specifically:

\begin{itemize}
    \item If $a > 0$, the parabola opens upwards.
    \item If $a < 0$, the parabola opens downwards.
    \item The coefficient $a$ affects the curvature of the graph; a larger magnitude of $a$ results in a more sharply curved graph.
\end{itemize}

The coefficients $b$ and $a$ together control the axis of symmetry of the parabola (the $x$-coordinate of the vertex and the $h$ parameter in the vertex form), which is located at:

$$
x = -\frac{b}{2a}
$$

The coefficient $c$ controls the height of the parabola where it intercepts the $y$-axis.

% [Image: graph]

\paragraph{Vertex}

The vertex of a parabola is the point where it turns, also known as the turning point. If the quadratic function is in vertex form, the vertex is $(h, k)$. By completing the square, the standard form:

$$
f(x) = ax^2 + bx + c
$$

can be converted into:

$$
f(x) = a(x - h)^2 + k = a\left(x - \frac{-b}{2a}\right)^2 + \left(c - \frac{b^2}{4a}\right)
$$

So the vertex $(h, k)$ in standard form is:

$$
\left(-\frac{b}{2a}, c - \frac{b^2}{4a}\right)
$$

If the quadratic function is in factored form:

$$
f(x) = a(x - r\_1)(x - r\_2)
$$

the average of the two roots, $\frac{r\_1 + r\_2}{2}$, is the $x$-coordinate of the vertex, and thus the vertex $(h, k)$ is:

$$
\left(\frac{r\_1 + r\_2}{2}, f\left(\frac{r\_1 + r\_2}{2}\right)\right)
$$

The vertex is the maximum point if $a < 0$, or the minimum point if $a > 0.

The vertical line:

$$
x = h = -\frac{b}{2a}
$$

that passes through the vertex is also the axis of symmetry of the parabola.

% [Image: graph]

% [Image: graph]

\paragraph{Maximum and Minimum Points}

Using calculus, the vertex point, being a maximum or minimum of the function, can be found by solving for the roots of the derivative:

$$
f(x) = ax^2 + bx + c \Rightarrow f''(x) = 2ax + b
$$

Setting $f''(x) = 0$ results in:

$$
x = -\frac{b}{2a}
$$

The corresponding function value is:

$$
f(x) = a\left(-\frac{b}{2a}\right)^2 + b\left(-\frac{b}{2a}\right) + c = c - \frac{b^2}{4a}
$$

Thus, the vertex point coordinates $(h, k)$ can be expressed as:

$$
\left(-\frac{b}{2a}, c - \frac{b^2}{4a}\right)
$$

\paragraph{Roots of the Univariate Function}

\textbf{Graph of $y = ax^2 + bx + c$}, where $a$ and the discriminant $b^2 - 4ac$ are positive, with:
\begin{itemize}
    \item Roots and y-intercept in red
    \item Vertex and axis of symmetry in blue
    \item Focus and directrix in pink
\end{itemize}

% [Image: graph]

\textbf{Visualisation of the complex roots of $y = ax^2 + bx + c$}: The parabola is rotated 180° about its vertex (orange). Its x-intercepts are rotated 90° around their midpoint, and the Cartesian plane is interpreted as the complex plane (green).

% [Image: graph]

\paragraph{Exact Roots}

The roots (or zeros), $r\_1$ and $r\_2$, of the univariate quadratic function

$$
f(x) = ax^2 + bx + c = a(x - r\_1)(x - r\_2),
$$

are the values of $x$ for which $f(x) = 0$.

When the coefficients $a$, $b$, and $c$ are real or complex, the roots are:

$$
r\_1 = \frac{-b - \sqrt{b^2 - 4ac}}{2a},
$$

$$
r\_2 = \frac{-b + \sqrt{b^2 - 4ac}}{2a}.
$$

\paragraph{Upper Bound on the Magnitude of the Roots}

The modulus of the roots of a quadratic $ax^2 + bx + c$ can be no greater than

$$
\frac{\max(|a|, |b|, |c|)}{|a|} \times \phi,
$$

where $\phi$ is the golden ratio $\frac{1 + \sqrt{5}}{2}$.

\paragraph{The Square Root of a Univariate Quadratic Function}

The square root of a univariate quadratic function gives rise to one of the four conic sections, almost always either an ellipse or a hyperbola.

If $a > 0$, then the equation

$$
y = \pm \sqrt{ax^2 + bx + c}
$$

describes a hyperbola, as can be seen by squaring both sides. The directions of the axes of the hyperbola are determined by the ordinate of the minimum point of the corresponding parabola $y\_p = ax^2 + bx + c$. If the ordinate is negative, then the hyperbola''s major axis (through its vertices) is horizontal, while if the ordinate is positive then the hyperbola''s major axis is vertical.

If $a < 0$, then the equation

$$
y = \pm \sqrt{ax^2 + bx + c}
$$

describes either a circle or another ellipse or nothing at all. If the ordinate of the maximum point of the corresponding parabola $y\_p = ax^2 + bx + c$ is positive, then its square root describes an ellipse, but if the ordinate is negative then it describes an empty locus of points.

\paragraph{Iteration}

To iterate a function $f(x) = ax^2 + bx + c$, one applies the function repeatedly, using the output from one iteration as the input to the next.

One cannot always deduce the analytic form of $f^{(n)}(x)$, which means the $n$th iteration of $f(x)$. (The superscript can be extended to negative numbers, referring to the iteration of the inverse of $f(x)$ if the inverse exists.) But there are some analytically tractable cases.

For example, for the iterative equation

$$
f(x) = a(x - c)^2 + c,
$$

one has

$$
f(x) = a(x - c)^2 + c = h^{-1}(g(h(x))),
$$

where

$$
g(x) = ax^2
$$

and

$$
h(x) = x - c.
$$

So by induction,

$$
f^{(n)}(x) = h^{-1}(g^{(n)}(h(x)))
$$

can be obtained, where $g^{(n)}(x)$ can be easily computed as

$$
g^{(n)}(x) = a^{2^n - 1} x^{2^n}.
$$

Finally, we have

$$
f^{(n)}(x) = a^{2^n - 1}(x - c)^{2^n} + c
$$

as the solution.

See Topological conjugacy for more detail about the relationship between $f$ and $g$. And see Complex quadratic polynomial for the chaotic behavior in the general iteration.

\paragraph{The Logistic Map}

$$
x\_{n+1} = rx\_n (1 - x\_n), \quad 0 \leq x\_0 < 1
$$

with parameter $2 < r < 4$ can be solved in certain cases, one of which is chaotic and one of which is not. In the chaotic case $r = 4$, the solution is

$$
x\_n = \sin^2(2^n \theta \pi)
$$

where the initial condition parameter $\theta$ is given by

$$
\theta = \frac{1}{\pi} \sin^{-1}(x\_0^{1/2}).
$$

For rational $\theta$, after a finite number of iterations $x\_n$ maps into a periodic sequence. But almost all $\theta$ are irrational, and for irrational $\theta$, $x\_n$ never repeats itself – it is non-periodic and exhibits sensitive dependence on initial conditions, so it is said to be chaotic.

The solution of the logistic map when $r = 2$ is

$$
x\_n = \frac{1}{2} - \frac{1}{2} (1 - 2x\_0)^{2^n}
$$

for $x\_0 \in [0, 1)$. Since $(1 - 2x\_0) \in (-1, 1)$ for any value of $x\_0$ other than the unstable fixed point 0, the term $(1 - 2x\_0)^{2^n}$ goes to 0 as $n$ goes to infinity, so $x\_n$ goes to the stable fixed point $\frac{1}{2}$.

\paragraph{Bivariate (Two-Variable) Quadratic Function}

A bivariate quadratic function is a second-degree polynomial of the form

$$
f(x, y) = Ax^2 + By^2 + Cx + Dy + Exy + F
$$

where $A$, $B$, $C$, $D$, and $E$ are fixed coefficients and $F$ is the constant term. Such a function describes a quadratic surface. Setting $f(x, y)$ equal to zero describes the intersection of the surface with the plane $z = 0$, which is a locus of points equivalent to a conic section.

\paragraph{Minimum/Maximum}

If $4AB - E^2 < 0$, the function has no maximum or minimum; its graph forms a hyperbolic paraboloid.

If $4AB - E^2 > 0$, the function has a minimum if $A > 0$, and a maximum if $A < 0$; its graph forms an elliptic paraboloid. In this case, the minimum or maximum occurs at $(x\_m, y\_m)$ where:

$$
x\_m = -\frac{2BC - DE}{4AB - E^2},
$$

$$
y\_m = -\frac{2AD - CE}{4AB - E^2}.
$$

If $4AB - E^2 = 0$ and $DE - 2CB = 2AD - CE \neq 0$, the function has no maximum or minimum; its graph forms a parabolic cylinder.

If $4AB - E^2 = 0$ and $DE - 2CB = 2AD - CE = 0$, the function achieves the maximum/minimum at a line—a minimum if $A > 0$ and a maximum if $A < 0$; its graph forms a parabolic cylinder.

\paragraph{Resources}

\begin{itemize}
    \item \href{https://mathresearch.utsa.edu/wikiFiles/MAT1053/Quadratic_Functions/MAT1053_M2.2Quadratic_Functions.pdf}{Quadratic Functions, Book Chapters}
    \item \href{https://mathresearch.utsa.edu/wikiFiles/MAT1053/Quadratic_Functions/MAT1053_M2.2Quadratic_FunctionsGN.pdf}{Guided Notes}
\end{itemize}

\paragraph{Licensing}

Content obtained and/or adapted from:

\begin{itemize}
    \item \href{https://en.wikipedia.org/wiki/Quadratic_function}{Quadratic function, Wikipedia under a CC BY-SA license}
\end{itemize}

\subsection{Learning Objectives}

\begin{enumerate}
    \item \textbf{Identify and Interpret the Vertex Form of a Quadratic Function} Students will be able to identify and convert between the standard form, factored form, and vertex form of a quadratic function. They will be able to:
    \item Recognize and use the vertex form of a quadratic function, $f(x) = a(x - h)^2 + k$, to determine the vertex $(h, k)$ of the parabola.
    \item Convert between the standard form $f(x) = ax^2 + bx + c$ and vertex form by completing the square.
    \item Understand how the vertex form reveals the minimum or maximum point of the quadratic function depending on the sign of $a$.
    \item \textbf{Calculate and Interpret the Roots of a Quadratic Function} Students will be able to:
    \item Find the roots of a quadratic function using the quadratic formula: $$ r\_1 = \frac{-b - \sqrt{b^2 - 4ac}}{2a}, \quad r\_2 = \frac{-b + \sqrt{b^2 - 4ac}}{2a} $$
    \item Interpret the nature of the roots based on the discriminant $b^2 - 4ac$, identifying real, repeated, or complex roots.
    \item Graph the quadratic function $f(x) = ax^2 + bx + c$ and accurately plot the roots and the y-intercept.
    \item \textbf{Apply Quadratic Functions to Model Real-World Situations} Students will be able to:
    \item Use quadratic functions to model real-world scenarios such as projectile motion, optimization problems, or financial calculations.
    \item Analyze the quadratic function to determine key features such as the vertex, axis of symmetry, and roots in the context of the application.
    \item Solve practical problems by setting up and solving quadratic equations, and interpret the solutions in terms of the original context.
\end{enumerate}

\subsection{Assessment Instrument}

\begin{enumerate}
    \item Instructions
    \item Answer all questions.
    \item Show all work where applicable.
    \item You may use a calculator, but ensure all steps are shown. Part 1: Definitions and Properties Question 1: Multiple Choice Which of the following represents the standard form of a quadratic function?
    \item $ f(x) = ax^2 + bx + c $
    \item $ f(x) = a(x - r\_1)(x - r\_2) $
    \item $ f(x) = a(x - h)^2 + k $
    \item $ f(x, y) = Ax^2 + By^2 + Cx + Dy + Exy + F $ Question 2: True or False The vertex of the parabola $ f(x) = ax^2 + bx + c $ can be found using the formula $ \left(-\frac{b}{2a}, c - \frac{b^2}{4a}\right) $. (True/False) Question 3: Matching Match the following forms of a quadratic function with their descriptions:
    \item \textbf{Standard Form}: $ f(x) = ax^2 + bx + c $
    \item \textbf{Factored Form}: $ f(x) = a(x - r\_1)(x - r\_2) $
    \item \textbf{Vertex Form}: $ f(x) = a(x - h)^2 + k $ \textbf{Descriptions:} a. Expresses the function with roots $ r\_1 $ and $ r\_2 $. b. Expresses the function with vertex coordinates $ (h, k) $. c. General form of the quadratic function. Part 2: Graphing and Analysis Question 4: Graphing Given the quadratic function $ f(x) = 2x^2 - 4x + 1 $:
    \item Plot the graph of the function.
    \item Identify the vertex of the parabola.
    \item Determine the axis of symmetry.
    \item Find the x-intercepts (roots) and y-intercept. Question 5: Vertex Form Conversion Convert the following quadratic function from standard form to vertex form: $$ f(x) = x^2 - 6x + 8 $$ Question 6: Roots Calculation For the quadratic function $ f(x) = 3x^2 - 12x + 7 $, find the exact roots using the quadratic formula. Part 3: Applications Question 7: Real-World Problem A projectile is launched with an initial velocity such that its height $ h(t) $ in meters after $ t $ seconds is given by: $$ h(t) = -4.9t^2 + 20t + 5 $$
    \item What is the maximum height reached by the projectile?
    \item After how many seconds does the projectile reach this maximum height?
    \item When does the projectile hit the ground? Question 8: Complex Roots Find the roots of the quadratic function: $$ f(x) = x^2 + 2x + 5 $$ Part 4: Higher-Level Concepts Question 9: Bivariate Quadratic Function Consider the bivariate quadratic function: $$ f(x, y) = 2x^2 + 3y^2 - 4xy + 5x + 6y + 7 $$
    \item Determine if the function has a maximum or minimum value.
    \item If so, find the coordinates of the minimum or maximum point. Question 10: Iteration For the function $ f(x) = 2(x - 1)^2 + 3 $, compute the result of iterating the function 3 times starting with $ x\_0 = 2 $. Part 5: Extra Credit Question 11: Logistic Map Solve the logistic map for $ r = 3.5 $ and initial condition $ x\_0 = 0.4 $ after 5 iterations. Describe the behavior of the sequence.
\end{enumerate}

%%===============================================================
\section{Graphs of Polynomials}
\label{sec:graphs-of-polynomials}
%%===============================================================

\paragraph{Graphs of Polynomials}

\paragraph{Polynomial Functions}

\paragraph{Polynomial of Degree 0:}

$
f(x) = 2
$

% [Image: graph]

\paragraph{Polynomial of Degree 1:}

$
f(x) = 2x + 1
$

% [Image: graph]

\paragraph{Polynomial of Degree 2:}

$
f(x) = x^2 - x - 2
$
$
= (x + 1)(x - 2)
$

% [Image: graph]

\paragraph{Polynomial of Degree 3:}

$
f(x) = \frac{1}{4}x^3 + \frac{3}{4}x^2 - \frac{3}{2}x - 2
$
$
= \frac{1}{4} (x + 4)(x + 1)(x - 2)
$

% [Image: graph]

\paragraph{Polynomial of Degree 4:}

$
f(x) = \frac{1}{14} (x + 4)(x + 1)(x - 1)(x - 3) + 0.5
$

% [Image: graph]

\paragraph{Polynomial of Degree 5:}

$
f(x) = \frac{1}{20} (x + 4)(x + 2)(x + 1)(x - 1)(x - 3) + 2
$

% [Image: graph]

\paragraph{Polynomial of Degree 6:}

$
f(x) = \frac{1}{100} (x^6 - 2x^5 - 26x^4 + 28x^3 + 145x^2 - 26x - 80)
$

% [Image: graph]

\paragraph{Polynomial of Degree 7:}

$
f(x) = (x - 3)(x - 2)(x - 1)x(x + 1)(x + 2)(x + 3)
$

% [Image: graph]

\paragraph{Polynomial Functions Definition}

A polynomial function is a function that can be defined by evaluating a polynomial. More precisely, a function $f$ of one argument from a given domain is a polynomial function if there exists a polynomial:

$$
a\_n x^n + a\_{n-1} x^{n-1} + \cdots + a\_2 x^2 + a\_1 x + a\_0
$$

that evaluates to $f(x)$ for all $x$ in the domain of $f$ (here, $n$ is a non-negative integer and $a\_0, a\_1, a\_2, \ldots, a\_n$ are constant coefficients). Generally, unless otherwise specified, polynomial functions have complex coefficients, arguments, and values. In particular, a polynomial, restricted to have real coefficients, defines a function from the complex numbers to the complex numbers. If the domain of this function is also restricted to the reals, the resulting function is a real function that maps reals to reals.

For example, the function $f$, defined by:

$$
f(x) = x^3 - x
$$

is a polynomial function of one variable. Polynomial functions of several variables are similarly defined, using polynomials in more than one indeterminate, as in:

$$
f(x, y) = 2x^3 + 4x^2 y + xy^5 + y^2 - 7
$$

According to the definition of polynomial functions, there may be expressions that obviously are not polynomials but nevertheless define polynomial functions. An example is the expression:

$$
\left(\sqrt{1 - x^2}\right)^2
$$

which takes the same values as the polynomial:

$$
1 - x^2
$$

on the interval:

$$
[-1, 1]
$$

and thus both expressions define the same polynomial function on this interval.

Every polynomial function is continuous, smooth, and entire.

\paragraph{Graphs}

A polynomial function in one real variable can be represented by a graph.

\begin{itemize}
    \item The graph of the zero polynomial $f(x) = 0$ is the x-axis.
    \item The graph of a degree 0 polynomial $f(x) = a\_0$, where $a\_0 \neq 0$, is a horizontal line with y-intercept $a\_0$.
    \item The graph of a degree 1 polynomial (or linear function) $f(x) = a\_0 + a\_1 x$, where $a\_1 \neq 0$, is an oblique line with y-intercept $a\_0$ and slope $a\_1$.
    \item The graph of a degree 2 polynomial $f(x) = a\_0 + a\_1 x + a\_2 x^2$, where $a\_2 \neq 0$, is a parabola.
    \item The graph of a degree 3 polynomial $f(x) = a\_0 + a\_1 x + a\_2 x^2 + a\_3 x^3$, where $a\_3 \neq 0$, is a cubic curve.
    \item The graph of any polynomial with degree 2 or greater $f(x) = a\_0 + a\_1 x + a\_2 x^2 + \cdots + a\_n x^n$, where $a\_n \neq 0$ and $n \geq 2$, is a continuous non-linear curve.
\end{itemize}

A non-constant polynomial function tends to infinity when the variable increases indefinitely (in absolute value). If the degree is higher than one, the graph does not have any asymptote. It has two parabolic branches with vertical direction (one branch for positive $x$ and one for negative $x$).

Polynomial graphs are analyzed in calculus using intercepts, slopes, concavity, and end behavior.

\paragraph{Resources}

\begin{itemize}
    \item \href{https://tutorial.math.lamar.edu/classes/alg/graphingpolynomials.aspx}{Graphing Polynomials, Paul''s Online Notes from Lamar University}
    \item \href{https://courses.lumenlearning.com/collegealgebra2017/chapter/graphs-of-polynomial-functions/}{Graphs of Polynomial Functions, Lumen Learning College Algebra}
    \item \href{https://en.wikipedia.org/wiki/Polynomial}{Polynomial, Wikipedia}
\end{itemize}

\paragraph{Licensing}

Content obtained and/or adapted from:

\begin{itemize}
    \item \href{https://en.wikipedia.org/wiki/Polynomial}{Polynomial, Wikipedia under a CC BY-SA license}
\end{itemize}

\subsection{Learning Objectives}

\begin{enumerate}
    \item SRO 1: Identify Polynomial Degrees \textbf{Identify the degree of the polynomial function given its equation and graph.} SRO 2: Analyze Polynomial Graphs \textbf{Analyze the graph of a polynomial function and determine its degree and general shape.} SRO 3: Polynomial Functions Definition \textbf{Define and identify polynomial functions based on their algebraic expressions and graph representations.}
\end{enumerate}

\subsection{Assessment Instrument}

\begin{enumerate}
    \item Instructions Answer the following questions based on the definition and properties of polynomial functions. --- Question 1 What is a polynomial function? Provide the general form of a polynomial function and describe what each term represents. $$ a\_n x^n + a\_{n-1} x^{n-1} + \cdots + a\_2 x^2 + a\_1 x + a\_0 $$ \textbf{Questions:}
    \item What does each coefficient $a\_n, a\_{n-1}, \ldots, a\_0$ represent in a polynomial function?
    \item What does the term $n$ denote in the polynomial function? --- Question 2 Determine whether the following function is a polynomial function: $$ f(x) = \frac{1}{1 + x^2} $$ \textbf{Questions:}
    \item Explain why or why not the given function is a polynomial function.
    \item If the function is not a polynomial, specify what type of function it is. --- Question 3 Consider the polynomial function: $$ f(x) = x^3 - x $$ \textbf{Questions:}
    \item Identify the degree of this polynomial function.
    \item Describe the general shape of the graph of this polynomial function. --- Question 4 What is the general shape of the graph for a degree 2 polynomial function, such as: $$ f(x) = a\_0 + a\_1 x + a\_2 x^2 $$ where $a\_2 \neq 0$? \textbf{Questions:}
    \item Describe the shape of the graph for a degree 2 polynomial.
    \item What are the characteristics of the graph (e.g., the presence of symmetry, shape, etc.)? --- Question 5 For a polynomial function of degree 3, such as: $$ f(x) = a\_0 + a\_1 x + a\_2 x^2 + a\_3 x^3 $$ \textbf{Questions:}
    \item How does the graph of a degree 3 polynomial function generally behave as $x$ increases or decreases indefinitely?
    \item What can be inferred about the end behavior of the graph for a degree 3 polynomial function? --- Scoring Rubric
    \item \textbf{Definition and General Form:} 2 points
    \item \textbf{Function Identification and Description:} 2 points
    \item \textbf{Graph Description:} 2 points
    \item \textbf{End Behavior and Characteristics:} 2 points
    \item \textbf{Total Points Per Question:} 6 points
\end{enumerate}

%%===============================================================
\section{Dividing Polynomials}
\label{sec:dividing-polynomials}
%%===============================================================

\paragraph{Polynomial Long Division}

Suppose we would like to divide one polynomial by another. The procedure is similar to long division of numbers and is illustrated in the following example:

\paragraph{Example 1}

Divide $x^2-2x-15$ (the dividend or numerator) by \\(x+3\\) (the divisor or denominator). Similar to long division of numbers, we set up our problem as follows:

$$
\begin{array}{rl}
         &       \text{\textunderscore\textunderscore\textunderscore\textunderscore\textunderscore\textunderscore\textunderscore\textunderscore\textunderscore\textunderscore\textunderscore\textunderscore}  \\\\
   x + 3 & \big) x^2-2x-15
\end{array}
$$

First, we have to answer the question, how many times does \(x+3\) go into \(x^2\)? To find out, divide the leading term of the dividend by the leading term of the divisor. So it goes in \(x\) times. We record this above the leading term of the dividend:

$$
\begin{array}{rl}
         &       x   \\\\
         &       \text{\textunderscore\textunderscore\textunderscore\textunderscore\textunderscore\textunderscore\textunderscore\textunderscore\textunderscore\textunderscore\textunderscore\textunderscore}  \\\\
   x + 3 & \big) x^2-2x-15
\end{array}
$$

Next, we multiply \(x+3\) by \(x\) and write this below the dividend as follows:
$$
\begin{array}{rl}
         &       x   \\\\
         &       \text{\textunderscore\textunderscore\textunderscore\textunderscore\textunderscore\textunderscore\textunderscore\textunderscore\textunderscore\textunderscore\textunderscore\textunderscore}  \\\\
   x + 3 & \big) x^2-2x-15 \\\\
         & -(x^2+3x)   \\\\
         &       \text{\textunderscore\textunderscore\textunderscore\textunderscore\textunderscore\textunderscore\textunderscore\textunderscore\textunderscore\textunderscore\textunderscore\textunderscore}  \\\\
\end{array}
$$

Now we perform the subtraction, bringing down any terms in the dividend that aren''''''''t matched in our subtrahend:
$$
\begin{array}{rl}
         &       x   \\\\
         &       \text{\textunderscore\textunderscore\textunderscore\textunderscore\textunderscore\textunderscore\textunderscore\textunderscore\textunderscore\textunderscore\textunderscore\textunderscore}  \\\\
   x + 3 & \big) x^2-2x-15 \\\\
         & -(x^2+3x)   \\\\
         &       \text{\textunderscore\textunderscore\textunderscore\textunderscore\textunderscore\textunderscore\textunderscore\textunderscore\textunderscore\textunderscore\textunderscore\textunderscore}  \\\\
         &  -5x-15
\end{array}
$$

Now we repeat, treating the bottom line as our new dividend:
$$
\begin{array}{rl}
         &       x - 5  \\\\
         &       \text{\textunderscore\textunderscore\textunderscore\textunderscore\textunderscore\textunderscore\textunderscore\textunderscore\textunderscore\textunderscore\textunderscore\textunderscore}  \\\\
   x + 3 & \big) x^2-2x-15 \\\\
         & -(x^2+3x)   \\\\
         &       \text{\textunderscore\textunderscore\textunderscore\textunderscore\textunderscore\textunderscore\textunderscore\textunderscore\textunderscore\textunderscore\textunderscore\textunderscore}  \\\\
         &  -5x-15 \\\\
         &  -(-5x-15) \\\\
         &       \text{\textunderscore\textunderscore\textunderscore\textunderscore\textunderscore\textunderscore\textunderscore\textunderscore\textunderscore\textunderscore\textunderscore\textunderscore}  \\\\
         & \phantom{---}0
\end{array}
$$

In this case, we have no remainder.

\paragraph{Example 2}

What about non-divisible polynomials like this?

$$(3x^2 + 3x - 4) / (x - 4)$$

1. We first consider the highest-degree terms from both the dividend and divisor, the result is the first term of our quotient.
   $$\frac{3x^2}{x} = 3x$$
   $$
   \begin{array}{r|ccc}
   x - 4 & 3x^2 & 3x & -4 \\\\
   \end{array}
   $$

2. Then we multiply this by our divisor.
   $$(3x) \times (x - 4) = 3x^2 - 12x$$

3. And subtract the result from our dividend.
   $$(3x^2 + 3x - 4) - (3x^2 - 12x) = 15x - 4$$
   $$
   \begin{array}{r|ccc}
   x - 4 & 3x^2 & 3x & -4 \\\\
   3x(x - 4) & 3x^2 & -12x & \\\\
   \end{array}
   $$

4. Now once again with the highest-degree terms of the remaining polynomial, and we got the second term of our quotient.
   $$\frac{15x}{x} = 15$$

5. Multiplying...
   $$(15) \times (x - 4) = 15x - 60$$

6. Subtracting...
   $$(15x - 4) - (15x - 60) = 56$$
   $$
   \begin{array}{r|ccc}
   x - 4 & 3x^2 & 3x & -4 \\\\
   3x(x - 4) & 3x^2 & -12x & \\\\
   15(x - 4) & 15x & -60 \\\\
   \end{array}
   $$

7. We are left with a constant term - our remainder:
   $$Q(x) = 3x + 15, \quad R = 56$$

So finally:
$$(3x^2 + 3x - 4) = (3x + 15) \times (x - 4) + 56$$

\paragraph{Resources}

\paragraph{Dividing Polynomials With Long Division}

\begin{itemize}
    \item \href{https://www.khanacademy.org/math/algebra2/x2ec2f6f830c9fb89:poly-div/x2ec2f6f830c9fb89:quad-div-by-linear/v/polynomial-division}{Intro to Polynomial Long Division}, Khan Academy
    \item \href{https://courses.lumenlearning.com/waymakercollegealgebra/chapter/polynomial-long-division/}{Polynomial Long Division}, Lumen Learning
    \item \href{https://www.purplemath.com/modules/polydiv2.htm}{Polynomial Long Division}, Purple Math
    \item \href{https://www.youtube.com/watch?v=_FSXJmESFmQ}{Long Division With Polynomials}, The Organic Chemistry Tutor
    \item \href{https://www.youtube.com/watch?v=l6_ghhd7kwQ}{Long Division of Polynomials}, patrickJMT
\end{itemize}

\paragraph{Dividing Polynomials with Synthetic Division}

\begin{itemize}
    \item \href{https://www.khanacademy.org/math/algebra-home/alg-polynomials/alg-synthetic-division-of-polynomials/v/synthetic-division}{Synthetic Division of Polynomials}, Khan Academy
    \item \href{https://www.purplemath.com/modules/synthdiv.htm}{Synthetic Division}, Purple Math
    \item \href{https://www.wtamu.edu/academic/anns/mps/math/mathlab/col_algebra/col_alg_tut37_syndiv.htm}{Synthetic Division}, WTAMU VirtualMathLab
    \item \href{https://www.youtube.com/watch?v=bZoMz1Cy1T4}{Synthetic Division Example}, patrickJMT
\end{itemize}

\paragraph{Licensing}

Content obtained and/or adapted from:
\begin{itemize}
    \item \href{https://en.wikibooks.org/wiki/Calculus/Algebra}{Algebra, Wikibooks: Calculus} under a CC BY-SA license
    \item \href{https://en.wikibooks.org/wiki/High_School_Mathematics_Extensions/Supplementary/Polynomial_Division}{Polynomial Division, Wikibooks: High School Mathematics Extensions} under a CC BY-SA license
\end{itemize}

\subsection{Learning Objectives}

\begin{enumerate}
    \item SLO 1: Understand the Concept of Polynomial Long Division
    \item Students will be able to explain the steps involved in polynomial long division, including how to set up the division problem and identify the divisor and dividend. SLO 2: Perform Polynomial Long Division
    \item Students will be able to divide a given polynomial by another polynomial using the long division method, accurately performing each step to find the quotient and remainder. SLO 3: Apply Polynomial Long Division to Various Problems
    \item Students will be able to apply the polynomial long division technique to solve problems involving both divisible and non-divisible polynomials, demonstrating their ability to handle different scenarios.
\end{enumerate}

\subsection{Assessment Instrument}

\begin{enumerate}
    \item Instructions Solve the following problems using polynomial long division. Show all your work clearly. Problem 1 Divide $x^2 - 2x - 15$ by $x + 3$. Problem 2 Divide $\frac{3x^2 + 3x - 4}{x - 4}$ Problem 3 (Challenge problem) Divide $\frac{5x^3 - 2x + 3}{x + 2}$
\end{enumerate}

%%===============================================================
\section{Remainder and Factor Theorem}
\label{sec:remainder-and-factor-theorem}
%%===============================================================

\paragraph{Remainder Theorem}

In algebra, the polynomial remainder theorem or little Bézout''s theorem (named after Étienne Bézout) is an application of Euclidean division of polynomials. It states that the remainder of the division of a polynomial $f(x)$ by a linear polynomial $x - r$ is equal to $f(r)$. In particular, $x - r$ is a divisor of $f(x)$ if and only if $f(r) = 0$, a property known as the factor theorem.

\paragraph{Examples}

\paragraph{Example 1}

Let $f(x) = x^3 - 12x^2 - 42$. Polynomial division of $f(x)$ by $(x - 3)$ gives the quotient $x^2 - 9x - 27$ and the remainder $-123$. Therefore, $f(3) = -123$.

\paragraph{Example 2}

Show that the polynomial remainder theorem holds for an arbitrary second degree polynomial $f(x) = ax^2 + bx + c$ by using algebraic manipulation:

$$
\frac{f(x)}{x - r} = \frac{ax^2 + bx + c}{x - r} = \frac{ax^2 - arx + arx + bx + c}{x - r} = \frac{ax(x - r) + (b + ar)x + c}{x - r} = ax + \frac{(b + ar)(x - r) + c + r(b + ar)}{x - r} = ax + b + ar + \frac{c + r(b + ar)}{x - r} = ax + b + ar + \frac{ar^2 + br + c}{x - r}
$$

Multiplying both sides by $(x - r)$ gives

$$
f(x) = ax^2 + bx + c = (ax + b + ar)(x - r) + ar^2 + br + c
$$

Since $R = ar^2 + br + c$ is the remainder, we have indeed shown that $f(r) = R$.

\paragraph{Proof}

The polynomial remainder theorem follows from the theorem of Euclidean division, which, given two polynomials $f(x)$ (the dividend) and $g(x)$ (the divisor), asserts the existence (and the uniqueness) of a quotient $Q(x)$ and a remainder $R(x)$ such that

$$
f(x) = Q(x)g(x) + R(x) \quad \text{and} \quad R(x) = 0 \ \text{ or } \deg(R) < \deg(g).
$$

If the divisor is $g(x) = x - r$, where $r$ is a constant, then either $R(x) = 0$ or its degree is zero; in both cases, $R(x)$ is a constant that is independent of $x$; that is

$$
f(x) = Q(x)(x - r) + R.
$$

Setting $x = r$ in this formula, we obtain:

$$
f(r) = R.
$$

A slightly different proof, which may appear to some people as more elementary, starts with an observation that $f(x) - f(r)$ is a linear combination of terms of the form $x^k - r^k$, each of which is divisible by $x - r$ since $x^k - r^k = (x - r)(x^{k-1} + x^{k-2}r + \dots + xr^{k-2} + r^{k-1})$.

\paragraph{Applications}

The polynomial remainder theorem may be used to evaluate $f(r)$ by calculating the remainder, $R$. Although polynomial long division is more difficult than evaluating the function itself, synthetic division is computationally easier. Thus, the function may be more "cheaply" evaluated using synthetic division and the polynomial remainder theorem.

The factor theorem is another application of the remainder theorem: if the remainder is zero, then the linear divisor is a factor. Repeated application of the factor theorem may be used to factorize the polynomial.

\paragraph{Factor Theorem}

In algebra, the factor theorem is a theorem linking factors and zeros of a polynomial. It is a special case of the polynomial remainder theorem.

The factor theorem states that a polynomial $f(x)$ has a factor $(x - k)$ if and only if $f(k) = 0$ (i.e. $k$ is a root).

\paragraph{Factorization of Polynomials}

Two problems where the factor theorem is commonly applied are those of factoring a polynomial and finding the roots of a polynomial equation; it is a direct consequence of the theorem that these problems are essentially equivalent.

The factor theorem is also used to remove known zeros from a polynomial while leaving all unknown zeros intact, thus producing a lower degree polynomial whose zeros may be easier to find. Abstractly, the method is as follows:

1. "Guess" a zero $a$ of the polynomial $f$. (In general, this can be very hard, but math textbook problems that involve solving a polynomial equation are often designed so that some roots are easy to discover.)
2. Use the factor theorem to conclude that $(x - a)$ is a factor of $f(x)$.
3. Compute the polynomial $g(x) = \frac{f(x)}{(x - a)}$, for example using polynomial long division or synthetic division.
4. Conclude that any root $x \neq a$ of $f(x) = 0$ is a root of $g(x) = 0$. Since the polynomial degree of $g$ is one less than that of $f$, it is "simpler" to find the remaining zeros by studying $g$.

\paragraph{Example}

Find the factors of

$$
x^3 + 7x^2 + 8x + 2.
$$

To do this one would use trial and error (or the rational root theorem) to find the first $x$ value that causes the expression to equal zero. To find out if $(x - 1)$ is a factor, substitute $x = 1$ into the polynomial above:

$$
x^3 + 7x^2 + 8x + 2 = (1)^3 + 7(1)^2 + 8(1) + 2 = 1 + 7 + 8 + 2 = 18
$$

As this is equal to 18 and not 0, this means $(x - 1)$ is not a factor of $x^3 + 7x^2 + 8x + 2$. So, we next try $(x + 1)$ (substituting $x = -1$ into the polynomial):

$$
(-1)^3 + 7(-1)^2 + 8(-1) + 2.
$$

This is equal to 0. Therefore $x - (-1)$, which is to say $x + 1$, is a factor, and $-1$ is a root of $x^3 + 7x^2 + 8x + 2$.

The next two roots can be found by algebraically dividing $x^3 + 7x^2 + 8x + 2$ by $(x + 1)$ to get a quadratic:

$$
\frac{x^3 + 7x^2 + 8x + 2}{x + 1} = x^2 + 6x + 2,
$$

and therefore $(x + 1)$ and $x^2 + 6x + 2$ are factors of $x^3 + 7x^2 + 8x + 2$. Of these, the quadratic factor can be further factored using the quadratic formula, which gives as roots of the quadratic $-3 \pm \sqrt{7}$. Thus the three irreducible factors of the original polynomial are $x + 1$, $x - (-3 + \sqrt{7})$, and $x - (-3 - \sqrt{7})$.

\paragraph{Resources}

\begin{itemize}
    \item \href{https://tutorial.math.lamar.edu/classes/alg/dividingpolynomials.aspx}{Dividing Polynomials, Paul''s Online Notes}
\end{itemize}

\paragraph{Licensing}

Content obtained and/or adapted from:

\begin{itemize}
    \item \href{https://en.wikipedia.org/wiki/Polynomial_remainder_theorem}{Polynomial remainder theorem, Wikipedia under a CC BY-SA license}
    \item \href{https://en.wikipedia.org/wiki/Factor_theorem}{Factor theorem, Wikipedia under a CC BY-SA license}
\end{itemize}

\subsection{Learning Objectives}

\begin{enumerate}
    \item 1.\textbf{Identify and apply the Remainder Theorem}: Students will be able to determine the remainder of a polynomial $f(x)$ when divided by a linear polynomial $x - r$ and evaluate $f(r)$. 2.\textbf{Prove the Remainder Theorem}: Students will demonstrate an understanding of the Remainder Theorem by proving it using algebraic manipulation and Euclidean division of polynomials. 3.\textbf{Use the Factor Theorem to factorize polynomials}: Students will use the Factor Theorem to find factors and roots of polynomials, simplifying polynomial expressions by identifying and removing known zeros.
\end{enumerate}

\subsection{Assessment Instrument}

\begin{enumerate}
    \item Instructions Complete the following problems. Show all your work and provide explanations where necessary. Problem 1: Basic Application of the Remainder Theorem Evaluate $f(2)$ using the Remainder Theorem for the polynomial $f(x) = 3x^3 - 4x^2 + 5x - 6$. Problem 2: Proving the Remainder Theorem Prove the Remainder Theorem for the polynomial $f(x) = 2x^2 - 3x + 1$ when divided by $x - 1$ using algebraic manipulation. Problem 3: Factor Theorem and Polynomial Factorization Given the polynomial $f(x) = x^3 - 6x^2 + 11x - 6$, use the Factor Theorem to:
    \item Find a factor of $f(x)$.
    \item Factorize $f(x)$ completely. Problem 4: Application of Synthetic Division Evaluate $f(3)$ using synthetic division for the polynomial $f(x) = 2x^4 - 3x^3 + 4x^2 - 5x + 6$. Problem 5: Verifying the Factor Theorem Verify that $x + 2$ is a factor of the polynomial $f(x) = x^3 + x^2 - 4x - 4$ and find all the factors of $f(x)$. Problem 6: General Proof of the Remainder Theorem Provide a general proof of the Remainder Theorem using Euclidean division for an arbitrary polynomial $f(x)$ and a linear divisor $x - r$. --- \textbf{Note}: Be sure to check your work for accuracy and clarity. If needed, review the polynomial division and synthetic division methods to ensure correct solutions.
\end{enumerate}

%%===============================================================
\section{Systems of Equations in Two Variables}
\label{sec:systems-of-equations-in-two-variables}
%%===============================================================

\paragraph{Systems of Equations in Two Variables}

\paragraph{Introduction}

A system of equations consists of two or more equations made up of two or more variables such that all equations in the system are considered simultaneously. To find the unique solution to a system of equations, we must find a numerical value for each variable in the system that will satisfy all equations in the system at the same time. Some systems may not have a solution (for example, two distinct lines that are parallel to one another) and others may have an infinite number of solutions (for example, $y = 0$ and $y = \sin(x)$, since $\sin(x) = 0$ for an infinite number of $x$-values). In order for a linear system to have a unique solution, there must be at least as many equations as there are variables. Even so, this does not guarantee a unique solution.

Systems of equations can be solved in a number of ways. We can use elimination and/or substitution to solve for points of intersection shared by all equations in the system. We can also sometimes graph the equations to find the points shared by all equations, given that our system is graphable.

Here are some examples of systems of equations with two variables:

\begin{itemize}
    \item \textbf{One solution}: $y = 2x$ and $y = x + 5$. By use of substitution, elimination, or graphing, we can find that the solution to this system is $(5, 10)$, as this is the only point shared by both equations in the system.
\end{itemize}

\begin{itemize}
    \item \textbf{Multiple solutions}: $y = x^2$ and $y = 4$. Since $x^2 = 4$ when $x = 2$ or $-2$, we know there are two solutions to this system: $(2, 4)$ and $(-2, 4)$. This is also clear when we graph these two equations, as $y = 4$ intersects the parabola $y = x^2$ when $x = -2$, and when $x = 2$.
\end{itemize}

\begin{itemize}
    \item \textbf{No solutions}: $y = x^2$ and $y = -3$. Since $x^2$ is positive for all values of $x$, it cannot equal $-3$. When graphed, we can see that these two equations never intersect. Thus there are no solutions to this system of equations.
\end{itemize}

\begin{itemize}
    \item \textbf{Infinite solutions}: $y = \sin(x)$ and $y = 0$. $\sin(x) = 0$ for all $x = \pi k$, where $k$ is some integer. So, this system has an infinite number of solutions of the form $(\pi k, 0)$; that is, the solution set for this system of equations is {..... $( -2\pi , 0 )$, $( -\pi , 0 )$, $( 0 , 0 )$, $( \pi , 0 )$, $( 2\pi , 0 )$ .....}.
\end{itemize}

\paragraph{Linear Equations with Two Variables}

In the previous module, linear equations with two variables were discussed. A single linear equation having two unknown variables is practically insufficient to solve or even narrow down the solutions for the two variables, although it does establish a relationship between them. The relationship is shown graphically as a line. Another linear equation with the same two variables may be enough to narrow down the solution to the two equations to one value for the first variable and one value for the second variable, i.e., to solve the system of two simultaneous linear equations. Let''s see how two linear equations with the same two unknowns might be related to each other. Since we said it was given that both equations were linear, the graphs of both equations would be lines in the same two-dimensional coordinate plane (for a system with two variables). The lines could be related to each other in the following three ways:

1. The graphs of both equations could coincide giving the same line. This means that the two equations are providing the same information about how the variables are related to each other. The two equations are basically the same, perhaps just different versions or forms of each other. Either one could be mathematically manipulated to produce the other one. Both lines would have the same slope and the same y-intercept. Such equations are considered dependent on each other. Since no new information is provided, the addition of the second equation does not solve the problem by narrowing the solution set down to one solution.

   \textbf{Example: Dependent linear equations}

   $6x - 3y = 12$

   $y = 2x - 4$

   The above two equations provide the same information and result in the same graph, i.e., lines which coincide as shown in the following image.

% [Image: graph]

Let''s see how these equations can be mathematically manipulated to show they are basically the same.

Divide both sides of the first equation $6x - 3y = 12$ by 3 to give

$$
2x - y = 4
$$

Now add $y$ to both sides

$$
2x = 4 + y
$$

Now subtract 4 from both sides

$$
y = 2x - 4
$$

This is the same as the second equation in the example. This is the slope-intercept form of the equation, from which a slope and a $y$-intercept unique to the line can be compared with any other equations in the slope-intercept form.

2. The graphs of two lines could be parallel although not the same. The two lines do not intersect each other at any point. This means there is no solution which satisfies both equations simultaneously, i.e., at the same time. The solution set for this system of simultaneous linear equations is the empty set. Such equations are considered inconsistent with each other and actually give contradictory information if it is claimed they are both true at the same time in the same problem. The parallel lines have equal slopes but different $y$-intercepts.

Sets of equations which have at least one common point which might provide a solution set are consistent with each other. For example, the dependent equations mentioned previously are consistent with each other.

\textbf{Example: Inconsistent linear equations}

$$
3x - 2y = -2
$$

$$
3x - 2y = 2
$$

To compare slopes and $y$-intercepts for these two linear equations, we place them in the slope-intercept forms. Subtract $3x$ from both sides of both equations.

$$
3x - 2y = -2
$$

$$
3x - 2y = 2
$$

$$
-2y = -3x - 2
$$

$$
-2y = -3x + 2
$$

Divide both sides of both equations by $-2$ and simplify to get slope-intercept forms for comparison.

$$
\frac{-2y}{-2} = \frac{-3x - 2}{-2}
$$

$$
\frac{-2y}{-2} = \frac{-3x + 2}{-2}
$$

$$
y = \frac{3}{2}x + 1
$$

$$
y = \frac{3}{2}x - 1
$$

Now, both slopes are equal at $\frac{3}{2}$, but the $y$-intercepts at 1 and -1 are different. The lines are parallel. The graphs are shown here:

% [Image: graph]

3. If the two lines are not the same and are not parallel, then they would intersect at one point because they are graphed in the same two-dimensional coordinate plane. The one point of intersection is the ordered pair of numbers which is the solution to the system of two linear equations and two unknowns. The two equations provide enough information to solve the problem and further equations are not needed. Such equations intersecting at a point providing a solution to the problem are considered independent of each other. The lines have different slopes but may or may not have the same $y$-intercept. Because such equations provide at least one solution point, they are consistent with each other.

\textbf{Example: Consistent independent linear equations}

$$
y = 3x - 5
$$

$$
y = -x - 1
$$

Both of these equations are given in the slope-intercept form, so it is easy to compare slopes and $y$-intercepts. For these two linear functions, both slopes are different and both $y$-intercepts are different. This means the lines are neither dependent nor inconsistent, so on a two-dimensional graph they must intersect at some point. In fact, the graph shows the lines intersecting at $(1, -2)$, which is the ordered pair solution to this system of independent simultaneous equations. Visual inspection of a graph cannot be relied on to give perfectly accurate coordinates every time, so either the point is tested with both equations or one of the following two methods is used to determine accurate coordinates for the intersection point.

% [Image: graph]

\paragraph{Resources}

\begin{itemize}
    \item \href{https://tutorial.math.lamar.edu/classes/alg/systemstwovrble.aspx}{Linear Systems with Two Variables, Paul''s Online Notes}
    \item \href{https://openstax.org/books/college-algebra/pages/7-1-systems-of-linear-equations-two-variables}{Systems of Linear Equations with Two Variables, OpenStax}
    \item \href{https://openstax.org/books/college-algebra/pages/7-3-systems-of-nonlinear-equations-and-inequalities-two-variables}{Systems of Nonlinear Equations with Two Variables, OpenStax}
    \item \href{https://www.youtube.com/watch?v=ej25myhYcSg}{Solving Systems of Equations with Elimination, patrickJMT}
    \item \href{https://www.youtube.com/watch?v=cwHR_B9zK7k}{Solving Systems of Equations with Substitution, patrickJMT}
    \item \href{https://www.youtube.com/watch?v=WNkPKv0OTGI}{Solving Systems of Equations with Graphing, patrickJMT}
\end{itemize}

\paragraph{Licensing}

Content obtained and/or adapted from:

\begin{itemize}
    \item \href{https://openstax.org/books/college-algebra/pages/7-1-systems-of-linear-equations-two-variables}{Systems of Linear Equations: Two Variables, OpenStax: College Algebra under a CC BY license}
    \item \href{https://en.wikibooks.org/wiki/Algebra/Systems_of_Equations}{Systems of Equations, Wikibooks: Algebra under a CC BY-SA license}
\end{itemize}

\subsection{Learning Objectives}

\begin{enumerate}
    \item \textbf{Identify and Solve Systems with Unique Solutions}: Solve a system of two linear equations to find the unique point of intersection. Example: Given $y = 2x$ and $y = x + 5$, find the solution $(5, 10)$.
    \item \textbf{Determine Consistency and Parallelism}: Analyze systems of linear equations to determine if they are consistent, dependent, or inconsistent. For example, identify that $3x - 2y = -2$ and $3x - 2y = 2$ are inconsistent due to parallel lines with the same slope but different $y$-intercepts.
    \item \textbf{Graph and Interpret Systems with Multiple or Infinite Solutions}: Graph and interpret systems where equations intersect at multiple points or have infinite solutions. For example, recognize that $y = x^2$ and $y = 4$ intersect at $(2, 4)$ and $(-2, 4)$, showing multiple solutions.
\end{enumerate}

\subsection{Assessment Instrument}

\begin{enumerate}
    \item Instructions Solve the following problems related to systems of equations in two variables. Provide your answers in the specified format and show all your work where applicable. Questions 1. Solve the System of Equations Solve the following system of linear equations and find the point of intersection. $y = 3x + 2$ and $y = -2x + 4$ 2. Determine the Nature of the System For each pair of equations below, determine whether the system is consistent and independent, consistent and dependent, or inconsistent. Provide a brief explanation for your answer.
    \item $4x - 5y = 10$ and  $8x - 10y = 20$
    \item $x + y = 3$ and $x - y = 1$ 3. Graphical Interpretation Graph the following systems of equations and determine if they have one solution, no solution, multiple solutions, or infinitely many solutions.
    \item $y = 3x - 8$ and $6y - 18x = -48$
    \item $y = \sin(x)$ and $y = 0$ 4. Solve Using Substitution or Elimination Solve the following system using either the substitution method or the elimination method: $2x - y = 1$ and $3x + 2y = 10$ 5. Application Problem A farmer has two types of crops to plant in a field. The first type requires 3 hours of labor per acre and yields `$`200 per acre. The second type requires 2 hours of labor per acre and yields `$`150 per acre. The farmer has a total of 30 hours of labor available and wants to maximize the revenue from the crops. Set up a system of linear equations to represent the situation and solve for the number of acres of each crop the farmer should plant to maximize revenue.
\end{enumerate}

%%===============================================================
\section{Systems of Inequalities}
\label{sec:systems-of-inequalities}
%%===============================================================

\paragraph{Systems of Inequalities}

In mathematics, a linear inequality is an inequality that involves a linear function. A linear inequality contains one of the symbols of inequality: It shows the data which is not equal in graph form.

\begin{itemize}
    \item $<$ less than
    \item $>$ greater than
    \item $\leq$ less than or equal to
    \item $\geq$ greater than or equal to
    \item $\neq$ not equal to
    \item $=$ equal to
\end{itemize}

A linear inequality looks exactly like a linear equation, with the inequality sign replacing the equality sign.

\paragraph{Contents}

\paragraph{Linear inequalities of real numbers}

\paragraph{Two-dimensional linear inequalities}

Two-dimensional linear inequalities are expressions in two variables of the form:

$$
a x + b y < c \text{ and } a x + b y \geq c
$$

where the inequalities may either be strict or not. The solution set of such an inequality can be graphically represented by a half-plane (all the points on one "side" of a fixed line) in the Euclidean plane. Technically, for this statement to be correct both $a$ and $b$ cannot simultaneously be zero. In that situation, the solution set is either empty or the entire plane. The line that determines the half-planes ($a x + b y = c$) is not included in the solution set when the inequality is strict. A simple procedure to determine which half-plane is in the solution set is to calculate the value of $a x + b y$ at a point $(x\_0, y\_0)$ which is not on the line and observe whether or not the inequality is satisfied.

For example, to draw the solution set of $x + 3y < 9$, one first draws the line with equation $x + 3y = 9$ as a dotted line, to indicate that the line is not included in the solution set since the inequality is strict. Then, pick a convenient point not on the line, such as $(0,0)$. Since $0 + 3(0) = 0 < 9$, this point is in the solution set, so the half-plane containing this point (the half-plane "below" the line) is the solution set of this linear inequality.

Graph of linear inequality: $x + 3y < 9$

% [Image: graph]

\paragraph{Linear inequalities in general dimensions}

In $\mathbb{R}^n$, linear inequalities are the expressions that may be written in the form:

$$
f(\bar{x}) < b \text{ or } f(\bar{x}) \leq b
$$

where $f$ is a linear form (also called a linear functional), $\bar{x} = (x\_1, x\_2, \ldots, x\_n)$, and $b$ is a constant real number.

More concretely, this may be written out as:

$$
a\_1 x\_1 + a\_2 x\_2 + \cdots + a\_n x\_n < b
$$

or

$$
a\_1 x\_1 + a\_2 x\_2 + \cdots + a\_n x\_n \leq b
$$

Here $x\_1, x\_2, \ldots, x\_n$ are called the unknowns, and $a\_1, a\_2, \ldots, a\_n$ are called the coefficients.

Alternatively, these may be written as:

$$
g(x) < 0
$$

or

$$
g(x) \leq 0
$$

where $g$ is an affine function.

That is:

$$
a\_0 + a\_1 x\_1 + a\_2 x\_2 + \cdots + a\_n x\_n < 0
$$

or

$$
a\_0 + a\_1 x\_1 + a\_2 x\_2 + \cdots + a\_n x\_n \leq 0
$$

Note that any inequality containing a "greater than" or a "greater than or equal" sign can be rewritten with a "less than" or "less than or equal" sign, so there is no need to define linear inequalities using those signs.

\paragraph{Systems of linear inequalities}

A system of linear inequalities is a set of linear inequalities in the same variables:

$ a\_{11} x\_1 + a\_{12} x\_2  + \cdots + a\_{1n} x\_n \leq b\_1 $

$ a\_{21} x\_1 + a\_{22} x\_2 + \cdots + a\_{2n} x\_n \leq b\_2 $ 

$             \vdots                         \vdots                          \vdots                          \vdots                          \vdots $

$ a\_{m1} x\_1 + a\_{m2} x\_2 + \cdots + a\_{mn} x\_n \leq b\_m $ 

Here $x\_1, x\_2, \ldots, x\_n$ are the unknowns, $a\_{11}, a\_{12}, \ldots, a\_{mn}$ are the coefficients of the system, and $b\_1, b\_2, \ldots, b\_m$ are the constant terms.

This can be concisely written as the matrix inequality:

$$
Ax \leq b
$$

where $A$ is an $m \times n$ matrix, $x$ is an $n \times 1$ column vector of variables, and $b$ is an $m \times 1$ column vector of constants.

In the above systems, both strict and non-strict inequalities may be used.

Not all systems of linear inequalities have solutions. Variables can be eliminated from systems of linear inequalities using Fourier–Motzkin elimination.

\paragraph{Applications}

\paragraph{Polyhedra}

The set of solutions of a real linear inequality constitutes a half-space of the $n$-dimensional real space, one of the two defined by the corresponding linear equation.

The set of solutions of a system of linear inequalities corresponds to the intersection of the half-spaces defined by individual inequalities. It is a convex set, since the half-spaces are convex sets, and the intersection of a set of convex sets is also convex. In the non-degenerate cases, this convex set is a convex polyhedron (possibly unbounded, e.g., a half-space, a slab between two parallel half-spaces, or a polyhedral cone). It may also be empty or a convex polyhedron of lower dimension confined to an affine subspace of the $n$-dimensional space $\mathbb{R}^n$.

\paragraph{Linear programming}

A linear programming problem seeks to optimize (find a maximum or minimum value) a function (called the objective function) subject to a number of constraints on the variables which, in general, are linear inequalities. The list of constraints is a system of linear inequalities.

\paragraph{Resources}

\begin{itemize}
    \item \href{https://www.khanacademy.org/math/algebra/x2f8bb11595b61c86:inequalities-systems-graphs/x2f8bb11595b61c86:graphing-two-variable-inequalities/v/graphing-systems-of-inequalities-2}{Graphing Systems of Inequalities, Khan Academy}
    \item \href{https://www.youtube.com/watch?v=5xQqwgS3O4U}{System of Linear Inequalities, Khan Academy}
\end{itemize}

\paragraph{Licensing}

Content obtained and/or adapted from:

\begin{itemize}
    \item \href{https://en.wikipedia.org/wiki/Linear_inequality}{Linear inequality, Wikipedia} under a CC BY-SA license
\end{itemize}

\subsection{Learning Objectives}

\begin{enumerate}
    \item \textbf{SLO 1:} Students will be able to define and identify linear inequalities and their symbols ($<$, $>$, $\leq$, $\geq$, $\neq$, $=$).
    \item \textbf{SLO 2:} Students will be able to graphically represent two-dimensional linear inequalities and determine the solution set of such inequalities.
    \item \textbf{SLO 3:} Students will be able to solve systems of linear inequalities and express the solution sets using matrix notation.
\end{enumerate}

\subsection{Assessment Instrument}

\begin{enumerate}
    \item Multiple-Choice Questions
    \item \textbf{Question 1:} Which of the following represents a linear inequality?
    \item A) $2x + 3 = 7$
    \item B) $4y - 5 > 2$
    \item C) $x^2 + y < 5$
    \item D) $7x \geq 4x - 9$
    \item \textbf{Question 2:} What is the solution set for the inequality $x + 3y < 9$ in the coordinate plane?
    \item A) The line $x + 3y = 9$
    \item B) The half-plane below the line $x + 3y = 9$
    \item C) The half-plane above the line $x + 3y = 9$
    \item D) The entire coordinate plane Short Answer Questions
    \item \textbf{Question 3:} Explain how you can determine which half-plane is included in the solution set of a linear inequality.
    \item \textbf{Question 4:} Write the system of linear inequalities in matrix form: $$ 2x + 3y \leq 5 $$ $$ -x + 4y > 3$$ Problem Solving
    \item \textbf{Question 5:} Solve the system of linear inequalities and graph the solution set: $$ x + y \leq 4 $$ $$ x - y \geq 1 $$
\end{enumerate}

%%===============================================================
\section{Exponential Functions}
\label{sec:exponential-functions}
%%===============================================================

\paragraph{Exponential Functions}

\paragraph{Operations With Exponential Functions}

\paragraph{Multiplication}

Firstly,
$$
b^x \times b^{2x} = b^{(x + 2x)} = b^{3x}
$$
So, when you multiply a base by the same base, you add the exponents. To clarify, here is an example with numbers:

$$
\begin{array}{|c|c|c|c|c|} \hline
x & 2^{x} & 2^{2x} & 2^{x} \times 2^{2x} & 2^{3x} \\\\ 
\hline
1 & 2 & 4 & 8 & 8 \\\\ 
\hline
2 & 4 & 16 & 64 & 64 \\\\ 
\hline
\end{array}
$$

\paragraph{Division}

Secondly,
$$
\frac{b^{2x}}{b^y} = b^{(2x - y)}
$$
So, when dividing a base by the same base, you subtract the exponents.

\textbf{Example with numbers:}

$$
\frac{2^4}{2^2} = \frac{16}{4} = 4 = 2^2
$$

\paragraph{Base Raised to Two Powers}

Thirdly,
$$
\left(b^{2x}\right)^{3x} = b^{(2x \times 3x)} = b^{6x^2}
$$
So, when a base raised to a variable is itself raised to another variable, you multiply the exponents.

\textbf{Example with numbers (when $x = 1$):}

$$
\left(2^2\right)^3 = 4^3 = 64 = 2^6
$$

\paragraph{Multiple Bases}

Fourthly,
$$
a^2 \times b^2 = (ab)^2
$$
Here, multiplying two different bases raised to the same power is equivalent to raising their product to the same power. Similarly, for division:

$$
\left(\frac{a}{b}\right)^2 = \frac{a^2}{b^2}
$$

\textbf{Example with numbers:}

$$
2^2 \times 3^2 = 36 = 6^2
$$

\paragraph{Fractional Exponents}

When $x$ is a fraction, you can express it as a root function. For example,
$$
b^{\frac{1}{x}} = \sqrt[x]{b}
$$
Similarly, when the fraction has a constant (denoted as $c$) other than 1 in the numerator,
$$
b^{\frac{c}{x}} = \left(\sqrt[x]{b}\right)^c
$$

\paragraph{The Laws of Exponents}

The rules for exponents are known as the laws of exponents and can be summarized as:

$$
b^x \times b^y = b^{x+y}
$$

$$
\frac{b^x}{b^y} = b^{x-y}
$$

$$
\left(b^x\right)^y = b^{xy}
$$

$$
a^n \times b^n = (ab)^n
$$

$$
\left(\frac{a}{b}\right)^n = \frac{a^n}{b^n}
$$

$$
b^{-n} = \frac{1}{b^n}
$$

$$
b^{\frac{c}{x}} = \left(\sqrt[x]{b}\right)^c \text{ where } c \text{ is a constant}
$$

$$
b^1 = b
$$

$$
b^0 = 1
$$

\paragraph{Graphing an Exponential Function}

When graphing an exponential function, use the same methods as with other functions. 

% [Image: Exponential Graph]

\paragraph{Solving Exponential Equations}

To solve an exponential equation, ensure all bases are the same. Then, you can remove the base and solve for the variable. 

\textbf{Example:}

Solve for $x$:
$$
2^{(x-1)} = 16
$$

Convert 16 to a base 2 raised to a number:
$$
2^{(x-1)} = 2^4
$$

Remove the base:
$$
x - 1 = 4
$$

Finally, solve for $x$:
$$
x = 5
$$

\paragraph{Resources}

\begin{itemize}
    \item \href{https://mathresearch.utsa.edu/wikiFiles/MAT1053/Exponential\%20Functions/MAT1053_M5.1Exponential_Functions.pdf}{Exponential Functions, Book Chapter}
    \item \href{https://mathresearch.utsa.edu/wikiFiles/MAT1053/Exponential\%20Functions/MAT1053_M5.1Exponential_FunctionsGN.pdf}{Guided Notes}
\end{itemize}

\paragraph{Licensing}

Content obtained and/or adapted from:

\begin{itemize}
    \item \href{https://en.wikibooks.org/wiki/A-level_Mathematics/OCR/C2/Logarithms_and_Exponentials}{Logarithms and Exponentials, Wikibooks: A-level Mathematics/OCR/C2 under a CC BY-SA license}
\end{itemize}

\subsection{Learning Objectives}

\begin{enumerate}
    \item 1.\textbf{Understand and Apply Exponential Multiplication}: Students will be able to multiply exponential functions with the same base and simplify the resulting expression by adding the exponents, as shown by the equation $b^x \times b^{2x} = b^{3x}$. 2.\textbf{Solve Exponential Division Problems}: Students will be able to divide exponential functions with the same base and simplify the result by subtracting the exponents, exemplified by $\frac{b^{2x}}{b^y} = b^{2x - y}$. 3.\textbf{Utilize Laws of Exponents in Complex Expressions}: Students will be able to apply the laws of exponents to simplify expressions involving multiple bases, fractional exponents, and raising a base to multiple powers, such as $\left(b^{2x}\right)^{3x} = b^{6x^2}$ and $b^{\frac{c}{x}} = \left(\sqrt[x]{b}\right)^c$.
\end{enumerate}

\subsection{Assessment Instrument}

\begin{enumerate}
    \item Assessment: Exponential Functions Instructions Complete the following questions to assess your understanding of exponential functions and their operations. Show all your work and simplify your answers. 1. Multiplication of Exponential Functions Given the expression: $$ b^3 \times b^5 $$ a. Simplify the expression using the properties of exponents. b. If $b = 2$, calculate the simplified result. 2. Division of Exponential Functions Evaluate the following expression: $$ \frac{b^{4x}}{b^{2x}} $$ a. Simplify the expression using the properties of exponents. b. If $b = 3$ and $x = 2$, calculate the result. 3. Base Raised to Two Powers Simplify the expression: $$ \left(b^x\right)^{2x} $$ a. Express the result in terms of a single exponent. b. If $b = 5$ and $x = 1$, compute the simplified result. 4. Multiple Bases Given the expression: $$ a^3 \times b^3 $$ a. Simplify the expression using the properties of exponents. b. If $a = 4$ and $b = 2$, calculate the simplified result. 5. Fractional Exponents Simplify the following expression: $$ b^{\frac{3}{2}} $$ a. Express the result as a root function. b. If $b = 8$, compute the simplified result. 6. Applying the Laws of Exponents Use the laws of exponents to simplify: $$ \frac{a^5 \times b^2}{a^2 \times b^3} $$ a. Express the result in terms of a single exponent. b. If $a = 2$ and $b = 4$, calculate the simplified result. 7. Graphing Exponential Functions Consider the exponential function: $$ f(x) = 2^x $$ a. Sketch the graph of $f(x)$ using standard graphing methods. b. Identify the y-intercept and the behavior as $x \to \infty$. 8. Solving Exponential Equations Solve for $x$ in the following equation: $$ 3^{x+1} = 81 $$ a. Express 81 as a power of 3. b. Solve for $x$.
\end{enumerate}

%%===============================================================
\section{Logarithmic Functions}
\label{sec:logarithmic-functions}
%%===============================================================

\paragraph{Logarithmic Functions}

% [Image: Logarithmic Functions]

\paragraph{Introduction}

In mathematics, you can find the inverse of an exponential function by switching $x$ and $y$ around: $y = b^x \text{ becomes } x = b^y$.  The problem arises in finding the value of $y$. The logarithmic function solves this problem. All conversions of a logarithmic function into an exponential function follow the same pattern:  $x = b^y \text{ becomes } y = \log\_b x$.  If a logarithm is given without a specified base, assume $b = 10$. Also, for logarithmic functions, $b > 0$ and $b \neq 1$. There are 2 cases where the log equals $x$:  $\log\_b b^X = X \text{ and } b^{\log\_b X} = X$.

To recap, a logarithm is the inverse function of an exponent. For example, the inverse of the function 

$$
f(x) = 3^x
$$

is 

$$
f^{-1}(x) = \log\_3 x.
$$

In general,

$$
y = b^x \iff x = \log\_b y
$$

given that $b > 0$.

\paragraph{Laws of Logarithmic Functions}

When $X$ and $Y$ are positive:

1. \textbf{Product Rule}:
   $$
   \log\_b (XY) = \log\_b X + \log\_b Y
   $$

2. \textbf{Quotient Rule}:
   $$
   \log\_b \left(\frac{X}{Y}\right) = \log\_b X - \log\_b Y
   $$

3. \textbf{Power Rule}:
   $$
   \log\_b (X^k) = k \log\_b X
   $$

\paragraph{Change of Base}

When $x$ and $b$ are positive real numbers and not equal to 1, you can write 

$$
\log\_a x
$$

as 

$$
\frac{\log\_b x}{\log\_b a}.
$$

This works for natural logs as well. 

\textbf{Example}:

Solving a Logarithmic Equation:

A logarithmic equation is an equation wherein one or more of the terms is a logarithm. 

\textbf{Example Problem}:

Solve

$$
\log x + \log(x+2) = 2
$$

\textbf{Solution}:
$$
\log x + \log(x+2) = 2
$$

$$
\log (x(x+2)) = 2
$$

$$
x(x+2) = 100
$$

$$
x^2 + 2x = 100
$$

$$
(x+1)^2 = 101
$$

$$
x+1 = \pm \sqrt{101}
$$

$$
x = -1 \pm \sqrt{101}
$$

\textbf{Another Example}:

Calculate

$$
\log\_2 8 = \frac{\log 8}{\log 2} = \frac{0.9}{0.3} = 3.
$$

Verify:

$$
2^3 = 8.
$$

\paragraph{Resources}

\begin{itemize}
    \item \href{https://mathresearch.utsa.edu/wikiFiles/MAT1053/Logarithmic_Functions/MAT1053_M5.2Logarithmic_Functions.pdf}{Logarithmic Functions, Book Chapter}
    \item \href{https://mathresearch.utsa.edu/wikiFiles/MAT1053/Logarithmic_Functions/MAT1053_M5.2Logarithmic_FunctionsGN.pdf}{Guided Notes}
\end{itemize}

\paragraph{Licensing}

Content obtained and/or adapted from:

\begin{itemize}
    \item \href{https://en.wikibooks.org/wiki/A-level_Mathematics/OCR/C2/Logarithms_and_Exponentials}{Logarithms and Exponentials, Wikibooks: A-level Mathematics/OCR/C2 under a CC BY-SA license}
\end{itemize}

\subsection{Learning Objectives}

\begin{enumerate}
    \item 1.\textbf{Understand the Inverse Relationship}: Given an exponential function $y = b^x$, convert it into its logarithmic form $x = \log\_b y$ and identify the base $b$ and the argument $y$. 2.\textbf{Apply Logarithmic Laws}: Use the product rule, quotient rule, and power rule for logarithms to simplify expressions involving logarithms, such as applying $$ \log\_b (XY) = \log\_b X + \log\_b Y, $$ $$ \log\_b \left(\frac{X}{Y}\right) = \log\_b X - \log\_b Y, $$ and $$ \log\_b (X^k) = k \log\_b X. $$ 3.\textbf{Solve Logarithmic Equations}: Solve equations involving logarithms by applying logarithmic properties, such as $$ \log\_b x + \log\_b (x + 2) = 2, $$ to find the value of $x$ through conversion and algebraic manipulation.
\end{enumerate}

\subsection{Assessment Instrument}

\begin{enumerate}
    \item Instructions Answer the following questions based on your understanding of logarithmic functions. Show all work for full credit. 1. Converting Exponential to Logarithmic Form \textbf{Question:} Convert the following exponential equation to its logarithmic form: $$ 5^x = 125 $$ 2. Applying Logarithmic Laws \textbf{Question:} Simplify the following expression using the laws of logarithms: $$ \log\_2 16 + \log\_2 4 $$ 3. Solving a Logarithmic Equation \textbf{Question:} Solve the following logarithmic equation for $x$: $$ \log\_3 (x) + \log\_3 (x + 3) = 2 $$ 4. Change of Base Formula \textbf{Question:} Calculate $\log\_7 49$ using the change of base formula: $$ \log\_7 49 = \frac{\log 49}{\log 7} $$ 5. Verifying a Logarithmic Calculation \textbf{Question:} Verify that $\log\_2 32 = 5$ is correct by converting it to an exponential form.
\end{enumerate}

%%===============================================================
\section{Logarithmic Properties}
\label{sec:logarithmic-properties}
%%===============================================================

\paragraph{Logarithmic Properties}

\paragraph{Equality Rule}

If 
$$
\log\_a(x) = \log\_a(y)
$$
then 
$$
x = y
$$

\paragraph{Product Rule}

$$
\log\_a(xy) = \log\_a(x) + \log\_a(y)
$$

\paragraph{Quotient Rule}

$$
\log\_a \left(\frac{x}{y}\right) = \log\_a(x) - \log\_a(y)
$$

\paragraph{Power Rule}

$$
\log\_a(x^n) = n \log\_a(x)
$$

\paragraph{Change-of-Base Rule}

$$
\log\_a(x) = \frac{\log\_b(x)}{\log\_b(a)}
$$
where $b$ is any positive number not equal to 1.

\paragraph{Resources}

\begin{itemize}
    \item \href{https://mathresearch.utsa.edu/wikiFiles/MAT1053/Logarithmic\%20Properties/MAT1053_M6.1Logarithmic_Properties.pdf}{Logarithmic Properties, Book Chapter}
    \item \href{https://mathresearch.utsa.edu/wikiFiles/MAT1053/Logarithmic\%20Properties/MAT1053_M6.1Logarithmic_PropertiesGN.pdf}{Guided Notes}
\end{itemize}

\subsection{Learning Objectives}

\begin{enumerate}
    \item Student Learning Objectives (SLOs)
    \item \textbf{Equality Rule}: Given the equation $$ \log\_a(x) = \log\_a(y) $$ students will be able to solve for $x$ and $y$ by setting $$ x = y. $$
    \item \textbf{Product Rule}: Students will apply the product rule to simplify expressions involving logarithms by using $$ \log\_a(xy) = \log\_a(x) + \log\_a(y). $$
    \item \textbf{Quotient and Power Rules}: Students will demonstrate the ability to use the quotient rule $$ \log\_a \left(\frac{x}{y}\right) = \log\_a(x) - \log\_a(y) $$ and the power rule $$ \log\_a(x^n) = n \log\_a(x) $$ to solve and simplify logarithmic expressions.
\end{enumerate}

\subsection{Assessment Instrument}

\begin{enumerate}
    \item Assessment Instrument for Logarithmic Properties Instructions Answer the following questions based on the logarithmic properties provided. Show all work where applicable. --- \textbf{Question 1: Equality Rule} Given that $$ \log\_3(x) = \log\_3(27), $$ solve for $x$. --- \textbf{Question 2: Product Rule} Simplify the following expression using the product rule: $$ \log\_5(20) + \log\_5(4). $$ --- \textbf{Question 3: Quotient Rule} Simplify the following expression using the quotient rule: $$ \log\_2(16) - \log\_2(4). $$ --- \textbf{Question 4: Power Rule} Simplify the following expression using the power rule: $$ \log\_7(49^3). $$ --- \textbf{Question 5: Change-of-Base Rule} Convert the following logarithm using the change-of-base rule: $$ \log\_2(32). $$ Use base 10 for the conversion. --- \textbf{Question 6: Mixed Application} Solve the following expression using the appropriate logarithmic properties: $$ \log\_4(64) - \log\_4(4) + \log\_4(2^3). $$ --- \textbf{Question 7: Application Problem} Given that $$ \log\_2(3x) = 5 $$ solve for $x$. --- \textbf{Question 8: Verification} Verify the following equality: $$ \log\_{10}(100) = \frac{\log\_{10}(10^2)}{\log\_{10}(10)}. $$
\end{enumerate}

%%===============================================================
\section{Logarithmic and Exponential Equations}
\label{sec:logarithmic-and-exponential-equations}
%%===============================================================

\paragraph{Logarithmic and Exponential Equations}

\paragraph{Logarithms and Exponents}

A logarithm is the inverse function of an exponent.

e.g. The inverse of the function $f(x) = 3^x$ is $f^{-1}(x) = \log_3 x$.

In general, $y = b^x \iff x = \log_b y$ given that $b > 0$.

\paragraph{Laws of Logarithms}

The laws of logarithms can be derived from the laws of exponentiation:

1. Product Rule:
   $$
   x^{a+b} = x^a \times x^b \iff \log a + \log b = \log(ab)
   $$

2. Quotient Rule:
   $$
   x^{a-b} = \frac{x^a}{x^b} \iff \log a - \log b = \log\left(\frac{a}{b}\right)
   $$

3. Power Rule:
   $$
   (x^a)^b = x^{ab} \iff \log(a^b) = b \log a
   $$

These laws apply to logarithms of any given base.

\paragraph{Natural Logarithms}

The natural logarithm is a logarithm with base $e$ where $e$ is a constant such that the function $e^x$ is its own derivative.

The natural logarithm has a special symbol: $\ln x$

The graph $y = e^{kx}$ exhibits exponential growth when $k > 0$ and exponential decay when $k < 0$. The inverse graph is $y = \frac{1}{k} \ln x$. 

\paragraph{Solving Logarithmic and Exponential Equations}

An exponential equation is an equation in which one or more of the terms is an exponential function. e.g. 
$$
5^x = 2^{x+2}.
$$
Exponential equations can be solved with logarithms.

e.g. Solve 
$$
3^{x+1} = 4^{2x-1}.
$$

% [Image: solution]

A logarithmic equation is an equation wherein one or more of the terms is a logarithm.

e.g. Solve 
$$
\lg x + \lg(x+2) = 2
$$

% [Image: solution 2]

\paragraph{Converting Relationships to a Linear Form}

In maths and science, it is easier to deal with linear relationships than non-linear relationships. Logarithms can be used to convert some non-linear relationships into linear relationships.

\paragraph{Exponential Relationships}

An exponential relationship is of the form 

$$
y = ab^x
$$

If we take the natural logarithm of both sides, we get 

$$
\ln y = \ln a + x \ln b
$$

We now have a linear relationship between $\ln y$ and $x$.

\textbf{e.g.} The following data is related with an exponential relationship. Determine this exponential relationship, then convert it to linear form.

(x,y)

(0,5)

(2, 45)

(4, 405)

% [Image: solution 3]

\paragraph{Converting to Linear Form}

Now convert it to linear form by taking the natural logarithm of both sides:

$$
y = 5 (3^x)
$$

$$
\ln y = \ln 5 + x \ln 3
$$

\paragraph{Power Relationships}

A power relationship is of the form 

$$
y = ax^b
$$

If we take the natural logarithm of both sides, we get 

$$
\ln y = \ln a + b \ln x
$$

This is a linear relationship between $\ln y$ and $\ln x$.

\textbf{e.g.} The amount of time that a planet takes to travel around the sun (its orbital period) and its distance from the sun are related by a power law. Use the following data to deduce this power law:

\paragraph{Planetary Data}

\begin{itemize}
    \item \textbf{Earth}: Distance from Sun = $149.6 \times 10^6$ km, Orbital Period = 365.2 days
    \item \textbf{Mars}: Distance from Sun = $227.9 \times 10^6$ km, Orbital Period = 687.0 days
    \item \textbf{Jupiter}: Distance from Sun = $778.6 \times 10^6$ km, Orbital Period = 4331 days
\end{itemize}

% [Image: solution 4]

\paragraph{Resources}

\begin{itemize}
    \item \href{https://mathresearch.utsa.edu/wikiFiles/MAT1093/Logarithmic\%20and\%20Exponential\%20Equations/Esparza\%201093\%20Notes\%207.4.pdf}{Logarithmic and Exponential Equations. Written notes created by Professor Esparza, UTSA.}
\end{itemize}

\paragraph{Licensing}

Content obtained and/or adapted from:

\begin{itemize}
    \item \href{https://en.wikibooks.org/wiki/A-level_Mathematics/CIE/Pure_Mathematics_2/Logarithmic_and_Exponential_Functions}{Logarithmic and Exponential Functions, Wikibooks: A-level Mathematics under a CC BY-SA license}
\end{itemize}

\subsection{Learning Objectives}

\begin{enumerate}
    \item \textbf{Convert Exponential Relationships to Linear Form}: Given an exponential relationship of the form $y = ab^x$, use the natural logarithm to convert the relationship into a linear form and identify the resulting linear equation. For example, convert $y = 5 (3^x)$ to its linear form.
    \item \textbf{Apply Logarithmic Laws}: Utilize the laws of logarithms (Product Rule, Quotient Rule, and Power Rule) to simplify and solve logarithmic expressions. For instance, simplify $\log\_a (xy)$ using the Product Rule and solve $\log\_a (x^n)$ using the Power Rule.
    \item \textbf{Solve Logarithmic and Exponential Equations}: Solve logarithmic and exponential equations by applying appropriate methods. For example, solve the equation $3^{x+1} = 4^{2x-1}$ using logarithmic properties and solve $\log x + \log(x+2) = 2$ to find $x$.
\end{enumerate}

\subsection{Assessment Instrument}

\begin{enumerate}
    \item Instructions Answer each question below to demonstrate your understanding of the properties and applications of logarithmic and exponential equations. Show all work for full credit. Question 1: Convert Exponential Relationships to Linear Form Given the exponential relationship: $$ y = 5 (3^x) $$
    \item Convert this exponential relationship into its linear form using the natural logarithm.
    \item Identify the resulting linear equation. Question 2: Apply Logarithmic Laws Simplify and solve the following logarithmic expressions using the laws of logarithms:
    \item Simplify the expression: $$ \log\_a (xy) $$ using the Product Rule.
    \item Solve the expression: $$ \log\_a \left(\frac{x}{y}\right) $$ using the Quotient Rule.
    \item Simplify the expression: $$ \log\_a (x^n) $$ using the Power Rule. Question 3: Solve Logarithmic and Exponential Equations Solve the following equations by applying appropriate methods:
    \item Solve the exponential equation: $$ 3^{x+1} = 4^{2x-1} $$ Show your work and provide the solution for $x$.
    \item Solve the logarithmic equation: $$ \log x + \log(x+2) = 2 $$ Show your work and provide the solution for $x$.
\end{enumerate}

%%===============================================================
\section{Exponential Growth and Decay Models}
\label{sec:exponential-growth-and-decay-models}
%%===============================================================

\paragraph{Exponential Growth and Decay Models}





\paragraph{The Natural Functions}



\paragraph{The Natural Exponent Function}



The base $e$ is a transcendental function defined by an infinite series, which approximately is $2.71828$:



$$

e = \sum\_{n=0}^{\infty} \frac{1}{n!} = \frac{1}{0!} + \frac{1}{1!} + \frac{1}{2!} + \frac{1}{3!} + \frac{1}{4!} + \cdots

$$



$e$ is one of the most important numbers in mathematics because of its use in most natural growth processes calculations. Since it is an exponential function, it follows all the laws of exponential functions. The natural exponent is the inverse function of the natural logarithm, so that:



$$

e^{\ln x} = x

$$



On the right is the graph of $e^x$.



\paragraph{The Natural Logarithm}



The natural log is a logarithmic function with the base $e$:



$$

\log\_e x

$$



It has a special notation of:



$$

\ln x

$$



Since the natural logarithm is a logarithmic function, it follows all the laws of logarithmic functions. The natural logarithm is the inverse function of the natural exponent, so that:



$$

\ln e^x = x

$$



On the right is the graph of $\ln x$.



\paragraph{Exponential Growth and Decay}



One of the primary uses of the natural functions is the calculation of natural growth or decay, which can be calculated using the following formula:



$$

y(t) = y\_0 e^{kt}

$$



where $y(t)$ is the final value, $y\_0$ is the initial value, $k$ is the growth constant, and $t$ is the elapsed time. One special decay is the half-life of an element where $k$ is defined as:



$$

k = -\frac{\ln 2}{\text{half-life}}

$$



Exponential growth and decay have a very broad range of applications, from measuring the growth of a bacteria colony to calculating interest.



\paragraph{Example of Exponential Growth}



A bacterial colony was started with 200 bacteria. In 2 hours, the colony has grown to 600 bacteria. Predict the quantity of bacteria in the colony in 6 hours.



From this, we know:



\begin{itemize}
    \item $y(t) = 600$
    \item $y\_0 = 200$
    \item $k = x$
    \item $t = 2$
\end{itemize}

Now we substitute the values into the formula:



$$

600 = 200 e^{2k}

$$



To isolate $e$:



$$

\frac{600}{200} = e^{2x}

$$



Since $\ln$ is the inverse function, we use it to remove the base $e$:



$$

\ln 3 = \ln e^{2x}

$$



Solving for $x$:



$$

1.1 \approx 2x

$$



So, the growth factor for the colony is approximately $0.55$.



Now we can predict the quantity of bacteria after 6 hours:



$$

x = 200 e^{6 \times 0.55}

$$



$$

x = 5423\ \text{bacteria}

$$



\paragraph{Example of Half-life Decay}



The element Carbon-14 has a half-life of 5730 years. When will a 100-gram sample of C-14 be reduced to 20\% C-14 and 80\% C-12?



\begin{itemize}
    \item $y(t) = 20\% \times 100 = 20$
    \item $y\_0 = 100$
    \item half-life = 5730
    \item $t = x$
\end{itemize}

We substitute all known values into the function. Note the special formula for the growth constant:



$$

20 = 100 e^{\left(-\frac{\ln 2}{5730} \times t\right)}

$$



To isolate $e$:



$$

0.2 = e^{\left(-\frac{\ln 2}{5730} \times t\right)}

$$



Since $\ln$ is the inverse function, we use it to remove the base $e$:



$$

\ln 0.2 = \ln e^{\left(-\frac{\ln 2}{5730} \times t\right)}

$$



$$

\ln 0.2 = -\frac{\ln 2}{5730} \times t

$$



Now we rearrange the equation to isolate $t$:



$$

t = -5730 \frac{\ln 0.2}{\ln 2}

$$



Solving for $t$:



$$

t \approx 13304.6\ \text{years}

$$



\paragraph{Resources}



\begin{itemize}
    \item \href{https://mathresearch.utsa.edu/wikiFiles/MAT1093/Exponential\%20growth\%20and\%20decay\%20models/Esparza\%201093\%20Notes\%207.6.pdf}{Exponential growth and decay models. Written notes created by Professor Esparza, UTSA.}
    \item \href{https://mathresearch.utsa.edu/wikiFiles/MAT1093/Exponential\%20growth\%20and\%20decay\%20models/Esparza\%201093\%20Notes\%207.6B.pdf}{Exponential growth and decay models Continued. Written notes created by Professor Esparza, UTSA.}
\end{itemize}

\paragraph{Licensing}



Content obtained and/or adapted from:



\begin{itemize}
    \item \href{https://en.wikibooks.org/wiki/A-level_Mathematics/OCR/C3/Special_Functions_and_Transformations}{Special Functions and Transformations, Wikibooks: A-level Mathematics} under a CC BY-SA license
\end{itemize}

\subsection{Learning Objectives}

\begin{enumerate}
    \item 1.\textbf{Understand the Natural Exponent Function} Students will be able to explain the definition of the natural exponent $e$, its importance in mathematics, and its relationship with the natural logarithm, including the formula $e^{\ln x} = x$. 2.\textbf{Apply Exponential Growth and Decay Models} Students will be able to use the exponential growth formula $y(t) = y_0 e^{kt}$ to solve real-world problems, such as predicting bacterial growth or calculating the decay of radioactive elements. 3.\textbf{Solve Problems Involving Half-Life Decay} Students will be able to calculate the time required for a substance to decay to a given percentage using the half-life decay formula $k = -\frac{\ln 2}{\text{half-life}}$ and $t = -5730 \frac{\ln 0.2}{\ln 2}$.
\end{enumerate}

\subsection{Assessment Instrument}

\begin{enumerate}
    \item Assessment Instrument 1. Understanding the Natural Exponent Function Question 1 Explain the definition of the natural exponent $e$ and its significance in mathematics. How is $e$ related to the natural logarithm? Use the formula $e^{\ln x} = x$ in your explanation. Question 2 Given the formula for the natural exponent $e$: $$ e = \sum\_{n=0}^{\infty} \frac{1}{n!} = \frac{1}{0!} + \frac{1}{1!} + \frac{1}{2!} + \frac{1}{3!} + \frac{1}{4!} + \cdots $$ Compute an approximation for $e$ using the first 5 terms of the series. 2. Applying Exponential Growth and Decay Models Question 3 A bacterial colony starts with 200 bacteria and grows to 600 bacteria in 2 hours. Predict the number of bacteria in the colony after 6 hours. Use the formula: $$ y(t) = y\_0 e^{kt} $$ Question 4 A population of bacteria grows according to the exponential growth model. If the growth constant $k$ is found to be 0.55, calculate the population size after 6 hours if the initial population size was 200. 3. Solving Problems Involving Half-Life Decay Question 5 Carbon-14 has a half-life of 5730 years. If you start with a 100-gram sample of Carbon-14, calculate when it will be reduced to 20\% of its original amount. Use the half-life decay formula: $$ k = -\frac{\ln 2}{\text{half-life}} $$ and solve: $$ t = -5730 \frac{\ln 0.2}{\ln 2} $$ Question 6 If a substance decays according to the formula $y(t) = y\_0 e^{\left(-\frac{\ln 2}{5730} \times t\right)}$ and you need to determine the time required for the substance to reduce to 10\% of its original amount, how would you proceed? Calculate the time $t$.
\end{enumerate}
